\documentclass[11pt,a4paper]{book}

% ============================================================
% Preamble
% ============================================================

\usepackage[margin=1.15in]{geometry}
\usepackage{amsmath,amssymb,amsthm,mathtools}
\usepackage{bm}
\usepackage{physics}
\usepackage{bbm}
\usepackage{enumitem}
\usepackage{hyperref}
\usepackage{cleveref}

% ------------------------------------------------------------
% Theorem-like environments
% ------------------------------------------------------------

\newtheorem{definition}{Definition}[chapter]
\newtheorem{example}[definition]{Example}
\newtheorem{proposition}[definition]{Proposition}
\newtheorem{theorem}[definition]{Theorem}
\newtheorem{warning}[definition]{Warning}
\newtheorem{remark}[definition]{Remark}

% ------------------------------------------------------------
% Pedagogical environments
% ------------------------------------------------------------

\newenvironment{intuition}
  {\begin{quote}\textbf{Intuition.}}
  {\end{quote}}

\newenvironment{lesson}
  {\begin{quote}\textbf{Lesson.}}
  {\end{quote}}

\newenvironment{checkpoint}
  {\begin{quote}\textbf{Checkpoint.}}
  {\end{quote}}

\newenvironment{workedexample}
  {\begin{quote}\textbf{Worked example.}}
  {\end{quote}}

% ------------------------------------------------------------
% Notation
% ------------------------------------------------------------

\newcommand{\M}{\mathcal{M}}      % model space
\newcommand{\D}{\mathcal{D}}      % data space
\newcommand{\Hh}{\mathcal{H}}     % Hilbert space
\newcommand{\B}{\mathcal{B}}      % Banach space
\newcommand{\N}{\mathcal{N}}
\newcommand{\R}{\mathbb{R}}
\newcommand{\E}{\mathbb{E}}
\newcommand{\Cov}{\operatorname{Cov}}
\newcommand{\Id}{I}

\newcommand{\prior}{\mu_0}
\newcommand{\post}{\mu^y}
\newcommand{\noise}{\eta}
\newcommand{\obs}{y}

\newcommand{\G}{G}
\newcommand{\Gadj}{G^\ast}

% ============================================================
% Document
% ============================================================

\begin{document}

\frontmatter

\title{Think First, Discretize Later:\\
A Pedagogical Introduction to Function-Space Bayesian Inverse Problems}
\author{Author Name}
\date{\today}

\maketitle

\tableofcontents

% ============================================================
% Front matter
% ============================================================

\chapter*{Preface}

\section*{What this document is about}

\section*{The central warning}

\begin{lesson}
Discretization should approximate the inverse problem.
It should not secretly define the inverse problem.
\end{lesson}

\section*{How to read this document}

\section*{Notation and conventions}

\mainmatter

% ============================================================
% ACT I
% ============================================================

\part{Act I: The Poster Child Problem}

\chapter{A Simple Inverse Problem}

\section{The continuous story we have not yet taken seriously}

We begin with a deliberately modest inverse problem. There is an unknown object,
which we call the \emph{model}, and there are finitely many noisy observations
of that model. The observations do not reveal the model directly. Instead, each
datum is obtained by applying a known measurement rule to the model and then
corrupting the result with noise.

At first, we will treat this problem in the most computationally convenient way:
we will discretize the model immediately, represent it using basis functions,
turn the forward map into a matrix, and compute a finite-dimensional Gaussian
posterior. This is the obvious route. It is also the route that will eventually
betray us.

The purpose of this first act is not to do things correctly. The purpose is to
do something plausible, watch it appear to work, and then prepare the stage for
the moment when the floorboards start to complain.

\subsection{The unknown model}

Let the physical domain be an interval
\[
    \Omega = [0,L].
\]
The unknown model is a function
\[
    m : \Omega \to \mathbb{R}.
\]
For example, $m(x)$ might describe a material property as a function of depth,
position, or time. At this stage we will not ask too carefully what kind of
function $m$ is allowed to be. Is it continuous? Square-integrable? Smooth?
Does it satisfy boundary conditions? Does it have derivatives?

For now, we only say
\[
    m \in \text{some space of functions on } \Omega.
\]

This vagueness is intentional. It reflects the way many inverse problems first
appear in practice. We know that there is a function we want to infer, and we
know that our data depend on it. The exact mathematical home of that function is
often postponed until later.

\begin{warning}
The phrase ``some space of functions'' is not harmless. It is a modelling
choice waiting to happen. If we do not make it explicitly, the discretization
will make it for us.
\end{warning}

\subsection{The observations}

Suppose we have $K$ observations,
\[
    y_1,\ldots,y_K \in \mathbb{R}.
\]
We collect them into a data vector
\[
    \bm{y}
    =
    \begin{pmatrix}
        y_1\\
        \vdots\\
        y_K
    \end{pmatrix}
    \in \mathbb{R}^K.
\]

Each observation is a noisy measurement of the unknown function. A typical
example has the form
\[
    y_i
    =
    \int_{\Omega} K_i(x)m(x)\,dx
    +
    \eta_i,
    \qquad
    i=1,\ldots,K,
\]
where $K_i$ is a known sensitivity kernel and $\eta_i$ is observational noise.

The sensitivity kernel $K_i$ tells us which parts of the model influence the
$i$th datum. If $K_i$ is concentrated near some point, then the datum is mostly
sensitive to the model near that point. If $K_i$ is broad, then the datum blends
information over a wider region.

\begin{intuition}
Each datum is a weighted average of the model, plus noise. The kernels
$K_i$ are the weights. They decide what the experiment can see and what it
must leave hidden.
\end{intuition}

\subsection{The forward operator}

We define the forward operator
\[
    \G : \M \to \mathbb{R}^K
\]
componentwise by
\[
    [\G(m)]_i
    =
    \int_{\Omega} K_i(x)m(x)\,dx,
    \qquad
    i=1,\ldots,K.
\]
Here $\M$ denotes the model space, although we have not yet chosen it
carefully. This is one of the central tensions of the story: we already know
how to write down the measurements, but we have not yet decided what
mathematical universe the unknown belongs to.

The map $\G$ is linear:
\[
    \G(\alpha m_1 + \beta m_2)
    =
    \alpha \G(m_1) + \beta \G(m_2),
\]
for scalars $\alpha,\beta \in \mathbb{R}$ and models $m_1,m_2 \in \M$.

This linearity is important. It means that we are studying a \emph{linear
inverse problem}. The difficulty will not come from nonlinearity. The difficulty
will come from a quieter source: the interaction between infinite-dimensional
unknowns, finite-dimensional data, priors, and discretization.

\subsection{Gaussian observational noise}

We assume that the noise vector
\[
    \bm{\eta}
    =
    \begin{pmatrix}
        \eta_1\\
        \vdots\\
        \eta_K
    \end{pmatrix}
\]
is Gaussian:
\[
    \bm{\eta} \sim \mathcal{N}(0,C_D),
\]
where $C_D \in \mathbb{R}^{K \times K}$ is symmetric positive definite.

The covariance matrix $C_D$ encodes the size and correlation structure of the
observational errors. If
\[
    C_D = \sigma_D^2 I_K,
\]
then the noise components are independent and have the same variance
$\sigma_D^2$. More generally, off-diagonal entries of $C_D$ allow different
data errors to be correlated.

Thus, the data model is
\[
    \bm{y}
    =
    \G(m) + \bm{\eta},
    \qquad
    \bm{\eta} \sim \mathcal{N}(0,C_D).
\]

\section{The inverse problem in its most innocent form}

\subsection{The data equation}

The inverse problem begins with the equation
\[
    \obs = \G m + \noise.
\]
Here
\[
    \obs \in \D,
    \qquad
    m \in \M,
    \qquad
    \noise \sim \mathcal{N}(0,C_D),
\]
where, for this first example,
\[
    \D = \mathbb{R}^K.
\]

This equation is compact, but it contains almost the entire drama. The unknown
$m$ belongs to a function space. The data $\obs$ belong to a finite-dimensional
space. The forward map $\G$ compresses a function into finitely many numbers.
Noise is then added.

The inverse problem asks us to travel in the opposite direction:
\[
    \obs
    \quad \leadsto \quad
    m.
\]
But the arrow is not reversible. Many different functions can give essentially
the same data, especially once noise is present.

\subsection{What we want to infer}

In a deterministic inverse problem, we might ask for a single best model
$m^\dagger$ satisfying
\[
    \G m^\dagger \approx \obs.
\]
In a Bayesian inverse problem, we ask for something richer: a probability
distribution over possible models after observing the data.

That is, we want the posterior measure
\[
    \post = \mathbb{P}(m \in \cdot \mid \obs).
\]

This posterior should tell us not only which models fit the data well, but also
how uncertain we remain. The posterior mean may summarize the central tendency,
but the posterior covariance, credible intervals, and samples reveal the spread.

\begin{lesson}
The output of a Bayesian inverse problem is not just a curve. It is a
probability measure on possible curves.
\end{lesson}

\subsection{Why the problem is indirect}

The difficulty is that the data are finite and indirect. Each datum is only a
projection, average, or blurred measurement of the model. If the model is a
function, then it has infinitely many degrees of freedom, while the data vector
has only $K$ components.

Even before noise enters, the equation
\[
    \obs = \G m
\]
usually has many solutions or no exact solution. With noise, exact equality is
not the right goal anyway. Instead, we must combine three ingredients:
\begin{enumerate}[label=\arabic*.]
    \item the information supplied by the data;
    \item the uncertainty in the observations;
    \item prior information about plausible models.
\end{enumerate}

The Bayesian framework gives us a systematic way to combine these ingredients.
But only if the ingredients themselves are chosen coherently.

\begin{intuition}
The data do not show us the model directly.
They show us shadows of the model cast through the forward operator.
\end{intuition}

% ------------------------------------------------------------

\chapter{Discretizing First}

\section{Choosing hat functions}

Now we take the computational shortcut. Instead of first deciding what the
function space $\M$ really is, we immediately introduce a finite-dimensional
basis. This is natural, practical, and familiar. It is also the seed of the
problem.

\subsection{A mesh on the physical domain}

Let
\[
    0 = x_1 < x_2 < \cdots < x_N = L
\]
be a mesh on $\Omega = [0,L]$. For simplicity, we may imagine a uniform mesh
with spacing
\[
    h = \frac{L}{N-1}.
\]

The mesh converts the continuous domain into finitely many nodes. We will
represent the unknown model by its coefficients relative to basis functions
associated with these nodes.

\subsection{Piecewise-linear basis functions}

Let
\[
    \{\phi_j\}_{j=1}^N
\]
be the standard piecewise-linear hat functions associated with the mesh. Thus
\[
    \phi_j(x_i) = \delta_{ij},
\]
and each $\phi_j$ is supported only near the node $x_j$.

The hat functions form a basis for the finite-dimensional space
\[
    \M_N
    =
    \operatorname{span}\{\phi_1,\ldots,\phi_N\}.
\]
Every function in $\M_N$ is continuous and piecewise linear.

At this point, the discretization feels like a harmless approximation. We are
not solving for an arbitrary function anymore. We are solving for a function
inside $\M_N$.

\subsection{The finite-dimensional approximation}

We write
\[
    m_N(x) = \sum_{j=1}^{N} u_j \phi_j(x).
\]
The unknown function has now been replaced by the coefficient vector
\[
    \bm{u}
    =
    \begin{pmatrix}
        u_1\\
        \vdots\\
        u_N
    \end{pmatrix}
    \in \mathbb{R}^N.
\]

The coefficient $u_j$ controls the amplitude of the $j$th hat function. Since
$\phi_j(x_i)=\delta_{ij}$, for nodal hat functions the coefficient $u_j$ is also
the value of the approximate function at node $x_j$:
\[
    m_N(x_j) = u_j.
\]

This makes the representation easy to interpret. It also quietly encourages us
to think of $\bm{u}$ as the actual unknown, rather than as a coordinate vector
representing a function.

\begin{warning}
A coefficient vector is not the same thing as a function. The vector
$\bm{u}$ becomes a function only after we choose the basis that interprets it.
\end{warning}

\section{Turning the forward map into a matrix}

\subsection{Data predicted by basis functions}

For each basis function $\phi_j$, we can compute its contribution to each datum:
\[
    G_{ij} = \G_i \phi_j.
\]
In the integral-kernel example, this means
\[
    G_{ij}
    =
    \int_{\Omega} K_i(x)\phi_j(x)\,dx.
\]

Therefore, if
\[
    m_N(x)
    =
    \sum_{j=1}^{N} u_j\phi_j(x),
\]
then linearity gives
\[
    [\G(m_N)]_i
    =
    \sum_{j=1}^N u_j \G_i\phi_j
    =
    \sum_{j=1}^N G_{ij}u_j.
\]
In vector form,
\[
    \G(m_N) = G\bm{u},
\]
where
\[
    G \in \mathbb{R}^{K \times N}.
\]

\subsection{The naive matrix problem}

The continuous-looking inverse problem has now become the finite-dimensional
linear model
\[
    \bm{y} = G \bm{u} + \bm{\eta},
\]
with
\[
    \bm{\eta} \sim \mathcal{N}(0,C_D).
\]

This is a standard Bayesian linear regression problem. Everything appears to be
finite-dimensional. Everything appears to be familiar. We have a design matrix
$G$, unknown coefficients $\bm{u}$, data $\bm{y}$, and Gaussian noise.

At this stage, it is tempting to forget the function $m$ entirely and simply
declare victory over the infinite-dimensional problem. The beast has been put
in a matrix-shaped cage.

But we will soon find out the beast cannot be contained - it is after all an infinite dimensional beast.

\section{A Gaussian prior on coefficients}

\subsection{The tempting choice}

To complete the Bayesian formulation, we need a prior distribution on the
unknown coefficients. The most tempting choice is
\[
    \bm{u} \sim \mathcal{N}(0,\sigma^2 I_N).
\]
Equivalently,
\[
    C_N = \sigma^2 I_N.
\]

This says that the coefficients are independent, centred Gaussian random
variables:
\[
    u_j \sim \mathcal{N}(0,\sigma^2),
    \qquad
    j=1,\ldots,N,
\]
with no prior correlation between different nodes.

More generally, we may allow a nonzero prior mean
\[
    \bm{u}_0 \in \mathbb{R}^N
\]
and write
\[
    \bm{u} \sim \mathcal{N}(\bm{u}_0,C_N).
\]
For the naive example, however, the important case is
\[
    C_N = \sigma^2 I_N.
\]

\subsection{Why this looks reasonable at first glance}

The prior
\[
    \mathcal{N}(0,\sigma^2 I_N)
\]
looks reasonable for several reasons.

First, it is simple. It introduces only one scale parameter, $\sigma$, which
controls the typical size of the coefficients.

Second, it is symmetric in the coefficients. No node is preferred over any
other node.

Third, it is computationally convenient. The covariance matrix is diagonal, easy
to store, easy to invert, and easy to sample from.

Fourth, it resembles a familiar finite-dimensional default: when we do not know
much about a vector, we often place an isotropic Gaussian prior on it.

These arguments are not silly. In a genuinely finite-dimensional problem, they
may be defensible. The trouble begins when this prior is used as a stand-in for
a probability distribution over functions and then carried through mesh
refinement.

\begin{intuition}
A diagonal prior says that every coefficient is allowed to wiggle independently.
As the mesh becomes finer, we add more wiggles. Unless their size is controlled
in a way that depends on the function space, the wiggles do not politely vanish.
They accumulate.
\end{intuition}

\section{The naive posterior}

\subsection{Finite-dimensional Gaussian conditioning}

We now have the finite-dimensional Bayesian linear model
\[
    \bm{u} \sim \mathcal{N}(\bm{u}_0,C_N),
\]
and
\[
    \bm{y}\mid \bm{u}
    \sim
    \mathcal{N}(G\bm{u},C_D).
\]

Because both the prior and the noise are Gaussian, the posterior distribution
of $\bm{u}$ given $\bm{y}$ is also Gaussian:
\[
    \bm{u}\mid \bm{y}
    \sim
    \mathcal{N}(\bar{\bm{u}}_N^y,C_N^y).
\]

This is one of the great computational conveniences of linear Gaussian inverse
problems. The posterior is available in closed form. No sampling method is
needed to obtain the posterior mean or covariance.

\subsection{Posterior covariance}

The posterior covariance is
\[
    C_N^y
    =
    C_N
    -
    C_N G^T
    \left(
        G C_N G^T + C_D
    \right)^{-1}
    G C_N.
\]

The matrix
\[
    G C_N G^T + C_D
\]
is the covariance of the predicted data, including both prior uncertainty
pushed through the forward model and observational noise.

The term
\[
    C_N G^T
    \left(
        G C_N G^T + C_D
    \right)^{-1}
    G C_N
\]
is the amount of prior covariance removed by observing the data. In directions
where the data are informative, the posterior uncertainty shrinks. In directions
that the data cannot see, the posterior remains close to the prior.

This last sentence is the fuse. If the prior is badly behaved in the unseen
directions, then the posterior inherits that bad behaviour.

\subsection{Posterior mean}

The posterior mean is
\[
    \bar{\bm{u}}_N^y
    =
    \bm{u}_0
    +
    C_N G^T
    \left(
        G C_N G^T + C_D
    \right)^{-1}
    \left(
        \bm{y} - G\bm{u}_0
    \right).
\]

The term
\[
    \bm{y} - G\bm{u}_0
\]
is the data residual relative to the prior mean. The operator
\[
    C_N G^T
    \left(
        G C_N G^T + C_D
    \right)^{-1}
\]
maps this residual back into coefficient space.

For the zero-mean isotropic prior $C_N=\sigma^2 I_N$, this becomes
\[
    \bar{\bm{u}}_N^y
    =
    \sigma^2 G^T
    \left(
        \sigma^2 GG^T + C_D
    \right)^{-1}
    \bm{y}.
\]

Everything here is finite-dimensional and computable. We can assemble $G$,
choose $C_N$ and $C_D$, compute the posterior mean, and reconstruct a function
from the coefficient vector.

\section{The first plot: it looks fine}

\subsection{Reconstructing the posterior mean}

Once the posterior mean coefficients $\bar{\bm{u}}_N^y$ are computed, we convert
them back into a function:
\[
    \bar{m}_N^y(x)
    =
    \sum_{j=1}^N
    \bar{u}_{N,j}^y \phi_j(x).
\]

This reconstructed posterior mean is the curve we would naturally plot against
the true model. It is smooth enough to look respectable, flexible enough to
follow the main features of the truth, and stable enough that nothing appears
obviously broken.

In a practical workflow, this is often the first moment of relief: the code
runs, the matrix dimensions agree, the posterior mean has the right shape, and
the curve seems to pass the eye test.

\subsection{Comparing the mean with the truth}

Suppose the true model is $m_{\mathrm{true}}$. Then the synthetic data are
generated by
\[
    \bm{y}
    =
    \G m_{\mathrm{true}}
    +
    \bm{\eta}.
\]
After computing $\bar{m}_N^y$, we compare
\[
    \bar{m}_N^y(x)
    \quad
    \text{with}
    \quad
    m_{\mathrm{true}}(x).
\]

If the observations are informative enough and the noise level is moderate, the
posterior mean may recover the broad features of the truth. Peaks may appear in
approximately the right places. Long-wavelength structure may be captured.
The result may look scientifically plausible.

This is the dangerous part. The posterior mean can look sensible even when the
posterior distribution is not.

\subsection{The false comfort of a plausible curve}

The posterior mean is only one summary of the posterior. It does not tell us
whether typical posterior samples are reasonable. It does not tell us whether
posterior uncertainty is controlled under mesh refinement. It does not tell us
whether the finite-dimensional posterior is approaching a meaningful
infinite-dimensional object.

A plausible mean can therefore hide a pathological posterior.

The reason is simple: the data constrain only a finite number of directions. If
the number of basis functions $N$ is larger than the number of observations $K$,
then there are directions in coefficient space that are weakly constrained or
entirely invisible to the data. In those directions, the posterior covariance is
controlled mostly by the prior. If the prior was chosen without reference to a
continuous function-space model, there is no reason to expect that uncertainty
to behave sensibly as $N$ increases.

\begin{checkpoint}
At this stage, the posterior mean may look acceptable.
The trap is that a reasonable-looking mean can hide an unreasonable posterior measure.
\end{checkpoint}

\section{What we have done so far}

Let us summarize the naive workflow.

\begin{enumerate}[label=\arabic*.]
    \item We began with an unknown function $m$ on a domain $\Omega$.
    \item We introduced data through a linear forward map $\G$.
    \item We discretized the model immediately using hat functions.
    \item We replaced the function $m$ by coefficients $\bm{u}$.
    \item We replaced the forward operator by a matrix $G$.
    \item We placed a Gaussian prior directly on the coefficients.
    \item We computed the finite-dimensional Gaussian posterior.
    \item We reconstructed and plotted the posterior mean.
\end{enumerate}

Nothing in this list looks outrageous. That is precisely why the example is
useful. The mistake is not dramatic. It does not announce itself by crashing the
code or producing nonsense immediately. It hides behind a plausible curve.

\begin{lesson}
The first act ends with a posterior mean that looks alright.
The second act begins by asking what the posterior itself is doing.
\end{lesson}

% ============================================================
% ACT II
% ============================================================

\part{Act II: Something Is Wrong}

\chapter{Sampling the Naive Posterior}

\section{Why the mean is not the whole posterior}

In the previous act, the posterior mean looked respectable. It passed the first
visual test: it had the right broad shape, it responded to the data, and it did
not immediately look like numerical wreckage.

But a Bayesian inverse problem does not produce only a mean. It produces a
probability measure. The mean is only one summary of that measure, and a
dangerous summary when viewed alone. A posterior mean can look calm while the
posterior distribution around it is thrashing like a stormy sea.

\subsection{A posterior is a probability measure}

The finite-dimensional posterior obtained in the naive discretization is
\[
    \bm{u}\mid \bm{y}
    \sim
    \mathcal{N}(\bar{\bm{u}}_N^y,C_N^y).
\]
This is a probability measure on $\mathbb{R}^N$. Once we reconstruct functions
using the hat basis, it also induces a probability measure on the finite element
space
\[
    \M_N = \operatorname{span}\{\phi_1,\ldots,\phi_N\}.
\]
That is, a random coefficient vector
\[
    \bm{u}
    =
    \begin{pmatrix}
        u_1\\
        \vdots\\
        u_N
    \end{pmatrix}
\]
produces a random function
\[
    m_N(x)
    =
    \sum_{j=1}^N u_j\phi_j(x).
\]

The posterior mean is the function
\[
    \bar{m}_N^y(x)
    =
    \sum_{j=1}^N \bar{u}_{N,j}^y\phi_j(x).
\]
But the posterior itself consists of all possible functions generated by
sampling
\[
    \bm{u} \sim \mathcal{N}(\bar{\bm{u}}_N^y,C_N^y).
\]

\begin{lesson}
The posterior mean asks: ``Where is the centre of the cloud?''
The posterior measure asks: ``What is the whole cloud doing?''
\end{lesson}

\subsection{Uncertainty lives in the samples}

To understand whether the posterior is sensible, we must inspect more than the
mean. We must ask what typical posterior samples look like.

A sample from the posterior is a possible model compatible with three things:
\begin{enumerate}[label=\arabic*.]
    \item the observed data;
    \item the assumed noise model;
    \item the assumed prior.
\end{enumerate}

If the samples look physically plausible, then the posterior uncertainty may be
telling a coherent story. If the samples fluctuate wildly, then something is
wrong, even if the mean still looks acceptable.

This distinction is essential. A mean can smooth away pathologies. Samples do
not. Samples are less polite. They show the posterior with its socks off.

\section{How to sample a finite-dimensional Gaussian}

\subsection{The mathematical recipe}

Let
\[
    \bm{u} \sim \mathcal{N}(\bar{\bm{u}},C),
\]
where
\[
    \bar{\bm{u}} \in \mathbb{R}^N,
    \qquad
    C \in \mathbb{R}^{N\times N}.
\]
Assume first that $C$ is symmetric positive definite. Then we may factorize
\[
    C = LL^T,
\]
for example using a Cholesky factorization. Now draw
\[
    \bm{\xi} \sim \mathcal{N}(0,I_N),
\]
and set
\[
    \bm{u}
    =
    \bar{\bm{u}} + L\bm{\xi}.
\]

This gives the desired distribution. Indeed,
\[
    \mathbb{E}[\bm{u}]
    =
    \mathbb{E}[\bar{\bm{u}} + L\bm{\xi}]
    =
    \bar{\bm{u}} + L\mathbb{E}[\bm{\xi}]
    =
    \bar{\bm{u}},
\]
and
\[
    \operatorname{Cov}(\bm{u})
    =
    \operatorname{Cov}(L\bm{\xi})
    =
    L\operatorname{Cov}(\bm{\xi})L^T
    =
    L I_N L^T
    =
    C.
\]

Thus,
\[
    \bm{u}
    =
    \bar{\bm{u}} + L\bm{\xi}
\]
is a sample from
\[
    \mathcal{N}(\bar{\bm{u}},C).
\]

If $C$ is positive semidefinite rather than positive definite, we may use an
eigenvalue decomposition
\[
    C = Q\Lambda Q^T,
\]
where
\[
    \Lambda
    =
    \operatorname{diag}(\lambda_1,\ldots,\lambda_N),
    \qquad
    \lambda_j \geq 0.
\]
Then a covariance square-root is
\[
    C^{1/2}
    =
    Q\Lambda^{1/2}Q^T,
\]
and we may sample using
\[
    \bm{u}
    =
    \bar{\bm{u}} + C^{1/2}\bm{\xi}.
\]

\subsection{The practical recipe}

For the naive posterior, the sampling procedure is as follows.

\begin{enumerate}[label=\arabic*.]
    \item Compute the posterior mean $\bar{\bm{u}}_N^y$.
    \item Compute the posterior covariance $C_N^y$.
    \item Factorize $C_N^y$, for example by Cholesky or by an eigenvalue
    decomposition.
    \item Draw independent standard normal vectors
    \[
        \bm{\xi}^{(s)} \sim \mathcal{N}(0,I_N),
        \qquad
        s=1,\ldots,S.
    \]
    \item Form coefficient samples
    \[
        \bm{u}^{(s)}
        =
        \bar{\bm{u}}_N^y
        +
        L\bm{\xi}^{(s)}.
    \]
    \item Reconstruct sample functions using the hat basis:
    \[
        m_N^{(s)}(x)
        =
        \sum_{j=1}^N u_j^{(s)}\phi_j(x).
    \]
\end{enumerate}

Here $S$ is the number of posterior samples we want to inspect.

\begin{remark}
In exact arithmetic, the posterior covariance is symmetric positive
semidefinite. In floating-point arithmetic, tiny negative eigenvalues may appear
because of round-off error. A practical implementation often symmetrizes the
matrix,
\[
    C_N^y \leftarrow \frac{1}{2}(C_N^y + (C_N^y)^T),
\]
and discards eigenvalues that are negative only up to numerical precision.
\end{remark}

\section{What the samples reveal}

\subsection{The samples fluctuate wildly}

After drawing posterior samples, we reconstruct the corresponding functions
\[
    m_N^{(s)}(x)
    =
    \sum_{j=1}^N u_j^{(s)}\phi_j(x).
\]
Now the problem begins to show itself.

The posterior mean may track the main features of the true model, but the
samples can vary dramatically around it. They may contain rapid oscillations,
large local deviations, or high-frequency behaviour that the data never
justified.

The plot now tells a different story. The mean says:
\[
    \text{``Everything is fine.''}
\]
The samples say:
\[
    \text{``We have been ignoring the basement.''}
\]

\subsection{The data are too few}

The reason is structural. The data vector has only $K$ entries:
\[
    \bm{y} \in \mathbb{R}^K.
\]
The coefficient vector has $N$ entries:
\[
    \bm{u} \in \mathbb{R}^N.
\]
If
\[
    N > K,
\]
then there are more unknown coefficients than observations.

The forward matrix
\[
    G \in \mathbb{R}^{K\times N}
\]
can have rank at most $K$:
\[
    \operatorname{rank}(G) \leq K.
\]
Therefore, if $N>K$, the nullspace of $G$ has dimension at least
\[
    \dim \operatorname{null}(G)
    =
    N-\operatorname{rank}(G)
    \geq
    N-K.
\]

Every vector
\[
    \bm{v} \in \operatorname{null}(G)
\]
satisfies
\[
    G\bm{v} = 0.
\]
So moving in the direction $\bm{v}$ does not change the predicted data.

\subsection{The number of basis functions is too large}

Increasing the number of basis functions increases the number of unknowns. But
unless the number of observations increases in a compatible way, it also
increases the number of directions that the data cannot see.

This is not automatically a problem. Infinite-dimensional inverse problems
always contain directions that are weakly observed or unobserved. The real
question is: what controls those directions?

In a coherent Bayesian formulation, the prior controls them in a way that is
consistent with a probability measure on a function space. In the naive
formulation, the prior is merely
\[
    C_N = \sigma^2 I_N.
\]
So every new coefficient receives a fresh independent amount of variance
$\sigma^2$.

The mesh is refined. More basis functions appear. More independent prior
wiggles are introduced. The data cannot see most of them. The posterior inherits
them.

\begin{warning}
The problem is not simply that $N>K$. The problem is that the unseen directions
are filled with variance that has not been designed to have a meaningful
continuous limit.
\end{warning}

\section{The underdetermined system}

\subsection{More unknown coefficients than observations}

In the naive discretized problem,
\[
    \bm{y} = G\bm{u} + \bm{\eta},
\]
the matrix $G$ maps from $\mathbb{R}^N$ to $\mathbb{R}^K$. When $N>K$, the
system is underdetermined. There are more coefficient directions than data
directions.

The likelihood can constrain only directions that influence the data. These are
directions in the row space of $G$. Directions in the nullspace of $G$ are
invisible to the likelihood.

If
\[
    \bm{v} \in \operatorname{null}(G),
\]
then
\[
    G(\bm{u}+\alpha \bm{v})
    =
    G\bm{u}
    +
    \alpha G\bm{v}
    =
    G\bm{u}
\]
for every scalar $\alpha \in \mathbb{R}$.

Thus, all coefficient vectors along the affine line
\[
    \bm{u} + \alpha\bm{v},
    \qquad
    \alpha \in \mathbb{R},
\]
produce exactly the same noiseless data prediction.

\subsection{Invisible directions}

Let
\[
    \mathcal{N}(G)
    =
    \{\bm{v}\in\mathbb{R}^N : G\bm{v}=0\}
\]
denote the nullspace of $G$.

For any $\bm{v}\in \mathcal{N}(G)$, the data contain no information about the
component of $\bm{u}$ in the direction $\bm{v}$. Therefore, the posterior
variance in that direction is controlled entirely by the prior.

To see this clearly, consider the naive prior
\[
    C_N = \sigma^2 I_N.
\]
The posterior covariance is
\[
    C_N^y
    =
    \sigma^2 I_N
    -
    \sigma^4 G^T
    \left(
        \sigma^2 GG^T + C_D
    \right)^{-1}
    G.
\]
If $\bm{v}\in\mathcal{N}(G)$, then $G\bm{v}=0$, and hence
\[
    C_N^y\bm{v}
    =
    \sigma^2 \bm{v}.
\]
So every nullspace direction remains a posterior covariance eigenvector with
eigenvalue $\sigma^2$.

In words: the data remove none of the prior variance in directions they cannot
see.

\subsection{Directions unconstrained by the data}

Let
\[
    r_N = \operatorname{rank}(G).
\]
Then
\[
    r_N \leq K.
\]
The nullspace has dimension
\[
    N-r_N.
\]
On this nullspace, the posterior covariance keeps the value $\sigma^2$ in each
orthonormal direction.

Therefore, at least
\[
    N-K
\]
directions retain posterior variance $\sigma^2$.

This is the first precise sign that something is wrong. Refining the mesh adds
new directions. The data cannot constrain them. The naive prior gives them all
the same amount of variance. The posterior uncertainty grows a tail made of
unseen degrees of freedom.

\begin{intuition}
The data illuminate only a finite number of directions.
The remaining directions sit in the dark, and the naive prior lets them grow teeth.
\end{intuition}

% ------------------------------------------------------------

\chapter{The Divergence of the Naive Posterior}

\section{The quantity we should monitor}

The plot of the posterior mean is not enough. We need quantities that measure
the size of the posterior uncertainty. There are several natural choices,
depending on the geometry we use.

\subsection{Expected posterior energy}

One diagnostic is the expected squared size of a posterior sample:
\[
    \mathbb{E}_{\post_N}\|m_N\|^2.
\]
If the posterior mean is not zero, then this decomposes into
\[
    \mathbb{E}_{\post_N}\|m_N\|^2
    =
    \|\bar{m}_N^y\|^2
    +
    \mathbb{E}_{\post_N}\|m_N-\bar{m}_N^y\|^2.
\]
The first term measures the size of the posterior mean. The second term
measures the posterior spread.

In coefficient space, with the Euclidean norm, this becomes
\[
    \mathbb{E}_{\post_N}\|\bm{u}\|_2^2
    =
    \|\bar{\bm{u}}_N^y\|_2^2
    +
    \operatorname{Tr}(C_N^y).
\]

Thus, the trace of the posterior covariance measures the total posterior
variance in coefficient space.

\subsection{Pointwise variance}

Another diagnostic is the pointwise posterior variance:
\[
    \operatorname{Var}_{\post_N}(m_N(x)).
\]
Since
\[
    m_N(x)
    =
    \sum_{j=1}^N u_j\phi_j(x),
\]
we may define the vector of basis functions evaluated at $x$:
\[
    \bm{\phi}(x)
    =
    \begin{pmatrix}
        \phi_1(x)\\
        \vdots\\
        \phi_N(x)
    \end{pmatrix}.
\]
Then
\[
    m_N(x)
    =
    \bm{\phi}(x)^T\bm{u}.
\]
Therefore,
\[
    \operatorname{Var}_{\post_N}(m_N(x))
    =
    \bm{\phi}(x)^T C_N^y \bm{\phi}(x).
\]

This tells us how uncertain the posterior is at each point $x$.

\subsection{Integrated variance}

We can also integrate the pointwise variance over the physical domain:
\[
    \int_{\Omega}
    \operatorname{Var}_{\post_N}(m_N(x))\,dx.
\]
Using the expression above,
\[
    \int_{\Omega}
    \operatorname{Var}_{\post_N}(m_N(x))\,dx
    =
    \int_{\Omega}
    \bm{\phi}(x)^T C_N^y \bm{\phi}(x)\,dx.
\]
This can be written as
\[
    \operatorname{Tr}(M C_N^y),
\]
where
\[
    M_{ij}
    =
    \int_{\Omega}\phi_i(x)\phi_j(x)\,dx
\]
is the finite element mass matrix.

\begin{remark}
The coefficient-space trace $\operatorname{Tr}(C_N^y)$ and the
function-space quantity $\operatorname{Tr}(M C_N^y)$ are not the same object.
This distinction is one of the clues that geometry matters. For the moment,
however, we follow the naive workflow on its own terms and inspect the
coefficient-space posterior it has created.
\end{remark}

\section{What happens as the mesh is refined}

\subsection{Increasing the number of basis functions}

Mesh refinement means that we increase the number of basis functions:
\[
    N \to \infty.
\]
The finite-dimensional spaces
\[
    \M_N
    =
    \operatorname{span}\{\phi_1,\ldots,\phi_N\}
\]
become richer. They can represent finer and finer variations in the model.

This sounds like progress. In a well-designed numerical method, increasing $N$
should improve the approximation of a fixed continuous problem.

But here we did not begin with a fixed continuous prior measure. We began with
the sequence of finite-dimensional priors
\[
    \bm{u}_N \sim \mathcal{N}(0,\sigma^2 I_N).
\]
Each time $N$ changes, we have a new Gaussian distribution on a new coordinate
space. We have not shown that these distributions approximate any probability
measure on functions.

\subsection{The expected variance grows}

For the naive prior,
\[
    C_N = \sigma^2 I_N.
\]
The posterior covariance is
\[
    C_N^y
    =
    \sigma^2 I_N
    -
    \sigma^4 G^T
    \left(
        \sigma^2 GG^T + C_D
    \right)^{-1}
    G.
\]

The data can reduce variance only in directions that affect the data. Since
\[
    \operatorname{rank}(G) \leq K,
\]
there are at least $N-K$ directions that the data cannot constrain. In each of
these directions, the posterior variance remains exactly $\sigma^2$.

Consequently,
\[
    \operatorname{Tr}(C_N^y)
    \geq
    \sigma^2 (N-K).
\]
As $N$ increases while $K$ is fixed,
\[
    \operatorname{Tr}(C_N^y)
    \to
    \infty.
\]

This means that the total posterior variance in coefficient space grows without
bound.

\subsection{The posterior does not converge to a sensible limit}

If the posterior covariance grows in total variance as the mesh is refined, then
the sequence of posteriors is not settling down to a meaningful limiting
posterior. Instead of approximating a fixed probability measure on functions, we
are generating a sequence of increasingly large finite-dimensional probability
clouds.

The data may continue to control a small number of visible directions. But the
number of invisible directions keeps increasing, and the naive prior keeps
feeding each of them the same variance.

The posterior mean may still look acceptable because it lies mostly in the
directions informed by the data. The posterior samples reveal the full story:
there is too much uncertainty in directions that were introduced by the
discretization rather than by a coherent continuous model.

\begin{lesson}
Refinement should reveal the limit of the problem.
Here, refinement reveals that the problem was never properly posed.
\end{lesson}

\section{A mathematical diagnosis}

\subsection{The covariance trace}

For a centred Gaussian vector
\[
    \bm{u} \sim \mathcal{N}(0,C_N),
\]
the expected squared Euclidean norm is
\[
    \mathbb{E}\|\bm{u}\|_2^2
    =
    \operatorname{Tr}(C_N).
\]
Indeed,
\[
    \mathbb{E}\|\bm{u}\|_2^2
    =
    \mathbb{E}[\bm{u}^T\bm{u}]
    =
    \mathbb{E}[\operatorname{Tr}(\bm{u}\bm{u}^T)]
    =
    \operatorname{Tr}(\mathbb{E}[\bm{u}\bm{u}^T])
    =
    \operatorname{Tr}(C_N).
\]

For a non-centred Gaussian
\[
    \bm{u} \sim \mathcal{N}(\bar{\bm{u}},C_N),
\]
we have
\[
    \mathbb{E}\|\bm{u}\|_2^2
    =
    \|\bar{\bm{u}}\|_2^2
    +
    \operatorname{Tr}(C_N).
\]

Thus, the trace measures the total variance of the Gaussian cloud.

\subsection{Why $C_N = \sigma^2 I_N$ is dangerous}

For the naive coefficient prior,
\[
    C_N = \sigma^2 I_N.
\]
Therefore,
\[
    \operatorname{Tr}(C_N)
    =
    \operatorname{Tr}(\sigma^2 I_N)
    =
    \sigma^2\operatorname{Tr}(I_N)
    =
    \sigma^2 N.
\]

So the prior expected coefficient energy is
\[
    \mathbb{E}\|\bm{u}\|_2^2
    =
    \sigma^2 N.
\]
As $N$ grows, this quantity grows linearly.

The posterior improves this only in directions seen by the data. Since the data
space has dimension $K$, the data can remove variance from at most $K$
independent directions. At least $N-K$ directions keep their prior variance.

Therefore,
\[
    \operatorname{Tr}(C_N^y)
    \geq
    \sigma^2(N-K).
\]

\subsection{The limit is unbounded}

Taking the refinement limit,
\[
    N \to \infty,
\]
with $K$ fixed, we obtain
\[
    \operatorname{Tr}(C_N^y)
    \to
    \infty.
\]
In particular,
\[
    \mathbb{E}_{\post_N}
    \|\bm{u}\|_2^2
    =
    \|\bar{\bm{u}}_N^y\|_2^2
    +
    \operatorname{Tr}(C_N^y)
    \to
    \infty,
\]
provided the posterior mean does not cancel this behaviour in some artificial
way. Since the covariance term alone diverges, the posterior spread is already
unbounded.

This is the mathematical version of what the samples showed visually.

\begin{proposition}
Let $G\in\mathbb{R}^{K\times N}$ and let the prior covariance be
$C_N=\sigma^2 I_N$, with $\sigma>0$. Assume the observational noise covariance
$C_D$ is symmetric positive definite. Then the posterior covariance
\[
    C_N^y
    =
    \sigma^2 I_N
    -
    \sigma^4 G^T
    \left(
        \sigma^2 GG^T + C_D
    \right)^{-1}
    G
\]
satisfies
\[
    \operatorname{Tr}(C_N^y)
    \geq
    \sigma^2(N-K).
\]
Hence, if $K$ is fixed and $N\to\infty$, then
\[
    \operatorname{Tr}(C_N^y)\to\infty.
\]
\end{proposition}

\begin{proof}
Since $\operatorname{rank}(G)\leq K$, the nullspace of $G$ has dimension at
least $N-K$. Let
\[
    \{\bm{v}_1,\ldots,\bm{v}_{N-r}\}
\]
be an orthonormal basis for $\operatorname{null}(G)$, where
$r=\operatorname{rank}(G)$. For each such vector,
\[
    G\bm{v}_j = 0.
\]
Therefore,
\[
    C_N^y\bm{v}_j
    =
    \sigma^2\bm{v}_j.
\]
Thus $C_N^y$ has at least $N-r$ eigenvalues equal to $\sigma^2$. Since
$r\leq K$, it has at least $N-K$ such eigenvalues. Hence
\[
    \operatorname{Tr}(C_N^y)
    \geq
    \sigma^2(N-K).
\]
\end{proof}

\section{The core mistake}

The problem is not that Gaussian priors are bad. The problem is not that finite
elements are bad. The problem is not even that the posterior mean was computed
incorrectly.

The problem is that we chose the prior after discretization, directly on the
coefficient vector, without asking whether the resulting sequence of priors has
a meaningful function-space limit.

\begin{warning}
The prior was not chosen as a probability measure on a function space.
It was chosen as a sequence of unrelated finite-dimensional matrices.
\end{warning}

A matrix
\[
    C_N \in \mathbb{R}^{N\times N}
\]
can be a perfectly valid covariance matrix in $\mathbb{R}^N$ for every fixed
$N$. But this does not mean that the sequence
\[
    C_1,C_2,C_3,\ldots
\]
defines a valid covariance operator on a function space.

The continuous question is not merely:
\[
    \text{Is } C_N \text{ positive definite for each } N?
\]
The continuous question is:
\[
    \text{Do these finite-dimensional distributions approximate a coherent
    probability measure on functions?}
\]

For the naive diagonal prior with fixed variance per coefficient, the answer is
no.

\section{What this means philosophically}

\subsection{The discretized problem replaced the original problem}

At the beginning, we pretended to have a continuous inverse problem. But we did
not fully define it. We did not specify the model space. We did not specify a
prior measure on that space. We did not specify the geometry needed to define
the correct adjoint.

Instead, we discretized first and formulated the Bayesian problem on the
coefficient vector:
\[
    \bm{u} \in \mathbb{R}^N.
\]
So the actual problem we solved was not the continuous inverse problem. It was a
finite-dimensional problem created by the discretization.

The mesh did not approximate the model. The mesh became the model.

\subsection{Refinement did not improve the approximation}

In a healthy numerical method, refinement means that the discrete solution gets
closer to the solution of the original problem. But here, refinement changes the
probabilistic model itself.

For each $N$, we choose
\[
    \bm{u}_N \sim \mathcal{N}(0,\sigma^2 I_N).
\]
When $N$ increases, this is not a better approximation of the same prior. It is
a different prior on a different space, with more independent random
coefficients.

Thus, refinement does not only increase resolution. It also injects additional
prior variance.

\subsection{Refinement changed the problem}

The posterior for $N=10$ and the posterior for $N=50$ are not merely two
approximations of the same infinite-dimensional posterior. Under the naive
construction, they are posteriors for two different Bayesian models.

The notation
\[
    \post_N
\]
should therefore make us nervous. It is not automatically a convergent sequence.
It is only a sequence of probability measures produced by a sequence of
finite-dimensional choices.

To know whether
\[
    \post_N \to \post
\]
for some meaningful function-space posterior $\post$, we must first define
$\post$. That requires choosing the model space, data space, forward map, noise
model, and prior measure before discretizing.

\begin{lesson}
A discretization is not innocent.
If the continuous problem has not been defined, the discretization becomes the model.
\end{lesson}

\section{The moral of Act II}

The naive posterior did not fail because the algebra was wrong. The algebra was
perfectly competent. That is what makes the failure instructive.

The failure came from asking finite-dimensional linear algebra to make
infinite-dimensional modelling choices on our behalf.

We started with a function but placed the prior on coefficients. We refined the
mesh but did not refine an approximation to a fixed prior measure. We used the
posterior mean as evidence of success while ignoring the posterior samples. Once
we sampled, the hidden degrees of freedom marched out carrying little banners
that read:
\[
    \text{You forgot the function space.}
\]

\begin{checkpoint}
The next act begins by stepping back from the computation. Before we choose a
basis, before we build a matrix, and before we sample anything, we must ask:
what is the mathematical problem we are actually trying to solve?
\end{checkpoint}

% ============================================================
% ACT III
% ============================================================

\part{Act III: What Problem Are We Actually Solving?}

\chapter{Stepping Back}

\section{The real question}

The previous act ended with an unpleasant discovery. The naive posterior did
not fail because the Gaussian conditioning formula was wrong. It failed because
we had never properly said what problem the discretization was supposed to
approximate.

We began with a function, replaced it by coefficients, put a Gaussian prior on
those coefficients, and hoped that mesh refinement would behave like
mathematical progress. It did not. The posterior mean looked acceptable, but
the posterior measure itself became increasingly unreasonable.

So now we pause. Before choosing basis functions, matrices, quadrature rules,
or numerical solvers, we ask a more basic question:

\[
    \text{What is the inverse problem?}
\]

This question is not a request for code. It is a request for modelling choices.
The unknown must live somewhere. The data must live somewhere. The forward map
must connect those two places. The noise model must describe observational
uncertainty. The prior must be a probability measure on the unknown itself, not
merely a convenient matrix attached to a discretization.

\begin{lesson}
A Bayesian inverse problem is not born when we assemble a matrix.
It is born when we specify the mathematical objects whose relationship we want
to infer.
\end{lesson}

\subsection{What is the unknown object?}

The first question is deceptively simple:
\[
    \text{What is } m?
\]

In the toy example, $m$ is a function on a physical domain
\[
    \Omega = [0,L].
\]
But saying that $m$ is a function is not enough. There are many kinds of
functions. Some are continuous, some are merely square-integrable, some have
weak derivatives, some satisfy boundary conditions, and some are so irregular
that pointwise values are not even meaningful.

Thus, we must decide what class of objects we are willing to regard as possible
models.

For example, we might choose
\[
    m \in L^2(\Omega),
\]
if we only want square-integrability. Or we might choose
\[
    m \in H^1(\Omega),
\]
if we want one weak derivative to be square-integrable. Or we might choose a
more regular Sobolev space if smoothness is part of our prior information.

At this stage, the precise choice is not yet the point. The point is that there
must be a choice.

\begin{warning}
The unknown object is not a vector of coefficients. Coefficients are only a
coordinate representation after a basis has been chosen.
\end{warning}

The model space determines what counts as an admissible model, what convergence
means, what it means for two models to be close, and what kinds of uncertainty
can be placed on the model.

\subsection{What kind of data do we observe?}

The second question is:
\[
    \text{What is } \obs?
\]

In the toy problem, the data are finitely many real numbers:
\[
    \obs =
    \begin{pmatrix}
        y_1\\
        \vdots\\
        y_K
    \end{pmatrix}
    \in \mathbb{R}^K.
\]
Thus, a natural data space is
\[
    \D = \mathbb{R}^K.
\]

This looks simple, and in many ways it is. Finite-dimensional spaces are
mathematically friendly. All norms on $\mathbb{R}^K$ induce the same topology,
linear maps into and out of $\mathbb{R}^K$ are easier to control, and
completeness is automatic.

But even in finite dimensions, the geometry matters. The data covariance
$C_D$ determines how residuals are weighted. The quantity
\[
    \obs - \G m
\]
is not merely a vector of errors. It is a vector of errors measured relative to
the uncertainty in the observations.

A natural data misfit is therefore
\[
    \frac{1}{2}
    \|\obs-\G m\|_{C_D^{-1}}^2
    =
    \frac{1}{2}
    (\obs-\G m)^T C_D^{-1}(\obs-\G m).
\]

\begin{intuition}
The data space tells us what observations are. The noise covariance tells us
how loudly each component of the data is allowed to complain.
\end{intuition}

\subsection{What is the forward map?}

The third question is:
\[
    \text{How does } m \text{ produce data?}
\]

This is the role of the forward map
\[
    \G : \M \to \D.
\]
In the toy example,
\[
    [\G(m)]_i
    =
    \int_{\Omega} K_i(x)m(x)\,dx,
    \qquad
    i=1,\ldots,K.
\]

The forward map is the mathematical description of the experiment. It tells us
which features of the model are visible to the data. If the forward map is
smoothing, averaging, or compact, then the inverse problem will usually be
ill-posed in a deterministic sense. Small changes in data may correspond to
large changes in some model directions.

Before discretizing, we must know what kind of map $\G$ is. Is it linear? Is it
continuous? Is it bounded between the chosen model and data spaces? Does it
make sense for every $m\in\M$?

For a linear inverse problem, we require at least
\[
    \G(\alpha m_1+\beta m_2)
    =
    \alpha \G m_1 + \beta \G m_2.
\]
For a well-behaved function-space formulation, we usually also want $\G$ to be
continuous:
\[
    m_n \to m \text{ in } \M
    \quad \Longrightarrow \quad
    \G m_n \to \G m \text{ in } \D.
\]

\begin{warning}
A matrix is not the forward operator. A matrix is a representation of the
forward operator after choosing coordinates in the model and data spaces.
\end{warning}

\subsection{What does prior information mean?}

The fourth question is:
\[
    \text{What do we know before observing the data?}
\]

Prior information may include expected amplitude, smoothness, length scale,
boundary behaviour, correlations, sign constraints, or physical plausibility.
In a Bayesian formulation, this information is encoded by a prior probability
measure
\[
    \prior
\]
on the model space $\M$.

Thus, instead of placing a Gaussian directly on coefficients, we write
\[
    m \sim \prior.
\]

If the prior is Gaussian, we may write
\[
    \prior = \mathcal{N}(m_0,C_0),
\]
where $m_0$ is the prior mean and $C_0$ is a covariance operator on the model
space.

The word \emph{operator} matters. In function space, covariance is not merely a
matrix. It is a mathematical object that must act coherently on functions. It
must encode a genuine probability measure on the model space.

\begin{lesson}
Prior information is not a numerical stabilizer sprinkled on top of a matrix
problem. It is part of the definition of the inverse problem.
\end{lesson}

\section{The ingredients of a well-posed Bayesian inverse problem}

We now assemble the pieces. A Bayesian inverse problem should be formulated
before discretization by specifying the following ingredients.

\subsection{Model space}

The model space is the mathematical home of the unknown:
\[
    m \in \M.
\]

The choice of $\M$ determines which objects are admissible models. It also
determines the meaning of convergence, continuity, size, regularity, and
eventually the geometry needed to define adjoints.

Typical choices include Hilbert spaces such as
\[
    L^2(\Omega),
    \qquad
    H^1(\Omega),
    \qquad
    H^s(\Omega),
\]
possibly with boundary conditions built into the space.

We should choose $\M$ by asking modelling questions:
\begin{enumerate}[label=\arabic*.]
    \item What kind of object is physically meaningful?
    \item How much regularity should the model have?
    \item What boundary behaviour is appropriate?
    \item What perturbations should count as small?
    \item What kind of prior measure can live on this space?
\end{enumerate}

\subsection{Data space}

The data space is the mathematical home of the observations:
\[
    \obs \in \D.
\]

In the toy problem,
\[
    \D = \mathbb{R}^K.
\]
This finite-dimensional choice is convenient, but it does not remove all
modelling decisions. We still need to decide how data residuals are measured.

If the observational noise covariance is $C_D$, then the natural data norm is
\[
    \|d\|_{C_D^{-1}}^2
    =
    d^T C_D^{-1}d.
\]
This norm weights residuals according to observational uncertainty. A residual
in a highly uncertain datum should not be punished as severely as an equal-sized
residual in a very accurate datum.

Thus, even when $\D$ is finite-dimensional, the data geometry carries meaning.

\subsection{Forward map}

The forward map connects models to predicted data:
\[
    \G : \M \to \D.
\]

For the linear case,
\[
    \G(\alpha m_1+\beta m_2)
    =
    \alpha \G m_1+\beta \G m_2.
\]
In the integral-kernel example,
\[
    [\G(m)]_i
    =
    \int_\Omega K_i(x)m(x)\,dx.
\]

The forward map must be compatible with the model space. For instance, the
integral above must be well-defined for every $m\in\M$. If the kernels $K_i$
are square-integrable and
\[
    \M = L^2(\Omega),
\]
then each datum is well-defined by the Cauchy-Schwarz inequality:
\[
    \left|
    \int_\Omega K_i(x)m(x)\,dx
    \right|
    \leq
    \|K_i\|_{L^2(\Omega)}
    \|m\|_{L^2(\Omega)}.
\]

This compatibility is exactly the sort of issue that disappears from view if we
start with a matrix.

\subsection{Noise model}

The noise model describes the relationship between the true predicted data and
the observed data:
\[
    \noise \sim \mathcal{N}(0,C_D).
\]

The data equation is
\[
    \obs = \G m + \noise.
\]

Equivalently, conditional on $m$, the data are distributed as
\[
    \obs \mid m
    \sim
    \mathcal{N}(\G m,C_D).
\]

For finite-dimensional data, this gives the likelihood
\[
    \pi(\obs\mid m)
    \propto
    \exp
    \left(
        -
        \frac{1}{2}
        \|\obs-\G m\|_{C_D^{-1}}^2
    \right).
\]

It is useful to define the data-misfit potential
\[
    \Phi(m;\obs)
    =
    \frac{1}{2}
    \|\obs-\G m\|_{C_D^{-1}}^2.
\]
Then
\[
    \pi(\obs\mid m)
    \propto
    \exp(-\Phi(m;\obs)).
\]

\subsection{Prior measure}

The prior measure describes uncertainty about the model before the data are
observed:
\[
    m \sim \prior.
\]

This is where the naive construction failed. We did not choose a measure on
$\M$. We chose matrices
\[
    C_N=\sigma^2 I_N
\]
on coefficient spaces
\[
    \mathbb{R}^N.
\]
Each fixed matrix was a valid covariance matrix in finite dimensions, but the
sequence did not define a coherent function-space prior.

In the function-space formulation, the prior is chosen directly on $\M$:
\[
    \prior \in \mathcal{P}(\M),
\]
where $\mathcal{P}(\M)$ denotes the set of probability measures on $\M$.

For a Gaussian prior, we write
\[
    \prior = \mathcal{N}(m_0,C_0).
\]
Here $C_0$ must be a valid covariance operator for a Gaussian measure on the
chosen model space.

\begin{warning}
A covariance matrix answers a finite-dimensional question.
A covariance operator answers a function-space question.
Confusing the two is how the naive posterior slipped through the trapdoor.
\end{warning}

\section{The posterior measure}

Once the ingredients are specified, the posterior can be written in a way that
does not depend on any particular discretization.

For Gaussian observational noise, define
\[
    \Phi(m;\obs)
    =
    \frac{1}{2}
    \|\obs-\G m\|_{C_D^{-1}}^2.
\]
The posterior measure $\post$ is then formally given by
\[
    \frac{d\post}{d\prior}(m)
    =
    \frac{1}{Z(\obs)}
    \exp(-\Phi(m;\obs)),
\]
where
\[
    Z(\obs)
    =
    \int_{\M}
    \exp(-\Phi(m;\obs))
    \,d\prior(m)
\]
is the normalizing constant.

This formula is one of the cleanest ways to write Bayes' theorem in function
space. It says that the posterior is obtained by reweighting the prior
according to how well each model predicts the observed data.

\begin{intuition}
The prior says which models are possible before the experiment.
The likelihood tilts that probability toward models that explain the data.
The posterior is the tilted measure.
\end{intuition}

For this expression to define a probability measure, we need
\[
    0 < Z(\obs) < \infty.
\]
If this holds, then $\post$ is a well-defined probability measure on the same
model space $\M$ as the prior.

\section{The order matters}

We can now state the central rule.

\begin{lesson}
Choose
\[
    \M,\qquad \D,\qquad \G,\qquad \prior
\]
before choosing a mesh, basis, quadrature rule, or matrix representation.
\end{lesson}

The order matters because each numerical object should approximate a
mathematical object that already exists.

A mesh should approximate the domain. A finite-dimensional subspace should
approximate the model space. A matrix should represent the forward map. A
discrete covariance should approximate the covariance operator. A discrete
posterior should approximate the function-space posterior.

The wrong order is
\[
    \text{basis}
    \quad\Rightarrow\quad
    \text{coefficients}
    \quad\Rightarrow\quad
    \text{matrix prior}
    \quad\Rightarrow\quad
    \text{posterior}.
\]

The better order is
\[
    \text{model space}
    \quad\Rightarrow\quad
    \text{prior measure}
    \quad\Rightarrow\quad
    \text{forward map}
    \quad\Rightarrow\quad
    \text{posterior}
    \quad\Rightarrow\quad
    \text{discretization}.
\]

\begin{checkpoint}
Discretization should enter as an approximation of a problem that already has a
meaning without the discretization.
\end{checkpoint}

\section{Well-posedness}

In deterministic inverse problems, well-posedness is often associated with
three requirements: existence, uniqueness, and stability. In Bayesian inverse
problems, these ideas are modified. We no longer demand that there is a unique
model matching the data. Instead, we ask whether the posterior distribution is
well-defined, meaningful, and stable.

\subsection{Existence}

Existence means that the posterior measure exists.

Using the Radon-Nikodym form
\[
    \frac{d\post}{d\prior}(m)
    =
    \frac{1}{Z(\obs)}
    \exp(-\Phi(m;\obs)),
\]
existence requires the normalizing constant
\[
    Z(\obs)
    =
    \int_{\M}
    \exp(-\Phi(m;\obs))
    \,d\prior(m)
\]
to satisfy
\[
    0 < Z(\obs) < \infty.
\]

The upper bound ensures that the posterior can be normalized. The lower bound
ensures that the data are not assigned zero probability under the model in a way
that destroys the conditioning.

For the linear Gaussian problem with finite-dimensional data, this condition is
usually benign when $\G$ is continuous, $C_D$ is positive definite, and the
prior is a genuine probability measure on $\M$.

\begin{lesson}
Existence means that Bayes' theorem produces an actual probability measure, not
just a decorative formula.
\end{lesson}

\subsection{Uniqueness in the probabilistic sense}

In an infinite-dimensional inverse problem, we should not usually expect a
unique model $m$ that explains the data. There may be infinitely many models
that fit the observations to within noise.

Bayesian uniqueness means something different. It means that, once we have
specified
\[
    \M,\qquad \D,\qquad \G,\qquad C_D,\qquad \prior,
\]
the posterior measure is uniquely determined.

That is, the object of inference is not
\[
    \text{the one true coefficient vector},
\]
but rather
\[
    \post = \mathbb{P}(m\in\cdot \mid \obs).
\]

There may be many plausible models, but there is one posterior distribution
describing their plausibility relative to the chosen assumptions.

\begin{intuition}
Bayesian inversion does not always collapse uncertainty into a single answer.
It gives uncertainty a precise shape.
\end{intuition}

\subsection{Stability with respect to data}

Stability means that small changes in the observed data should produce small
changes in the posterior.

Suppose we observe two nearby data vectors
\[
    \obs,\obs' \in \D.
\]
They produce two posterior measures
\[
    \mu^{\obs}
    \qquad
    \text{and}
    \qquad
    \mu^{\obs'}.
\]
A Bayesian inverse problem is stable if
\[
    \obs \approx \obs'
    \quad\Longrightarrow\quad
    \mu^{\obs} \approx \mu^{\obs'}.
\]

To make this precise, one chooses a distance between probability measures, such
as the Hellinger distance. The exact metric is less important here than the
principle: the posterior should not change violently when the data are perturbed
slightly.

For linear Gaussian inverse problems with finite-dimensional data, stability is
often obtained when the forward map is continuous and the likelihood potential
depends continuously on the data.

\begin{warning}
A stable posterior does not mean the inverse problem has become magically
well-conditioned in the deterministic sense. It means the probability measure
responds continuously to perturbations in the data.
\end{warning}

\subsection{Convergence under discretization}

Finally, the discretized posteriors should converge to the function-space
posterior.

Let
\[
    \M_N \subset \M
\]
be a sequence of finite-dimensional approximation spaces, and let
\[
    \post_N
\]
denote the posterior obtained by discretizing the already-formulated
function-space problem.

A good discretization should satisfy, in an appropriate sense,
\[
    \post_N \to \post
    \qquad
    \text{as}
    \qquad
    N\to\infty.
\]

This is very different from merely defining a new posterior for every $N$.
The limiting posterior $\post$ must already exist. The discretizations are then
judged by whether they approximate it.

The naive construction failed precisely here. It produced a sequence
\[
    \post_1,\post_2,\post_3,\ldots
\]
without first defining a target posterior $\post$ on function space.

\begin{lesson}
Convergence under discretization is only meaningful after the continuous problem
has been posed.
\end{lesson}

\section{The moral of Act III}

We now know what went wrong in the naive workflow. The order was upside down.

We used discretization to define the model space. We used a matrix to define
the prior. We used a transpose without first specifying the geometry. Then we
expected refinement to behave as if there were a continuous problem underneath.

The cure is to reverse the order.

First choose the mathematical ingredients:
\[
    \M,\qquad
    \D,\qquad
    \G,\qquad
    C_D,\qquad
    \prior.
\]
Then define the posterior measure:
\[
    \frac{d\post}{d\prior}(m)
    =
    \frac{1}{Z(\obs)}
    \exp
    \left(
        -
        \frac{1}{2}
        \|\obs-\G m\|_{C_D^{-1}}^2
    \right).
\]
Only after this do we discretize.

\begin{checkpoint}
The next act is a detour through functional analysis. We need to understand
what kinds of spaces are available, what structure they carry, and why Hilbert
spaces so often become the natural habitat for linear Bayesian inverse
problems.
\end{checkpoint}
% ============================================================
% ACT IV
% ============================================================

\part{Act IV: A Guided Detour Through Functional Analysis}

\chapter{A Zoo of Spaces}

\section{Why there are so many spaces}

We have reached the point where the inverse problem can no longer be treated as
a matrix problem wearing a function-shaped hat. We must decide where the unknown
model lives. This forces us into the world of functional analysis.

At first glance, functional analysis looks like a zoo. There are sets,
topological spaces, metric spaces, vector spaces, normed spaces, Banach spaces,
inner-product spaces, Hilbert spaces, Sobolev spaces, spaces of distributions,
spaces of continuous functions, spaces of measures, and many others.

This abundance is not mathematical decoration. Each type of space carries a
different kind of structure. Each structure allows us to ask different
questions.

For inverse problems, the relevant questions include:
\begin{enumerate}[label=\arabic*.]
    \item What objects are admissible models?
    \item What does it mean for two models to be close?
    \item What does it mean for a sequence of models to converge?
    \item Can we add perturbations to a model?
    \item Can we measure the size of a perturbation?
    \item Can we define angles, orthogonality, projections, and adjoints?
    \item Can we place a Gaussian prior on the space?
\end{enumerate}

Different spaces answer these questions with different levels of richness. The
path through this chapter is therefore not a catalogue for its own sake. It is
a guided climb: each step adds just enough structure to make more of the inverse
problem meaningful.

\begin{intuition}
A space is not only a container for objects. It is a rulebook for what we are
allowed to do with those objects.
\end{intuition}

\section{Sets}

\subsection{Membership}

The simplest mathematical structure is a set. A set is just a collection of
objects. If $X$ is a set and $x$ is one of its elements, we write
\[
    x \in X.
\]

For example, we might write
\[
    m \in \M
\]
to say that $m$ is an admissible model.

At the level of sets, the only basic question is membership:
\[
    \text{Is this object allowed or not?}
\]
A set tells us which objects belong to our model universe, but nothing more.

For example, suppose
\[
    \M =
    \{m:\Omega\to\R\}.
\]
This says that every real-valued function on $\Omega$ is admissible. But it
does not tell us whether one function is close to another, whether a sequence
of functions converges, whether small perturbations make sense, or whether
Gaussian probability measures can be placed on $\M$.

\subsection{No geometry yet}

A bare set has no geometry. There is no distance, no notion of size, no notion
of angle, and no notion of convergence.

If $m_1,m_2\in \M$, a set alone does not allow us to ask:
\[
    \text{Is } m_1 \text{ close to } m_2?
\]
Nor does it allow us to ask:
\[
    m_n \to m?
\]

This is too little structure for inverse problems. In inverse problems, we need
to compare models, perturb models, approximate models, and study the behaviour
of sequences of approximations.

\begin{lesson}
A set tells us what objects are allowed. It does not tell us how those objects
relate to one another.
\end{lesson}

\section{Topological spaces}

\subsection{Open sets}

The next layer of structure is topology. A topological space is a set $X$
together with a collection of subsets called open sets. We write this collection
as
\[
    \mathcal{T}.
\]
The pair
\[
    (X,\mathcal{T})
\]
is called a topological space.

The open sets tell us what it means for points to be near one another, but in a
very flexible way. A topology does not necessarily come from a distance. It only
tells us which subsets count as neighbourhoods.

The collection $\mathcal{T}$ must satisfy:
\begin{enumerate}[label=\arabic*.]
    \item $\varnothing \in \mathcal{T}$ and $X\in\mathcal{T}$;
    \item arbitrary unions of open sets are open;
    \item finite intersections of open sets are open.
\end{enumerate}

These rules are abstract, but the goal is familiar: topology lets us talk about
continuity and convergence without first committing to a particular formula for
distance.

\subsection{Continuity}

If $X$ and $Y$ are topological spaces, a map
\[
    F:X\to Y
\]
is continuous if the inverse image of every open set in $Y$ is open in $X$.

That is, if $U\subset Y$ is open, then
\[
    F^{-1}(U)
    =
    \{x\in X : F(x)\in U\}
\]
must be open in $X$.

For inverse problems, continuity matters because the forward map should not
behave like a mathematical trapdoor. If two models are close, their predicted
data should also be close:
\[
    m_n \to m
    \quad\Longrightarrow\quad
    \G m_n \to \G m.
\]

This kind of statement requires at least a topology on the model space and the
data space.

\subsection{Convergence without distances}

A topology gives a notion of convergence. We can say that a sequence
\[
    x_1,x_2,x_3,\ldots
\]
converges to $x$ if, eventually, the sequence lies in every neighbourhood of
$x$.

Symbolically,
\[
    x_n \to x.
\]

The important point is that convergence can be discussed without a distance
function. This is useful in many areas of analysis, especially when dealing with
weak convergence, distributions, and probability measures.

However, for the present story, topology alone is still too sparse. We will
want to measure sizes of perturbations, define covariance operators, and discuss
Gaussian priors. For this, it is usually helpful to have more structure.

\begin{intuition}
Topology tells us when things approach each other, but it does not necessarily
tell us how far apart they are.
\end{intuition}

\section{Metric spaces}

\subsection{Distances}

A metric space is a set $X$ equipped with a distance function
\[
    d:X\times X\to [0,\infty).
\]
For $x,y,z\in X$, the metric must satisfy:
\begin{enumerate}[label=\arabic*.]
    \item $d(x,y)\geq 0$;
    \item $d(x,y)=0$ if and only if $x=y$;
    \item $d(x,y)=d(y,x)$;
    \item $d(x,z)\leq d(x,y)+d(y,z)$.
\end{enumerate}

The last condition is the triangle inequality. It says that going from $x$ to
$z$ directly is no longer than going through an intermediate point $y$.

A metric allows us to quantify closeness:
\[
    d(x,y) \ll 1.
\]

In inverse problems, metrics let us ask whether reconstructed models are close
to one another, whether approximations converge, and whether the posterior
changes continuously when the data are perturbed.

\subsection{Convergence by distance}

In a metric space, convergence has a very concrete meaning:
\[
    x_n \to x
\]
if
\[
    d(x_n,x)\to 0.
\]

This is often easier to work with than purely topological convergence.

For example, if $X=C([0,L])$, the space of continuous functions on $[0,L]$, we
may define the uniform distance
\[
    d(f,g)
    =
    \sup_{x\in[0,L]} |f(x)-g(x)|.
\]
Then
\[
    f_n \to f
\]
means that $f_n$ converges uniformly to $f$.

Alternatively, if $X=L^2(\Omega)$, we may define
\[
    d(f,g)
    =
    \left(
        \int_\Omega |f(x)-g(x)|^2\,dx
    \right)^{1/2}.
\]
Then convergence means convergence in mean-square size.

These two notions of convergence are different. Choosing a space means choosing
which distinctions matter.

\begin{warning}
There is no single universal notion of closeness between functions. The chosen
space decides what kind of closeness is being measured.
\end{warning}

\section{Vector spaces}

\subsection{Addition and scalar multiplication}

A vector space is a set $V$ whose elements can be added and multiplied by
scalars. If $u,v\in V$ and $\alpha,\beta\in\R$, then
\[
    \alpha u+\beta v \in V.
\]

This is the basic algebraic structure behind linear inverse problems.

Examples include:
\[
    \R^N,
    \qquad
    L^2(\Omega),
    \qquad
    C([0,L]),
    \qquad
    H^1(\Omega).
\]

If $m_1$ and $m_2$ are models, vector space structure lets us form
\[
    m_1+m_2,
    \qquad
    \alpha m_1,
    \qquad
    m_1-m_2.
\]

The difference
\[
    m_1-m_2
\]
is especially important. It represents a perturbation, an update, or an error.

\subsection{Perturbations}

In inverse problems, we often reason in terms of perturbations:
\[
    m = m_0+\delta m.
\]
Here $m_0$ is a reference model and $\delta m$ is an update.

This expression only makes sense if the model space has a linear structure.
We need to be able to add a perturbation to a model and still obtain another
admissible model:
\[
    m_0,\delta m\in\M
    \quad\Longrightarrow\quad
    m_0+\delta m\in\M.
\]

Linear structure is also what allows us to describe Gaussian priors by adding a
random perturbation to a mean:
\[
    m = m_0 + \xi,
\]
where $\xi$ is a centred random element.

\begin{intuition}
Vector spaces are natural when models can be nudged, corrected, averaged, and
superposed.
\end{intuition}

\subsection{Linear inverse problems}

A linear inverse problem is built around a linear forward map
\[
    \G:\M\to\D.
\]
Linearity means
\[
    \G(\alpha m_1+\beta m_2)
    =
    \alpha \G m_1+\beta \G m_2.
\]

This structure is powerful. It allows us to understand how perturbations in the
model produce perturbations in the data:
\[
    \delta y = \G \delta m.
\]

It also allows Gaussian priors and Gaussian noise to produce Gaussian
posteriors, provided the rest of the formulation is coherent.

But vector spaces alone are still not enough. They let us add models, but they
do not tell us how large a model is, how far apart two models are, or what it
means for an infinite sequence of approximations to converge. For that, we need
norms.

\section{Manifolds}

\subsection{When the unknown does not live in a vector space}

Not every meaningful model space is naturally a vector space. Sometimes the
unknown lives on a curved space called a manifold.

Examples include:
\begin{enumerate}[label=\arabic*.]
    \item directions on a sphere;
    \item rotations in three dimensions;
    \item positive-definite matrices;
    \item shapes or interfaces;
    \item diffeomorphisms;
    \item fields constrained to have fixed magnitude.
\end{enumerate}

In such cases, adding two models may not make sense. If $m_1$ and $m_2$ are
points on a sphere, then
\[
    m_1+m_2
\]
need not lie on the sphere. The model space is curved, not linear.

\subsection{Tangent spaces}

Even when the model space is a manifold, local perturbations often live in a
linear space called a tangent space.

If $m$ is a point on a manifold $\mathcal{M}$, the tangent space at $m$ is
written
\[
    T_m\mathcal{M}.
\]
Elements of $T_m\mathcal{M}$ represent infinitesimal perturbations around $m$.

Thus, a manifold is not globally a vector space, but it has local linear
approximations. This is why optimization and inverse problems on manifolds often
involve tangent spaces, exponential maps, and local coordinates.

\subsection{Why we do not follow this path here}

Manifold-valued inverse problems are important, but they require additional
machinery. One must keep track of curvature, local charts, tangent spaces, and
how probability measures behave on nonlinear spaces.

In this document, we will not go in that direction. Our goal is already rich
enough: to understand why even linear inverse problems can go wrong if we
discretize before choosing the function-space formulation.

\begin{remark}
Manifolds are important, but this document stays in the linear world.
Our unknowns will live in vector spaces.
\end{remark}

% ------------------------------------------------------------

\chapter{Adding Structure One Layer at a Time}

\section{From set to vector space}

We now rebuild the model space slowly. The aim is to understand why each layer
of structure is useful and what it buys us.

A bare set tells us which objects are allowed. A vector space tells us that
objects can be added, scaled, and perturbed. This is the first major step toward
linear inverse problems.

\subsection{Why linear combinations matter}

Suppose $m_1$ and $m_2$ are admissible models. If the model space is a vector
space, then
\[
    \alpha m_1+\beta m_2
\]
is also an admissible model for every $\alpha,\beta\in\R$.

This allows us to construct approximations by basis expansion:
\[
    m_N
    =
    \sum_{j=1}^N u_j\phi_j.
\]

It also allows us to form averages:
\[
    \bar{m}
    =
    \mathbb{E}[m],
\]
and fluctuations around an average:
\[
    m-\bar{m}.
\]

These operations are essential in Gaussian Bayesian inversion. The mean,
covariance, posterior update, and adjoint machinery all lean on linear
structure.

\subsection{Perturbation thinking}

Many inverse problems are naturally phrased in terms of a reference model plus
a perturbation:
\[
    m = m_0+\delta m.
\]

The data predicted by the perturbed model are
\[
    \G(m_0+\delta m).
\]
If $\G$ is linear, then
\[
    \G(m_0+\delta m)
    =
    \G m_0+\G\delta m.
\]

The perturbation $\delta m$ is the object that the data update. This way of
thinking is particularly natural when prior information gives us a background
model $m_0$ and the inverse problem seeks deviations from it.

\begin{lesson}
Linear spaces are the natural habitat of additive perturbations.
If our inference problem is about updates, differences, and superposition,
linear structure is not decoration. It is the grammar.
\end{lesson}

\section{From vector space to normed space}

A vector space lets us add and scale. But it still does not tell us how large a
model is. To measure size, we introduce a norm.

A normed space is a vector space $V$ equipped with a function
\[
    \|\cdot\|:V\to[0,\infty)
\]
such that, for all $u,v\in V$ and $\alpha\in\R$,
\begin{enumerate}[label=\arabic*.]
    \item $\|u\|\geq 0$;
    \item $\|u\|=0$ if and only if $u=0$;
    \item $\|\alpha u\|=|\alpha|\,\|u\|$;
    \item $\|u+v\|\leq \|u\|+\|v\|$.
\end{enumerate}

\subsection{Measuring size}

The norm of a model,
\[
    \|m\|,
\]
measures its size according to the geometry of the chosen space.

For example, in $L^2(\Omega)$,
\[
    \|m\|_{L^2}
    =
    \left(
        \int_\Omega |m(x)|^2\,dx
    \right)^{1/2}.
\]
This norm measures the root-mean-square size of a function.

In $H^1(\Omega)$,
\[
    \|m\|_{H^1}
    =
    \left(
        \int_\Omega |m(x)|^2\,dx
        +
        \int_\Omega |m'(x)|^2\,dx
    \right)^{1/2},
\]
at least in one dimension for sufficiently regular functions. This norm
measures both the size of the function and the size of its derivative.

Thus, choosing a norm is not neutral. It decides which features of a model are
expensive.

\subsection{Convergence}

A norm automatically defines a distance:
\[
    d(m_1,m_2)
    =
    \|m_1-m_2\|.
\]
Therefore, it gives a notion of convergence:
\[
    m_n \to m
\]
if
\[
    \|m_n-m\|\to 0.
\]

This is crucial for discretization. When we say that finite-dimensional
approximations converge to a function-space model, we must mean convergence in
some norm:
\[
    \|m_N-m\|_{\M}\to 0.
\]

Different norms give different meanings to this statement. Convergence in
$L^2$ means average-square convergence. Convergence in $H^1$ also requires
convergence of derivatives. Uniform convergence requires pointwise closeness
everywhere.

\begin{warning}
A discretization cannot be called convergent until we say which norm or topology
defines convergence.
\end{warning}

\section{From normed space to Banach space}

Normed spaces are useful, but one more property is often essential:
completeness.

A complete normed vector space is called a Banach space.

\subsection{Cauchy sequences}

A sequence $(m_n)$ in a normed space is called Cauchy if, for every
$\varepsilon>0$, there exists an index $N_\varepsilon$ such that
\[
    \|m_n-m_k\| < \varepsilon
\]
whenever
\[
    n,k\geq N_\varepsilon.
\]

In words, the elements of the sequence eventually become arbitrarily close to
one another.

A convergent sequence is always Cauchy. If
\[
    m_n\to m,
\]
then for large $n$ and $k$, both $m_n$ and $m_k$ are close to the same limit
$m$, so they are close to each other.

The converse is not always true. A Cauchy sequence may be trying to converge to
an object that is missing from the space.

\subsection{Completeness}

A normed space is complete if every Cauchy sequence has a limit in the space.

That is, whenever $(m_n)$ is Cauchy, there exists some $m$ in the same space
such that
\[
    \|m_n-m\|\to 0.
\]

Completeness prevents sequences from slipping out of the space through invisible
cracks. This matters because analysis is full of limiting processes:
\[
    N\to\infty,
    \qquad
    h\to 0,
    \qquad
    n\to\infty.
\]

If our model space is not complete, then a sequence of increasingly refined
approximations may converge to something outside the space. That is awkward at
best and disastrous at worst.

\subsection{Why completeness prevents missing limits}

Consider a numerical method that produces approximations
\[
    m_1,m_2,m_3,\ldots.
\]
If the method is behaving well, these approximations should get closer and
closer to one another. In other words, they should form a Cauchy sequence.

But if the model space is not complete, there may be no model $m$ in the space
such that
\[
    m_N\to m.
\]

A Banach space rules out this kind of missing-limit behaviour.

\begin{intuition}
Completeness means that if a sequence behaves as though it is converging, the
limit has not been exiled from the space.
\end{intuition}

\section{From Banach space to inner-product space}

A Banach space gives us size and convergence. But for many parts of linear
inverse theory, we want more: angles, orthogonality, projections, and adjoints.

These come from an inner product.

An inner product on a real vector space $V$ is a map
\[
    \langle \cdot,\cdot\rangle:V\times V\to\R
\]
such that, for $u,v,w\in V$ and $\alpha,\beta\in\R$,
\begin{enumerate}[label=\arabic*.]
    \item $\langle u,v\rangle=\langle v,u\rangle$;
    \item $\langle \alpha u+\beta v,w\rangle
    =
    \alpha\langle u,w\rangle+\beta\langle v,w\rangle$;
    \item $\langle u,u\rangle\geq 0$;
    \item $\langle u,u\rangle=0$ if and only if $u=0$.
\end{enumerate}

An inner product defines a norm by
\[
    \|u\|
    =
    \sqrt{\langle u,u\rangle}.
\]

\subsection{Angles}

In Euclidean space, the inner product determines angles:
\[
    \langle u,v\rangle
    =
    \|u\|\|v\|\cos\theta.
\]

In function spaces, the same idea persists. The inner product measures how much
two functions align.

For example, in $L^2(\Omega)$,
\[
    \langle f,g\rangle_{L^2}
    =
    \int_\Omega f(x)g(x)\,dx.
\]
If this quantity is large and positive, the functions tend to be positive and
negative in the same regions. If it is near zero, they are largely uncorrelated
in the $L^2$ geometry.

\subsection{Orthogonality}

Two elements $u$ and $v$ are orthogonal if
\[
    \langle u,v\rangle = 0.
\]

Orthogonality is indispensable in approximation theory. If $\M_N$ is a
finite-dimensional subspace of $\M$, the best approximation to $m\in\M$ in
$\M_N$ is often characterized by an orthogonality condition:
\[
    \langle m-P_Nm,v_N\rangle_{\M}=0
    \qquad
    \text{for all } v_N\in\M_N.
\]
Here $P_Nm$ is the projection of $m$ onto $\M_N$.

This is one reason Hilbert spaces are so useful for numerical analysis. They
allow approximation to be expressed geometrically.

\subsection{Projection}

A projection maps an object onto a subspace. In a Hilbert-type geometry, the
orthogonal projection $P_Nm$ is the element of $\M_N$ closest to $m$:
\[
    P_Nm
    =
    \operatorname*{arg\,min}_{v_N\in\M_N}
    \|m-v_N\|_{\M}.
\]

The error
\[
    m-P_Nm
\]
is orthogonal to the approximation space:
\[
    \langle m-P_Nm,v_N\rangle_{\M}=0
    \qquad
    \forall v_N\in\M_N.
\]

Projection is one of the cleanest ways to connect continuous and discrete
problems. Rather than replacing the continuous problem with a coefficient
problem, we approximate continuous objects by projecting them into
finite-dimensional subspaces.

\section{From inner-product space to Hilbert space}

An inner-product space becomes a Hilbert space when it is complete with respect
to the norm induced by the inner product.

Thus, a Hilbert space is an inner-product space in which every Cauchy sequence
converges.

\subsection{Completeness again}

Completeness returns because infinite-dimensional analysis is full of limits.
We want limits of approximations, limits of basis expansions, limits of
posterior samples, and limits of discretized covariance operators to remain
inside the space.

If $\Hh$ is a Hilbert space and
\[
    m_n
\]
is Cauchy in the Hilbert norm, then there exists
\[
    m\in\Hh
\]
such that
\[
    \|m_n-m\|_{\Hh}\to 0.
\]

This is the analytical floor beneath the machinery. Without it, sequences may
wander off through the walls.

\subsection{The Hilbert-space toolkit}

Hilbert spaces are especially useful because they provide:
\begin{enumerate}[label=\arabic*.]
    \item a norm for measuring size;
    \item a topology for discussing convergence;
    \item an inner product for geometry;
    \item orthogonality;
    \item projections onto closed subspaces;
    \item adjoints of bounded linear operators;
    \item spectral theory for many important operators;
    \item a natural language for covariance operators;
    \item a convenient setting for Gaussian measures.
\end{enumerate}

For linear inverse problems, the adjoint is particularly important. If
\[
    \G:\M\to\D
\]
is a bounded linear map between Hilbert spaces, its adjoint
\[
    \Gadj:\D\to\M
\]
is defined by
\[
    \langle \G m,d\rangle_{\D}
    =
    \langle m,\Gadj d\rangle_{\M}.
\]

This definition depends on the inner products of both spaces. It is not merely
a transpose. The transpose appears only after coordinates and Euclidean
geometries have been chosen.

\begin{lesson}
Hilbert spaces are not the only possible model spaces, but they are often the
simplest setting where geometry, projection, adjoints, and Gaussian priors can
speak the same language.
\end{lesson}

% ------------------------------------------------------------

\chapter{Hilbert Spaces We Actually Use}

\section{$L^2$ spaces}

The first Hilbert space that appears almost everywhere in inverse problems is
$L^2(\Omega)$.

\subsection{Square-integrable functions}

The space $L^2(\Omega)$ consists of square-integrable functions:
\[
    L^2(\Omega)
    =
    \left\{
        f : \Omega \to \R
        :
        \int_\Omega |f(x)|^2\,dx < \infty
    \right\}.
\]

Strictly speaking, elements of $L^2(\Omega)$ are equivalence classes of
functions that agree almost everywhere. This means that changing a function at
a set of points of measure zero does not produce a different $L^2$ element.

The inner product is
\[
    \langle f,g\rangle_{L^2}
    =
    \int_\Omega f(x)g(x)\,dx.
\]
The corresponding norm is
\[
    \|f\|_{L^2}
    =
    \left(
        \int_\Omega |f(x)|^2\,dx
    \right)^{1/2}.
\]

\subsection{What $L^2$ encodes}

The $L^2$ norm measures average squared amplitude. If
\[
    \|f\|_{L^2}
\]
is finite, then the function has finite energy in an integral sense.

This is useful when:
\begin{enumerate}[label=\arabic*.]
    \item pointwise values are not essential;
    \item average magnitude matters more than exact pointwise behaviour;
    \item the forward map involves integrals against sensitivity kernels;
    \item the natural data are weighted averages of the model.
\end{enumerate}

For the toy inverse problem,
\[
    [\G(m)]_i
    =
    \int_\Omega K_i(x)m(x)\,dx,
\]
the $L^2$ setting is natural if
\[
    K_i\in L^2(\Omega).
\]
Then, by the Cauchy-Schwarz inequality,
\[
    |[\G(m)]_i|
    \leq
    \|K_i\|_{L^2}
    \|m\|_{L^2}.
\]
So each datum is a bounded linear functional on $L^2(\Omega)$.

\subsection{What $L^2$ does not encode}

The space $L^2(\Omega)$ does not directly encode smoothness. A function may be
in $L^2$ while oscillating violently, having jumps, or lacking classical
derivatives.

Also, pointwise values are not fundamental in $L^2$. Since functions are
identified when they agree almost everywhere, the expression
\[
    f(x_0)
\]
is not generally a well-defined operation on $L^2(\Omega)$.

This matters if the data are point samples:
\[
    y_i = m(x_i)+\eta_i.
\]
Point evaluation is not a continuous functional on $L^2(\Omega)$. Therefore,
if the data require pointwise values, $L^2$ may not be the right model space.

\begin{warning}
Choosing $L^2$ says that average-square size matters. It does not say that
functions are smooth, continuous, or pointwise meaningful.
\end{warning}

\section{Sobolev spaces}

Sobolev spaces add controlled derivatives to the story. They allow us to encode
regularity while still working with a Hilbert-space structure.

\subsection{Weak derivatives}

Classical derivatives require functions to be differentiable pointwise. Sobolev
spaces use a more flexible notion called weak derivatives.

The idea is that a function $f$ has a weak derivative $g$ if integration by
parts behaves as though $g$ were the derivative of $f$. In one dimension, this
means
\[
    \int_\Omega f(x)\varphi'(x)\,dx
    =
    -
    \int_\Omega g(x)\varphi(x)\,dx
\]
for every sufficiently smooth test function $\varphi$ with compact support.

If such a $g$ exists, we write
\[
    f' = g
\]
in the weak sense.

Weak derivatives allow us to discuss derivatives of functions that may not be
classically differentiable everywhere.

\subsection{Smoothness as square-integrability of derivatives}

For an integer $k\geq 0$, the Sobolev space $H^k(\Omega)$ is
\[
    H^k(\Omega)
    =
    \left\{
        f \in L^2(\Omega)
        :
        D^\alpha f \in L^2(\Omega),
        \ |\alpha|\leq k
    \right\}.
\]

In one dimension, this can be written more simply as
\[
    H^k(\Omega)
    =
    \left\{
        f\in L^2(\Omega)
        :
        f',f'',\ldots,f^{(k)}\in L^2(\Omega)
    \right\},
\]
where the derivatives are understood weakly.

The $H^k$ norm is
\[
    \|f\|_{H^k}^2
    =
    \sum_{|\alpha|\leq k}
    \|D^\alpha f\|_{L^2}^2.
\]

For example,
\[
    \|f\|_{H^1}^2
    =
    \|f\|_{L^2}^2
    +
    \|f'\|_{L^2}^2.
\]

\subsection{What Sobolev spaces encode}

Sobolev spaces encode smoothness by requiring derivatives to be square
integrable.

The choice of $k$ controls how much regularity is required:
\[
    H^0(\Omega)=L^2(\Omega),
\]
while $H^1(\Omega)$ requires one weak derivative in $L^2$, $H^2(\Omega)$
requires two, and so on.

This is extremely useful for inverse problems because prior information often
contains statements like:
\begin{enumerate}[label=\arabic*.]
    \item the model should not oscillate too rapidly;
    \item the model should vary smoothly with position;
    \item sharp jumps are unlikely;
    \item derivatives should have bounded energy.
\end{enumerate}

Such statements are naturally expressed using Sobolev-type norms.

\subsection{Smoothness, roughness, and regularity}

The larger $k$ is, the smoother the functions in $H^k(\Omega)$ must be. Thus,
choosing a Sobolev space is a modelling decision about regularity.

For example, in one dimension, functions in $H^1(\Omega)$ have more structure
than arbitrary $L^2$ functions. They have weak derivatives in $L^2$, and they
possess representatives with better continuity properties.

Higher-order Sobolev spaces encode stronger smoothness. Fractional Sobolev
spaces $H^s(\Omega)$, with non-integer $s$, allow intermediate degrees of
regularity.

In Bayesian inversion, this connects directly to the prior. A prior covariance
operator can be chosen so that samples have a particular Sobolev regularity.
This lets us say not merely that the model is random, but what kind of random
function it is expected to be.

\begin{intuition}
Sobolev spaces turn smoothness into bookkeeping. The derivatives that must be
square-integrable are the entries in the ledger.
\end{intuition}

\section{Boundary conditions}

Boundary conditions are often treated as numerical details. In a function-space
formulation, they are part of the mathematical model.

\subsection{Boundary conditions as part of the space}

Suppose
\[
    \Omega=[0,L].
\]
The space
\[
    H^1(\Omega)
\]
contains functions with one weak derivative in $L^2$. But we may want to impose
conditions at the boundary, such as
\[
    f(0)=0
\]
or
\[
    f(L)=0.
\]

This leads to subspaces such as
\[
    H_0^1(\Omega),
\]
which is commonly understood as the closure of smooth compactly supported
functions in the $H^1$ norm. Informally, it encodes zero boundary values.

Boundary conditions restrict which functions are admissible. Therefore, they
are part of the model space.

\subsection{Dirichlet conditions}

A Dirichlet condition fixes the value of the function at the boundary. For
example,
\[
    f(0)=0,
    \qquad
    f(L)=0.
\]

Such conditions may represent physical constraints. For instance, a perturbation
might be required to vanish at a boundary because the boundary value is known
exactly.

If the model space includes Dirichlet conditions, then every admissible model
must satisfy them. A discretization must respect this, and a prior measure must
place probability one on functions satisfying the same condition.

\subsection{Neumann conditions}

A Neumann condition fixes the derivative normal to the boundary. In one
dimension, this may look like
\[
    f'(0)=0,
    \qquad
    f'(L)=0.
\]

Such conditions may encode zero flux, symmetry, or a preference for flatness at
the boundary.

Neumann conditions often arise naturally through differential operators such as
the Laplacian. If a prior covariance is built from an operator like
\[
    \kappa^2 I - \Delta,
\]
then the boundary conditions used to define $\Delta$ affect the covariance, and
therefore affect the prior measure.

\subsection{Mixed boundary conditions}

Mixed boundary conditions combine different types of behaviour at different
parts of the boundary. For example, on the interval $[0,L]$, one might impose
\[
    f(0)=0,
    \qquad
    f'(L)=0.
\]

This is a Dirichlet condition at one end and a Neumann condition at the other.

Such choices can have visible effects on posterior samples, especially near the
boundary. If the prior says that functions must vanish at one end, samples will
behave differently than if the prior says that derivatives vanish there.

\begin{warning}
Boundary conditions are not numerical decoration.
They are part of the mathematical model and, often, part of the prior.
\end{warning}

\section{Choosing a model space}

We now return to the practical question:
\[
    \text{How should we choose } \M?
\]

There is no universal answer. The model space is a modelling decision. It must
reflect the physical problem, the data, the forward map, and the kind of prior
information we are willing to assert.

\subsection{What physical quantities should be finite?}

The first question is:
\[
    \text{What should have finite size?}
\]

If we only require the model itself to have finite average-square amplitude,
then $L^2(\Omega)$ may be natural:
\[
    \int_\Omega |m(x)|^2\,dx < \infty.
\]

If we also require the derivative to have finite average-square amplitude, then
$H^1(\Omega)$ may be more appropriate:
\[
    \int_\Omega |m(x)|^2\,dx
    +
    \int_\Omega |m'(x)|^2\,dx
    <\infty.
\]

If curvature or higher derivatives matter, then one may consider $H^2$ or
higher Sobolev spaces.

The chosen norm should correspond to a meaningful notion of model size.

\subsection{How smooth should the model be?}

The second question is:
\[
    \text{How regular should plausible models be?}
\]

If the model represents a quantity expected to vary smoothly in space, then the
prior and model space should reflect that. If sharp interfaces are possible,
then a highly smooth Sobolev space may be inappropriate.

There is a subtle point here. The model space and the prior are related, but
they are not identical. The model space defines the ambient mathematical
habitat. The prior measure says which parts of that habitat are probable.

For Gaussian priors, sample regularity is controlled by the covariance
operator. A prior may live on a rougher or smoother space depending on the decay
of the covariance eigenvalues.

\subsection{What boundary behaviour is meaningful?}

The third question is:
\[
    \text{What should happen at the boundary?}
\]

Possible answers include:
\begin{enumerate}[label=\arabic*.]
    \item the model value is known at the boundary;
    \item the perturbation should vanish at the boundary;
    \item the derivative should vanish at the boundary;
    \item the boundary should be periodic;
    \item no explicit boundary condition is justified.
\end{enumerate}

Each choice changes the model space and can change the prior covariance. If the
covariance is defined through a differential operator, then the boundary
conditions are part of the operator definition.

This is why boundary conditions cannot be postponed until the implementation.
They affect the mathematical identity of the prior.

\subsection{What should count as a small perturbation?}

Finally, we must ask:
\[
    \text{When is } \delta m \text{ small?}
\]

In an $L^2$ geometry,
\[
    \|\delta m\|_{L^2}
\]
is small when the perturbation has small average-square amplitude.

In an $H^1$ geometry,
\[
    \|\delta m\|_{H^1}
\]
is small only when both the perturbation and its derivative are small in
$L^2$.

Thus, the chosen geometry determines which perturbations are considered mild
and which are considered large. This matters for regularization, prior
modelling, posterior covariance, and adjoints.

\begin{lesson}
Choosing a model space is choosing what kind of object the unknown is, what
kind of perturbations are meaningful, and what kind of uncertainty can be
placed on the problem.
\end{lesson}

\section{The moral of Act IV}

The functional-analytic detour was not a luxury expedition. It was necessary
cartography.

We began with sets, which only provide membership. We added topology, which
allows continuity and convergence. We added metrics and norms, which measure
distance and size. We added completeness, producing Banach spaces. We added
inner products, giving angles, orthogonality, projections, and adjoints. With
completeness and an inner product together, we arrived at Hilbert spaces.

For linear Bayesian inverse problems, Hilbert spaces are often the sweet spot.
They are structured enough to support the geometry we need, but familiar enough
to connect cleanly to computation.

\begin{checkpoint}
The next act turns from the model space to the data space and the forward
operator. In particular, we will see why the adjoint $\Gadj$ is determined by
the inner products on both spaces, and why replacing it by a transpose can
quietly change the problem.
\end{checkpoint}
% ============================================================
% ACT V
% ============================================================

\part{Act V: The Data Space and the Adjoint}

\chapter{The Data Space}

\section{Finite-dimensional data}

In the toy inverse problem, the data consist of finitely many real numbers:
\[
    \obs =
    \begin{pmatrix}
        y_1\\
        \vdots\\
        y_K
    \end{pmatrix}.
\]
Thus, the natural data space is
\[
    \D = \R^K.
\]

This is already a major simplification. The model may live in an
infinite-dimensional function space, but the observations live in a
finite-dimensional vector space. The forward map therefore compresses a model
into finitely many numbers:
\[
    \G : \M \to \R^K.
\]

For example, if
\[
    [\G m]_i
    =
    \int_\Omega K_i(x)m(x)\,dx,
    \qquad
    i=1,\ldots,K,
\]
then the $i$th component of the data is a weighted average of the model, plus
noise.

The data equation is
\[
    \obs = \G m + \noise,
\]
where
\[
    \noise \sim \mathcal{N}(0,C_D).
\]

At first glance, finite-dimensional data seem completely harmless. We know how
to work with vectors in $\R^K$. We know how to multiply matrices. We know how to
write Gaussian densities. But even here, a modelling choice remains hidden:
we must decide how to measure lengths, angles, and residuals in the data space.

\begin{intuition}
The data space may be finite-dimensional, but it still has geometry.
Ignoring that geometry is how a transpose sneaks into places where an adjoint
was supposed to stand.
\end{intuition}

\section{Why finite dimensions are easier}

Finite-dimensional spaces are friendly creatures. Unlike function spaces, they
do not usually make us worry about missing limits, inequivalent norms, or
pathological convergence. This is one reason inverse problems are often
presented first in finite-dimensional form.

\subsection{All norms induce the same topology}

On $\R^K$, many different norms can be defined. For example,
\[
    \|d\|_1
    =
    \sum_{i=1}^K |d_i|,
\]
\[
    \|d\|_2
    =
    \left(
        \sum_{i=1}^K |d_i|^2
    \right)^{1/2},
\]
and
\[
    \|d\|_\infty
    =
    \max_{1\leq i\leq K}|d_i|.
\]

In finite dimensions, all norms induce the same topology. This means that they
give the same notion of convergence:
\[
    d_n \to d
    \quad
    \text{in one norm}
    \qquad
    \Longleftrightarrow
    \qquad
    d_n \to d
    \quad
    \text{in any other norm}.
\]

This is not true in infinite-dimensional spaces. In function spaces, the choice
of norm can radically change what convergence means. But in $\R^K$, convergence
is robust across norms.

\begin{remark}
Although all norms on $\R^K$ induce the same topology, they do not define the
same geometry. They may measure length and angle differently. For inverse
problems, this distinction matters.
\end{remark}

\subsection{Completeness is automatic}

Every finite-dimensional normed vector space is complete. Therefore, if
$(d_n)$ is a Cauchy sequence in $\R^K$, then there exists some
\[
    d\in\R^K
\]
such that
\[
    d_n \to d.
\]

This means that data-space limits do not escape the space. There is no
finite-dimensional version of a sequence of approximations converging to an
object outside the data space.

This makes the data side of the inverse problem much easier than the model
side. The model may require careful functional analysis, but the data live in a
compact little apartment with reliable plumbing.

\subsection{Linear maps are continuous}

Another simplification is that every linear map between finite-dimensional
normed spaces is continuous.

In particular, if
\[
    A:\R^N\to\R^K
\]
is linear, then there exists a constant $c>0$ such that
\[
    \|Au\|_{\R^K}
    \leq
    c\|u\|_{\R^N}
\]
for all $u\in\R^N$.

When the model space is infinite-dimensional, this automatic continuity is no
longer guaranteed. A linear map
\[
    \G:\M\to\D
\]
must be checked. It must make sense on the chosen model space, and it must be
continuous if we want stable dependence of predicted data on the model.

Since $\D=\R^K$ is finite-dimensional, checking continuity of $\G$ often reduces
to checking the continuity of its individual components:
\[
    \G_i:\M\to\R,
    \qquad
    i=1,\ldots,K.
\]

For example, if
\[
    \G_i m
    =
    \int_\Omega K_i(x)m(x)\,dx,
\]
and if
\[
    \M=L^2(\Omega),
    \qquad
    K_i\in L^2(\Omega),
\]
then
\[
    |\G_i m|
    \leq
    \|K_i\|_{L^2}\|m\|_{L^2}.
\]
Thus, each $\G_i$ is continuous on $L^2(\Omega)$.

\section{But geometry still matters}

Finite-dimensional data are easy topologically, but geometry still matters.
The data space is not just a box of $K$ numbers. It has an inner product, and
that inner product tells us how to measure residuals, define adjoints, and
interpret noise.

\subsection{Choosing an inner product}

A general inner product on $\R^K$ can be written as
\[
    \langle a,b\rangle_{\D}
    =
    a^T W_D b,
\]
where
\[
    W_D\in\R^{K\times K}
\]
is symmetric positive definite.

If
\[
    W_D=I_K,
\]
then this is the usual Euclidean inner product:
\[
    \langle a,b\rangle_{\D}
    =
    a^T b.
\]

If $W_D$ is not the identity, then different data directions are weighted
differently. The norm induced by this inner product is
\[
    \|a\|_{\D}
    =
    \sqrt{\langle a,a\rangle_{\D}}
    =
    \sqrt{a^T W_D a}.
\]

This choice affects the adjoint of the forward map. In particular, if the model
space also has a non-Euclidean geometry, then the adjoint is not represented by
the ordinary matrix transpose.

\begin{warning}
In finite dimensions, the transpose and the adjoint coincide only when the
coordinate bases are orthonormal for the inner products being used.
\end{warning}

\subsection{Noise covariance and data weighting}

The observational noise is modelled as
\[
    \noise \sim \mathcal{N}(0,C_D).
\]
Here $C_D$ describes the covariance of the noise in data coordinates.

If the data are independent with equal variance, then
\[
    C_D = \sigma_D^2 I_K.
\]
If the data have different variances, then $C_D$ may be diagonal:
\[
    C_D =
    \operatorname{diag}(\sigma_1^2,\ldots,\sigma_K^2).
\]
If the data errors are correlated, then $C_D$ has off-diagonal entries.

The inverse covariance
\[
    C_D^{-1}
\]
is called the precision matrix. It weights the data residual
\[
    r(m)=\obs-\G m.
\]
The usual Gaussian likelihood contains the quadratic form
\[
    r(m)^T C_D^{-1} r(m).
\]

This quantity penalizes residuals according to how surprising they are under the
noise model. A residual in a noisy datum is less serious than the same residual
in a highly accurate datum.

\begin{intuition}
The covariance says how the noise wiggles.
The precision says how strongly each residual is judged.
\end{intuition}

There are two related geometries here:
\begin{enumerate}[label=\arabic*.]
    \item the chosen data-space inner product $\langle\cdot,\cdot\rangle_{\D}$;
    \item the statistical weighting induced by the noise covariance $C_D$.
\end{enumerate}

In many practical presentations, one uses the Euclidean data inner product and
places all statistical weighting inside $C_D^{-1}$. Then
\[
    \|r\|_{C_D^{-1}}^2
    =
    r^T C_D^{-1} r.
\]

More abstractly, one may regard the inverse covariance as an operator whose
meaning depends on the chosen data-space geometry. The key lesson is that the
weights must be stated explicitly. They do not materialize from the matrix fog.

\subsection{Whitened and unwhitened data}

Sometimes it is convenient to transform the data so that the noise becomes
standard Gaussian.

Assume
\[
    C_D = LL^T,
\]
where $L$ is invertible. Define the whitened residual
\[
    \widetilde{r}(m)
    =
    L^{-1}(\obs-\G m).
\]
Then
\[
    \|\widetilde{r}(m)\|_2^2
    =
    (L^{-1}r)^T(L^{-1}r)
    =
    r^T L^{-T}L^{-1}r
    =
    r^T C_D^{-1} r.
\]

Thus, the weighted misfit in the original coordinates is the Euclidean misfit
in whitened coordinates.

Equivalently, define
\[
    \widetilde{\obs}=L^{-1}\obs,
    \qquad
    \widetilde{\G}=L^{-1}\G,
    \qquad
    \widetilde{\noise}=L^{-1}\noise.
\]
Then
\[
    \widetilde{\obs}
    =
    \widetilde{\G}m
    +
    \widetilde{\noise},
\]
and
\[
    \widetilde{\noise}
    \sim
    \mathcal{N}(0,I_K).
\]

Whitening is often computationally useful, but it is not magic. It is a change
of coordinates in data space. The underlying statistical model is the same.

\begin{remark}
Working with whitened data can make formulas look cleaner. Working with
unwhitened data can make the physical meaning of the errors more visible. Both
are acceptable when the transformation is handled consistently.
\end{remark}

\section{The data misfit}

The data misfit measures how well a model predicts the observed data. Let
\[
    r(m)
    =
    \obs-\G m.
\]
For Gaussian observational noise, the natural quadratic misfit is
\[
    \|r(m)\|_{C_D^{-1}}^2.
\]

With the Euclidean inner product on $\R^K$, this is
\[
    \|\obs - \G m\|_{C_D^{-1}}^2
    =
    (\obs-\G m)^T C_D^{-1}(\obs-\G m).
\]

Equivalently, writing the inner product explicitly,
\[
    \|\obs - \G m\|_{C_D^{-1}}^2
    =
    \langle \obs - \G m,
    C_D^{-1}(\obs - \G m)
    \rangle_{\D},
\]
provided $\langle\cdot,\cdot\rangle_{\D}$ is the Euclidean inner product.

If instead
\[
    \langle a,b\rangle_{\D}=a^T W_D b,
\]
then one must be careful about what is meant by $C_D^{-1}$ as an operator in
that geometry. In coordinates, the safest expression is the quadratic form
\[
    (\obs-\G m)^T C_D^{-1}(\obs-\G m),
\]
where $C_D$ is the covariance matrix of the coordinate noise.

The likelihood is then
\[
    \pi(\obs\mid m)
    \propto
    \exp
    \left(
        -
        \frac{1}{2}
        (\obs-\G m)^T C_D^{-1}(\obs-\G m)
    \right).
\]

Thus, the Bayesian update is driven by the potential
\[
    \Phi(m;\obs)
    =
    \frac{1}{2}
    \|\obs-\G m\|_{C_D^{-1}}^2.
\]

\begin{lesson}
The data misfit is not merely a least-squares penalty. It is the negative log
likelihood produced by the noise model.
\end{lesson}

% ------------------------------------------------------------

\chapter{The Forward Operator and Its Adjoint}

\section{The forward map}

The forward map sends a model to predicted data:
\[
    \G : \M \to \D.
\]

For the toy problem,
\[
    \D=\R^K,
\]
and
\[
    [\G m]_i
    =
    \int_\Omega K_i(x)m(x)\,dx.
\]

In words, each datum is obtained by testing the model against a sensitivity
kernel. The map $\G$ collects all these tests into a single data vector:
\[
    \G m
    =
    \begin{pmatrix}
        \G_1 m\\
        \vdots\\
        \G_K m
    \end{pmatrix}.
\]

If $\M$ is a Hilbert space and $\D$ is equipped with an inner product, then
$\G$ has an adjoint, provided $\G$ is bounded. The adjoint is one of the central
objects in linear inverse theory. It tells us how data-space residuals are
pulled back into model space.

In finite-dimensional Euclidean settings, this pullback is often written using
a transpose. But the transpose is only the simplest coordinate shadow of the
adjoint. The real object lives in the spaces.

\section{The adjoint is not just a transpose}

\subsection{Definition of the adjoint}

Let $\M$ and $\D$ be Hilbert spaces with inner products
\[
    \langle\cdot,\cdot\rangle_{\M}
    \qquad
    \text{and}
    \qquad
    \langle\cdot,\cdot\rangle_{\D}.
\]
The adjoint of a bounded linear operator
\[
    \G:\M\to\D
\]
is the unique bounded linear operator
\[
    \Gadj:\D\to\M
\]
satisfying
\[
    \langle \G m, d\rangle_{\D}
    =
    \langle m, \Gadj d\rangle_{\M}
\]
for all
\[
    m\in\M,
    \qquad
    d\in\D.
\]

This equation is the definition. It is not a computational trick. It says that
the same pairing can be evaluated either by pushing the model forward into data
space or by pulling the data vector back into model space.

\begin{intuition}
The forward map asks: ``What data does this model predict?''
The adjoint asks: ``What model-space direction is suggested by this data-space
quantity?''
\end{intuition}

\subsection{Why inner products determine the adjoint}

The adjoint depends on the inner products. Change the inner product, and the
adjoint changes.

Suppose the data-space inner product is
\[
    \langle a,b\rangle_{\D}=a^T W_D b.
\]
Suppose also that, after choosing a basis in a finite-dimensional model space,
the model-space inner product is represented by a symmetric positive definite
matrix $M$:
\[
    \langle u,v\rangle_{\M}=u^T M v.
\]

Let the forward map be represented by a matrix
\[
    G\in\R^{K\times N}.
\]
The adjoint matrix $A$ must satisfy
\[
    \langle Gu,d\rangle_{\D}
    =
    \langle u,Ad\rangle_{\M}
\]
for all
\[
    u\in\R^N,
    \qquad
    d\in\R^K.
\]

The left-hand side is
\[
    \langle Gu,d\rangle_{\D}
    =
    (Gu)^T W_D d
    =
    u^T G^T W_D d.
\]
The right-hand side is
\[
    \langle u,Ad\rangle_{\M}
    =
    u^T M A d.
\]
Since this must hold for every $u$ and $d$, we need
\[
    M A = G^T W_D.
\]
Therefore,
\[
    A = M^{-1}G^T W_D.
\]

Thus, the adjoint is represented by
\[
    G^\ast_{\mathrm{disc}}
    =
    M^{-1}G^T W_D.
\]

Only when
\[
    M=I_N
    \qquad
    \text{and}
    \qquad
    W_D=I_K
\]
does this reduce to
\[
    G^\ast_{\mathrm{disc}}=G^T.
\]

\subsection{Matrix transpose versus Hilbert adjoint}

The transpose is a coordinate operation. Given a matrix $G$, the transpose
$G^T$ is obtained by flipping rows and columns.

The Hilbert adjoint is a geometric operation. It is defined by the identity
\[
    \langle \G m,d\rangle_{\D}
    =
    \langle m,\Gadj d\rangle_{\M}.
\]

In Euclidean coordinates with orthonormal bases, these two operations coincide.
That coincidence is useful, but it is also dangerous. It can lead us to forget
that the transpose is not the fundamental object.

If the basis functions in the model space are not orthonormal, then the
coefficient vectors do not inherit the Euclidean inner product. For example,
with basis functions $\{\phi_j\}_{j=1}^N$, a coefficient vector
\[
    \bm{u}
\]
represents the function
\[
    m_N=\sum_{j=1}^N u_j\phi_j.
\]
The inner product between two such functions is
\[
    \langle m_N,n_N\rangle_{\M}
    =
    \sum_{i=1}^N\sum_{j=1}^N
    u_i v_j
    \langle \phi_i,\phi_j\rangle_{\M}.
\]
This is
\[
    \bm{u}^T M\bm{v},
\]
not generally
\[
    \bm{u}^T\bm{v}.
\]

\begin{warning}
Using $G^T$ as the adjoint when the basis is not orthonormal means silently
using the wrong geometry.
\end{warning}

\section{The role of Gram matrices}

\subsection{Non-orthonormal bases}

A basis
\[
    \{\phi_1,\ldots,\phi_N\}
\]
is orthonormal in $\M$ if
\[
    \langle \phi_i,\phi_j\rangle_{\M}
    =
    \delta_{ij}.
\]
Most finite element bases are not orthonormal. Hat functions, for example,
overlap with neighbouring hat functions, so their inner products are not zero
for neighbouring indices.

Therefore, the geometry of the function space is not represented by the
Euclidean dot product of coefficient vectors.

If
\[
    m_N=\sum_{j=1}^N u_j\phi_j
    \qquad
    \text{and}
    \qquad
    n_N=\sum_{j=1}^N v_j\phi_j,
\]
then
\[
    \langle m_N,n_N\rangle_{\M}
    =
    \bm{u}^T M\bm{v},
\]
where $M$ is the Gram matrix of the basis.

\subsection{The model-space Gram matrix}

The model-space Gram matrix is defined by
\[
    M_{ij}
    =
    \langle \phi_i,\phi_j\rangle_{\M}.
\]

For example, if
\[
    \M=L^2(\Omega),
\]
then
\[
    M_{ij}
    =
    \int_\Omega \phi_i(x)\phi_j(x)\,dx.
\]
This is the familiar finite element mass matrix.

If instead
\[
    \M=H^1(\Omega),
\]
then a common inner product is
\[
    \langle f,g\rangle_{H^1}
    =
    \int_\Omega f(x)g(x)\,dx
    +
    \int_\Omega f'(x)g'(x)\,dx.
\]
The Gram matrix becomes
\[
    M_{ij}
    =
    \int_\Omega \phi_i(x)\phi_j(x)\,dx
    +
    \int_\Omega \phi_i'(x)\phi_j'(x)\,dx.
\]

Thus, the same basis functions produce different Gram matrices depending on
the chosen model-space geometry.

\begin{intuition}
The Gram matrix is the memory of the function-space geometry after we choose
coordinates.
\end{intuition}

\subsection{The data-space Gram matrix}

Similarly, the data-space inner product may be written as
\[
    \langle a,b\rangle_{\D}
    =
    a^T W_D b.
\]
The matrix
\[
    W_D
\]
is the Gram matrix of the data-space coordinate basis under the chosen inner
product.

If the data coordinates are orthonormal, then
\[
    W_D=I_K.
\]
If the data coordinates carry weights, units, quadrature factors, or other
geometric structure, then $W_D$ may not be the identity.

In many finite-dimensional data problems, we choose Euclidean data coordinates
and therefore take
\[
    W_D=I_K.
\]
But this is still a choice, not an inevitability.

\subsection{The discrete adjoint}

Let
\[
    G\in\R^{K\times N}
\]
represent the forward map from model coefficients to data coordinates:
\[
    \bm{d}=G\bm{u}.
\]

Let the model-space inner product be represented by
\[
    M
\]
and the data-space inner product by
\[
    W_D.
\]

The discrete adjoint is the matrix
\[
    G^\ast_{\mathrm{disc}}
    =
    M^{-1}G^T W_D.
\]

Indeed, for all coefficient vectors $\bm{u}\in\R^N$ and data vectors
$\bm{d}\in\R^K$, we require
\[
    \langle G\bm{u},\bm{d}\rangle_{\D}
    =
    \langle \bm{u},G^\ast_{\mathrm{disc}}\bm{d}\rangle_{\M}.
\]
The left-hand side is
\[
    (G\bm{u})^T W_D\bm{d}
    =
    \bm{u}^T G^T W_D\bm{d}.
\]
The right-hand side is
\[
    \bm{u}^T M G^\ast_{\mathrm{disc}}\bm{d}.
\]
Therefore,
\[
    M G^\ast_{\mathrm{disc}}
    =
    G^T W_D,
\]
and hence
\[
    G^\ast_{\mathrm{disc}}
    =
    M^{-1}G^T W_D.
\]

This formula is one of the simplest places where the naive discretization
betrays itself. If we use $G^T$ instead of $M^{-1}G^T W_D$, then we are not
using the adjoint corresponding to the chosen model and data spaces.

\begin{proposition}
Let $G:\R^N\to\R^K$ be represented by a matrix $G$. Equip $\R^N$ with the inner
product
\[
    \langle u,v\rangle_{\M}=u^T Mv,
\]
where $M$ is symmetric positive definite, and equip $\R^K$ with
\[
    \langle a,b\rangle_{\D}=a^T W_D b,
\]
where $W_D$ is symmetric positive definite. Then the adjoint of $G$ with
respect to these inner products is represented by
\[
    G^\ast_{\mathrm{disc}}
    =
    M^{-1}G^T W_D.
\]
\end{proposition}

\begin{proof}
The adjoint $G^\ast_{\mathrm{disc}}$ is defined by
\[
    \langle Gu,d\rangle_{\D}
    =
    \langle u,G^\ast_{\mathrm{disc}}d\rangle_{\M}
\]
for all $u\in\R^N$ and $d\in\R^K$. Expanding both inner products gives
\[
    u^T G^T W_D d
    =
    u^T M G^\ast_{\mathrm{disc}} d.
\]
Since this holds for all $u$ and $d$, we have
\[
    G^T W_D
    =
    M G^\ast_{\mathrm{disc}}.
\]
Multiplying by $M^{-1}$ gives
\[
    G^\ast_{\mathrm{disc}}
    =
    M^{-1}G^T W_D.
\]
\end{proof}

\section{The geometry-aware discretization}

\subsection{What changes}

In the naive discretization, we often write posterior formulas using $G^T$:
\[
    C_N^y
    =
    C_N
    -
    C_N G^T
    \left(
        G C_N G^T + C_D
    \right)^{-1}
    G C_N.
\]

This formula is correct only under particular coordinate assumptions. It assumes
that the relevant adjoint is represented by $G^T$. That is true when the model
and data coordinates are orthonormal for the inner products being used.

With non-orthonormal basis functions, the geometry-aware adjoint is instead
\[
    G^\ast_{\mathrm{disc}}
    =
    M^{-1}G^T W_D.
\]

Therefore, the posterior formulas must be interpreted as discretizations of the
function-space formulas:
\[
    C^y
    =
    C_0
    -
    C_0\Gadj
    \left(
        \G C_0\Gadj + C_D
    \right)^{-1}
    \G C_0,
\]
and
\[
    m^y
    =
    m_0
    +
    C_0\Gadj
    \left(
        \G C_0\Gadj + C_D
    \right)^{-1}
    (\obs-\G m_0).
\]

The object to discretize is $\Gadj$, not the symbol $G^T$.

\subsection{What stays invariant}

A geometry-aware discretization keeps the continuous problem fixed. The model
space, data space, forward map, prior, and noise model are chosen first:
\[
    \M,
    \qquad
    \D,
    \qquad
    \G,
    \qquad
    \prior,
    \qquad
    C_D.
\]

Then the numerical representation is built to approximate those objects.

Changing the basis should not change the mathematical problem. If we represent
the same finite-dimensional subspace using a different basis, the coordinate
matrices may change, but the underlying functions, operators, and probability
measures should represent the same objects.

This is the point of including Gram matrices. They ensure that coordinate
vectors remember the geometry of the functions they represent.

\begin{intuition}
The coordinates may change costume. The function-space problem should remain
the same actor.
\end{intuition}

\subsection{Why this prevents changing the problem by accident}

When we use the wrong adjoint, we are not merely making a small algebraic
mistake. We are changing the inverse problem.

The adjoint appears in:
\begin{enumerate}[label=\arabic*.]
    \item posterior mean formulas;
    \item posterior covariance formulas;
    \item gradient calculations;
    \item normal equations;
    \item data-to-model backprojection;
    \item covariance updates.
\end{enumerate}

If the adjoint is wrong, then residuals are pulled back into the wrong
model-space directions. The posterior mean can shift. The uncertainty can be
misrepresented. The discretized problem can stop approximating the intended
continuous problem.

This is why the model-space Gram matrix $M$ and data-space Gram matrix $W_D$
are not numerical clutter. They are the finite-dimensional fingerprints of the
chosen spaces.

\begin{lesson}
The transpose belongs to Euclidean coordinates.
The adjoint belongs to the spaces.
\end{lesson}

\section{The moral of Act V}

The data space was easy, but not empty of choices. Its finite-dimensional nature
gave us automatic completeness and friendly convergence, but its inner product
still mattered. The noise covariance told us how to weight residuals. Whitening
showed that data weighting can be moved into coordinates, provided we do so
consistently.

The forward map then forced us to confront the adjoint. The adjoint is defined
by
\[
    \langle \G m,d\rangle_{\D}
    =
    \langle m,\Gadj d\rangle_{\M},
\]
which means it depends on the inner products of both spaces. Once basis
functions enter, those inner products become Gram matrices. For a
non-orthonormal model basis and a weighted data geometry, the discrete adjoint
is
\[
    G^\ast_{\mathrm{disc}}
    =
    M^{-1}G^T W_D,
\]
not simply $G^T$.

\begin{checkpoint}
We have now chosen the model space, the data space, the forward map, and the
geometry needed to define the adjoint. The next act turns to the remaining
ingredient: the prior measure. This is where the phrase ``the prior is not just
a matrix'' finally gets its full teeth.
\end{checkpoint}
% ============================================================
% ACT VI
% ============================================================

\part{Act VI: The Prior Is Not Just a Matrix}

\chapter{A Detour Through Measure Theory}

\section{Why probabilities are measures}

We now reach the last major ingredient of the function-space Bayesian inverse
problem: the prior.

In the naive discretized problem, the prior was introduced as a covariance
matrix,
\[
    C_N = \sigma^2 I_N.
\]
For every fixed $N$, this is a valid covariance matrix on $\R^N$. But the
failure of Act II showed us that this is not enough. A prior for a function
space inverse problem must be a probability measure on a function space.

This forces us to step briefly into measure theory.

The reason is simple: probability is not fundamentally about densities.
Probability is fundamentally about assigning mass to events.

\subsection{Events}

Let $X$ be a space of possible outcomes. In our setting, $X$ might be a model
space,
\[
    X = \M.
\]
An event is a subset of $X$. For example,
\[
    A \subset \M
\]
could be the event
\[
    A =
    \{m\in\M : \|m\|_{\M}\leq 1\}.
\]
If $m$ is random, then we may ask for the probability
\[
    \mathbb{P}(m\in A).
\]

In finite dimensions, we often imagine events as regions in $\R^N$. For
example, if
\[
    \bm{u}\in\R^N,
\]
then an event might be a ball,
\[
    A =
    \{\bm{u}\in\R^N : \|\bm{u}\|_2\leq r\}.
\]

In function space, events can be more subtle. We might ask whether a random
function has small norm, whether it is positive on a subinterval, whether its
derivative is bounded in some weak sense, or whether it lies near a reference
model.

The prior must assign probabilities to such events.

\subsection{Sigma-algebras}

Not every subset of an infinite-dimensional space can be assigned a probability
in a consistent way. Instead, we choose a collection of admissible events called
a sigma-algebra.

A sigma-algebra $\mathcal{A}$ on a set $X$ is a collection of subsets of $X$
such that:
\begin{enumerate}[label=\arabic*.]
    \item $X\in\mathcal{A}$;
    \item if $A\in\mathcal{A}$, then $X\setminus A\in\mathcal{A}$;
    \item if $A_1,A_2,A_3,\ldots\in\mathcal{A}$, then
    \[
        \bigcup_{n=1}^{\infty} A_n \in \mathcal{A}.
    \]
\end{enumerate}

The elements of $\mathcal{A}$ are the events whose probabilities we are allowed
to discuss.

For a topological space, a common choice is the Borel sigma-algebra, generated
by the open sets. If $\M$ is a Hilbert space, we usually take its Borel
sigma-algebra and define probability measures on it.

\begin{intuition}
A sigma-algebra is the list of questions the probability measure agrees to
answer.
\end{intuition}

\subsection{Probability measures}

A probability measure on $(X,\mathcal{A})$ is a map
\[
    \mu:\mathcal{A}\to[0,1]
\]
such that
\[
    \mu(X)=1
\]
and, for any countable collection of disjoint events
\[
    A_1,A_2,A_3,\ldots,
\]
we have
\[
    \mu\left(\bigcup_{n=1}^{\infty}A_n\right)
    =
    \sum_{n=1}^{\infty}\mu(A_n).
\]

In Bayesian inverse problems, the prior is a probability measure
\[
    \prior
\]
on the model space $\M$. The posterior is another probability measure
\[
    \post
\]
on the same space, obtained by updating the prior using the data.

Thus, at the most fundamental level, we should write
\[
    m\sim\prior,
\]
not
\[
    \bm{u}\sim\mathcal{N}(0,\sigma^2 I_N),
\]
unless $\bm{u}$ is truly the unknown object and the problem is genuinely
finite-dimensional.

\begin{lesson}
A prior is a probability measure on possible models.
A covariance matrix is only one finite-dimensional way of describing part of
such a measure.
\end{lesson}

\section{Why densities are not fundamental}

In finite-dimensional probability, we often describe a distribution by a
density. For example, the Gaussian density on $\R^N$ is
\[
    \rho(\bm{u})
    =
    \frac{1}{(2\pi)^{N/2}\det(C)^{1/2}}
    \exp
    \left(
        -
        \frac{1}{2}
        (\bm{u}-\bar{\bm{u}})^T C^{-1}
        (\bm{u}-\bar{\bm{u}})
    \right).
\]

This formula is familiar, but it hides an important fact. A density is not a
probability measure by itself. A density is always a density with respect to
some reference measure.

\subsection{Densities relative to a reference measure}

Let $\mu$ and $\nu$ be measures on the same measurable space. If $\mu$ is
absolutely continuous with respect to $\nu$, then there exists a density
$\rho$ such that
\[
    d\mu = \rho\,d\nu.
\]
This means that, for every event $A$,
\[
    \mu(A)
    =
    \int_A \rho\,d\nu.
\]

The density $\rho$ is also written as
\[
    \rho
    =
    \frac{d\mu}{d\nu},
\]
the Radon-Nikodym derivative of $\mu$ with respect to $\nu$.

Thus, the expression
\[
    d\mu = \rho\,d\nu
\]
does not define $\rho$ in isolation. It defines $\mu$ by weighting the reference
measure $\nu$.

\begin{warning}
A density is always relative to a background measure. Without the background
measure, the density is only half a sentence.
\end{warning}

\subsection{The role of Lebesgue measure in finite dimensions}

In $\R^N$, the usual background measure is Lebesgue measure. When we write
\[
    \bm{u}\sim\mathcal{N}(\bar{\bm{u}},C),
\]
we often mean that the probability measure has density
\[
    \rho(\bm{u})
    =
    \frac{1}{(2\pi)^{N/2}\det(C)^{1/2}}
    \exp
    \left(
        -
        \frac{1}{2}
        \|\bm{u}-\bar{\bm{u}}\|_{C^{-1}}^2
    \right)
\]
with respect to Lebesgue measure
\[
    d\bm{u}.
\]

This is why finite-dimensional Bayesian formulas are often written as if
probability measures and densities were the same thing:
\[
    \pi(\bm{u}\mid \bm{y})
    \propto
    \pi(\bm{y}\mid \bm{u})\pi(\bm{u}).
\]

The hidden reference measure is Lebesgue measure. It is so familiar that it
becomes invisible.

\subsection{Why there is no infinite-dimensional Lebesgue measure}

In infinite-dimensional function spaces, there is no analogue of Lebesgue
measure with the same useful properties. There is no translation-invariant,
locally finite measure that plays the role of volume on an infinite-dimensional
Hilbert space.

This changes everything.

If there is no infinite-dimensional Lebesgue measure, then we cannot define a
function-space prior by writing a density with respect to such a measure:
\[
    \rho(m)\,dm.
\]
There is no meaningful infinite-dimensional volume element $dm$ of this kind.

Therefore, in function space, it is safer and more fundamental to define the
prior directly as a measure:
\[
    \prior.
\]

The posterior is then often defined by specifying its density relative to the
prior:
\[
    \frac{d\post}{d\prior}(m)
    =
    \frac{1}{Z(\obs)}
    \exp(-\Phi(m;\obs)).
\]

Here the prior is the reference measure. The likelihood reweights it.

\begin{lesson}
In finite dimensions, Lebesgue measure quietly carries the density formalism.
In function space, that quiet support disappears.
\end{lesson}

\section{What changes on function spaces}

\subsection{A prior must be a measure}

In a function-space inverse problem, the prior must be a probability measure on
the model space:
\[
    \prior \in \mathcal{P}(\M).
\]

This means that the prior assigns probabilities to events such as
\[
    \{m\in\M : \|m\|_{\M}\leq r\},
\]
or
\[
    \{m\in\M : \|m-m_0\|_{\M}\leq r\}.
\]

The prior must be defined before discretization. A discretized prior should then
be an approximation of this measure, not a new probability model invented at
each resolution.

This is exactly where the naive prior failed. The sequence
\[
    \mathcal{N}(0,\sigma^2 I_N)
\]
is a perfectly legal sequence of finite-dimensional Gaussian measures, but it
does not define a sensible Gaussian probability measure on a Hilbert space
whose covariance is $\sigma^2 I$.

\subsection{Gaussian measures on infinite-dimensional spaces}

A Gaussian measure on a Hilbert space $\M$ is a probability measure such that
every continuous linear functional applied to the random model is a
finite-dimensional Gaussian random variable.

That is, if
\[
    m\sim\mu
\]
and
\[
    \ell:\M\to\R
\]
is continuous and linear, then
\[
    \ell(m)
\]
is a Gaussian random variable.

A Gaussian measure is characterized by its mean $m_0$ and covariance operator
$C_0$:
\[
    \mu = \mathcal{N}(m_0,C_0).
\]

The mean satisfies
\[
    \ell(m_0)
    =
    \mathbb{E}[\ell(m)]
\]
for every continuous linear functional $\ell$. The covariance is characterized
by
\[
    \operatorname{Cov}(\ell_1(m),\ell_2(m))
    =
    \langle C_0 h_1,h_2\rangle_{\M},
\]
where $h_1,h_2\in\M$ represent the linear functionals through the Hilbert-space
inner product:
\[
    \ell_j(m)=\langle m,h_j\rangle_{\M}.
\]

In finite dimensions, almost any symmetric positive semidefinite matrix gives a
Gaussian covariance. In infinite dimensions, the covariance operator must
satisfy stronger conditions. Most importantly for our story, it must be trace
class if the Gaussian is to live on the Hilbert space itself.

\subsection{Typical samples versus pointwise intuition}

In function spaces, typical samples may be less regular than the covariance
kernel or mean function suggests at first glance. This is already visible in
familiar random processes.

For example, a Brownian path is continuous but almost surely nowhere
classically differentiable. A Gaussian field may have a smooth covariance
function while its samples are still rougher than one might expect.

This matters because a prior is not only about its mean or covariance as formal
objects. It is about the random functions it produces:
\[
    m\sim\prior.
\]

When choosing a prior, we should ask:
\begin{enumerate}[label=\arabic*.]
    \item Do typical samples belong to the intended model space?
    \item Do they have physically plausible roughness?
    \item Do they satisfy the intended boundary conditions?
    \item Does the prior uncertainty remain controlled under discretization?
\end{enumerate}

\begin{warning}
In function space, writing down a density without specifying the underlying
measure is usually smoke without a fireplace.
\end{warning}

% ------------------------------------------------------------

\chapter{Gaussian Priors on Function Spaces}

\section{Mean and covariance}

A Gaussian prior on the model space is written
\[
    m \sim \mathcal{N}(m_0,C_0).
\]

Here
\[
    m_0\in\M
\]
is the prior mean, and
\[
    C_0
\]
is the prior covariance operator.

The mean describes the central model before observing the data. The covariance
describes the allowed uncertainty around that mean.

A useful way to think about a Gaussian prior is through a random perturbation:
\[
    m = m_0+\xi,
\]
where
\[
    \xi \sim \mathcal{N}(0,C_0).
\]

Thus, the prior says:
\[
    \text{models are random perturbations of } m_0,
\]
where the geometry and regularity of those perturbations are controlled by
$C_0$.

\begin{intuition}
The prior mean says where the cloud is centred.
The covariance operator says what shape the cloud has in function space.
\end{intuition}

\section{Covariance as an operator}

In finite dimensions, covariance is a matrix. In a Hilbert space, covariance is
an operator:
\[
    C_0 : \M \to \M.
\]

For a centred Gaussian random model
\[
    \xi \sim \mathcal{N}(0,C_0),
\]
the covariance operator satisfies
\[
    \mathbb{E}
    \left[
        \langle \xi,f\rangle_{\M}
        \langle \xi,g\rangle_{\M}
    \right]
    =
    \langle C_0 f,g\rangle_{\M}
\]
for all
\[
    f,g\in\M.
\]

This formula tells us how the random model fluctuates along the directions
$f$ and $g$. In particular, the variance in direction $f$ is
\[
    \operatorname{Var}(\langle \xi,f\rangle_{\M})
    =
    \langle C_0 f,f\rangle_{\M}.
\]

Thus, $C_0$ encodes the geometry of uncertainty.

If $C_0$ has eigenpairs
\[
    C_0 e_j = \lambda_j e_j,
\]
where
\[
    \{e_j\}_{j=1}^{\infty}
\]
is an orthonormal basis of $\M$, then the prior can often be represented by the
Karhunen-Loeve expansion
\[
    m
    =
    m_0
    +
    \sum_{j=1}^{\infty}
    \sqrt{\lambda_j}\,\xi_j e_j,
\]
where
\[
    \xi_j\sim\mathcal{N}(0,1)
\]
are independent standard normal random variables.

This expansion makes the role of the eigenvalues clear. The number
$\lambda_j$ is the variance assigned to direction $e_j$.

\section{Required properties}

Not every linear operator is a valid covariance operator. To define a Gaussian
measure on a Hilbert space, the covariance must satisfy several important
properties.

\subsection{Linearity}

The covariance operator must be linear:
\[
    C_0(\alpha f+\beta g)
    =
    \alpha C_0f+\beta C_0g
\]
for all
\[
    f,g\in\M,
    \qquad
    \alpha,\beta\in\R.
\]

Linearity reflects the fact that covariance is bilinear in the directions being
tested.

\subsection{Symmetry}

The covariance must be symmetric:
\[
    \langle C_0 f,g\rangle_{\M}
    =
    \langle f,C_0 g\rangle_{\M}.
\]

This is the operator version of the symmetry of a covariance matrix:
\[
    C_{ij}=C_{ji}.
\]

Symmetry is natural because
\[
    \operatorname{Cov}(X,Y)
    =
    \operatorname{Cov}(Y,X).
\]

\subsection{Positivity}

The covariance must be positive:
\[
    \langle C_0 f,f\rangle_{\M} \geq 0.
\]

This is required because
\[
    \langle C_0 f,f\rangle_{\M}
    =
    \operatorname{Var}(\langle \xi,f\rangle_{\M}),
\]
and variances cannot be negative.

If
\[
    \langle C_0 f,f\rangle_{\M}=0,
\]
then the prior has no uncertainty in the direction $f$.

\subsection{Trace class}

In infinite-dimensional Hilbert spaces, positivity and symmetry are not enough.
The covariance must also be trace class:
\[
    \Tr(C_0) < \infty.
\]

If
\[
    C_0 e_j = \lambda_j e_j
\]
for an orthonormal basis $\{e_j\}$, then
\[
    \Tr(C_0)
    =
    \sum_{j=1}^{\infty}\lambda_j.
\]

Thus, trace class means
\[
    \sum_{j=1}^{\infty}\lambda_j < \infty.
\]

This condition says that the total variance assigned by the prior is finite.
The eigenvalues must decay quickly enough that the Gaussian perturbation has
finite expected squared norm in the model space.

\begin{warning}
A positive symmetric operator may still fail to define a Gaussian measure on
the Hilbert space if its trace is infinite.
\end{warning}

\section{Why trace class matters}

\subsection{Finite expected energy}

Let
\[
    m\sim\mathcal{N}(m_0,C_0)
\]
on a Hilbert space $\M$. If the covariance has eigenpairs
\[
    C_0 e_j = \lambda_j e_j,
\]
then the random perturbation can be written formally as
\[
    m-m_0
    =
    \sum_{j=1}^{\infty}
    \sqrt{\lambda_j}\,\xi_j e_j.
\]

The expected squared norm is
\[
    \mathbb{E}\|m-m_0\|_{\M}^2
    =
    \mathbb{E}
    \left\|
        \sum_{j=1}^{\infty}
        \sqrt{\lambda_j}\,\xi_j e_j
    \right\|_{\M}^2.
\]

Using orthonormality,
\[
    \left\|
        \sum_{j=1}^{\infty}
        \sqrt{\lambda_j}\,\xi_j e_j
    \right\|_{\M}^2
    =
    \sum_{j=1}^{\infty}
    \lambda_j\xi_j^2.
\]

Taking expectations gives
\[
    \mathbb{E}\|m-m_0\|_{\M}^2
    =
    \sum_{j=1}^{\infty}
    \lambda_j
    \mathbb{E}[\xi_j^2].
\]
Since
\[
    \mathbb{E}[\xi_j^2]=1,
\]
we obtain
\[
    \mathbb{E}\|m-m_0\|_{\M}^2
    =
    \sum_{j=1}^{\infty}\lambda_j
    =
    \Tr(C_0).
\]

Therefore,
\[
    \E\|m-m_0\|_{\M}^2
    =
    \Tr(C_0).
\]

This is the infinite-dimensional version of the finite-dimensional identity
\[
    \mathbb{E}\|\bm{u}-\bar{\bm{u}}\|_2^2
    =
    \Tr(C).
\]

\subsection{The failure of the identity covariance}

Now consider the tempting covariance
\[
    C_0 = \sigma^2 I.
\]
If $\M$ is infinite-dimensional and
\[
    \{e_j\}_{j=1}^{\infty}
\]
is an orthonormal basis, then
\[
    C_0 e_j = \sigma^2 e_j.
\]
Thus, every eigenvalue is
\[
    \lambda_j=\sigma^2.
\]

The trace is therefore
\[
    \Tr(\sigma^2 I)
    =
    \sum_{j=1}^{\infty}\sigma^2
    =
    \infty.
\]

So
\[
    \Tr(\sigma^2 I) = \infty.
\]

This means that the formal Gaussian perturbation
\[
    \xi
    =
    \sum_{j=1}^{\infty}
    \sigma \xi_j e_j
\]
does not have finite expected squared norm in $\M$:
\[
    \mathbb{E}\|\xi\|_{\M}^2
    =
    \infty.
\]

This is exactly the continuous analogue of the failure we saw in the naive
discretization:
\[
    \Tr(\sigma^2 I_N)=\sigma^2 N\to\infty.
\]

In finite dimensions, $\sigma^2 I_N$ is harmless for fixed $N$. In infinite
dimensions, $\sigma^2 I$ is not a valid covariance operator for a Gaussian
measure living on the Hilbert space.

\begin{lesson}
A covariance operator must produce actual random functions in the chosen model
space. That requirement is far stronger than being a positive matrix after
discretization.
\end{lesson}

\section{What the covariance really chooses}

The covariance operator chooses more than variance. It chooses the regularity,
correlation, length scale, and boundary behaviour of prior samples.

Suppose
\[
    C_0 e_j=\lambda_j e_j.
\]
Then the prior perturbation is built from modes $e_j$ with amplitudes controlled
by
\[
    \sqrt{\lambda_j}.
\]

If the eigenvalues decay rapidly, high-frequency modes are strongly suppressed,
and samples are smoother. If the eigenvalues decay slowly, high-frequency modes
retain more energy, and samples are rougher.

Thus, the covariance operator encodes statements such as:
\begin{enumerate}[label=\arabic*.]
    \item how large model perturbations are expected to be;
    \item how correlated nearby points are;
    \item how smooth or rough samples should be;
    \item what boundary behaviour samples should satisfy;
    \item which directions are plausible and which are unlikely.
\end{enumerate}

\begin{intuition}
A covariance matrix says how a finite vector wiggles.
A covariance operator says how a random function breathes.
\end{intuition}

% ------------------------------------------------------------

\chapter{How to Build Covariances}

\section{Covariances from differential operators}

One powerful way to build covariance operators is to start from differential
operators. This is natural when prior information is about smoothness, length
scale, and boundary behaviour.

\subsection{The Laplacian}

The Laplacian is the differential operator
\[
    \Delta.
\]
On an interval,
\[
    \Delta f = f''.
\]
The negative Laplacian,
\[
    -\Delta,
\]
is often positive under appropriate boundary conditions.

The eigenfunctions of $-\Delta$ provide natural spatial modes. For example, on
a bounded interval with suitable boundary conditions, one may have
\[
    -\Delta e_j = \rho_j e_j,
\]
where
\[
    0\leq \rho_1\leq \rho_2\leq\cdots,
    \qquad
    \rho_j\to\infty.
\]

High eigenvalues correspond to oscillatory modes. Therefore, functions of
$-\Delta$ can be used to control how much prior variance is assigned to
different spatial frequencies.

\subsection{Shifted elliptic operators}

A common construction uses the shifted elliptic operator
\[
    A = \kappa^2 I - \Delta,
\]
where
\[
    \kappa>0.
\]
The parameter $\kappa$ sets a length scale. Roughly speaking, larger $\kappa$
corresponds to shorter correlation lengths, while smaller $\kappa$ corresponds
to longer correlation lengths.

A covariance operator can then be defined as
\[
    C_0
    =
    \alpha
    \left(
        \kappa^2 I - \Delta
    \right)^{-s}.
\]

Here
\[
    \alpha>0
\]
sets the overall variance scale, and
\[
    s>0
\]
controls smoothness.

If
\[
    -\Delta e_j=\rho_j e_j,
\]
then
\[
    (\kappa^2 I-\Delta)e_j
    =
    (\kappa^2+\rho_j)e_j.
\]
Therefore,
\[
    C_0 e_j
    =
    \alpha
    (\kappa^2+\rho_j)^{-s}e_j.
\]

The covariance eigenvalues are
\[
    \lambda_j
    =
    \alpha(\kappa^2+\rho_j)^{-s}.
\]

As $j$ increases, $\rho_j$ grows, so $\lambda_j$ decays. High-frequency modes
receive less prior variance.

\subsection{Bessel-type priors}

Priors of the form
\[
    C_0
    =
    \alpha
    \left(
        \kappa^2 I - \Delta
    \right)^{-s}
\]
are often called Bessel-type priors or Matérn-type constructions, depending on
the domain, boundary conditions, and interpretation.

They are useful because they connect three things:
\begin{enumerate}[label=\arabic*.]
    \item the covariance operator;
    \item the smoothness of samples;
    \item the differential geometry of the domain through the Laplacian.
\end{enumerate}

The parameter $s$ controls how aggressively high-frequency modes are damped.
For the covariance to be trace class on the chosen Hilbert space, $s$ must be
large enough relative to the dimension of the domain and the geometry of the
operator.

In one dimension, the eigenvalues of $-\Delta$ typically grow like
\[
    \rho_j \sim j^2.
\]
Then
\[
    \lambda_j
    \sim
    j^{-2s}.
\]
The trace condition
\[
    \sum_{j=1}^{\infty}\lambda_j<\infty
\]
requires, in this simple heuristic,
\[
    2s>1.
\]
Thus,
\[
    s>\frac{1}{2}.
\]

This illustrates the general principle: the covariance eigenvalues must decay
fast enough that the total prior variance is finite.

\subsection{The role of boundary conditions}

The operator
\[
    -\Delta
\]
is not fully defined until its domain and boundary conditions are specified.

For example, on
\[
    \Omega=[0,L],
\]
one may impose Dirichlet conditions,
\[
    f(0)=f(L)=0,
\]
Neumann conditions,
\[
    f'(0)=f'(L)=0,
\]
or mixed conditions,
\[
    f(0)=0,
    \qquad
    f'(L)=0.
\]

These choices change the eigenfunctions and eigenvalues of the Laplacian.
Therefore, they change the covariance operator
\[
    C_0
    =
    \alpha(\kappa^2 I-\Delta)^{-s}.
\]

This means that boundary conditions are part of the prior.

If the prior is built from a differential operator, then specifying the
boundary conditions is as important as specifying $\alpha$, $\kappa$, and $s$.
The boundary behaviour of posterior samples is inherited from the prior unless
the data strongly constrain it.

\begin{warning}
A Laplacian without boundary conditions is not yet a covariance recipe. It is
only the beginning of one.
\end{warning}

\section{Covariances from kernels}

Another common way to build covariances is through covariance functions, also
called kernels.

\subsection{Covariance functions}

A covariance function is a function
\[
    k:\Omega\times\Omega\to\R
\]
defined by
\[
    k(x,x')
    =
    \Cov(m(x),m(x')).
\]

The value $k(x,x')$ describes how the random model values at $x$ and $x'$ vary
together.

If $k(x,x')$ is large and positive, then $m(x)$ and $m(x')$ tend to move in the
same direction. If it is close to zero, then they are weakly correlated. If it
is negative, then they tend to move in opposite directions.

The diagonal value
\[
    k(x,x)
\]
is the pointwise variance:
\[
    k(x,x)
    =
    \operatorname{Var}(m(x)).
\]

Common covariance kernels encode amplitude, length scale, smoothness, and
stationarity. For example, a squared-exponential-type kernel produces very
smooth samples, while Matérn-type kernels allow a tunable degree of roughness.

\subsection{From kernels to covariance operators}

A covariance kernel can define an integral operator:
\[
    (C_0 f)(x)
    =
    \int_\Omega k(x,x')f(x')\,dx'.
\]

This operator maps test functions to functions. Under suitable assumptions on
$k$, it is symmetric and positive.

Symmetry of the covariance kernel means
\[
    k(x,x')=k(x',x).
\]
Positivity means that, for any finite collection of points
\[
    x_1,\ldots,x_n\in\Omega
\]
and coefficients
\[
    a_1,\ldots,a_n\in\R,
\]
we have
\[
    \sum_{i=1}^n\sum_{j=1}^n
    a_i a_j k(x_i,x_j)
    \geq 0.
\]

When these conditions hold, the kernel can be used to build a covariance
operator.

\subsection{Positive-definite kernels}

A kernel $k$ is positive definite if, for every finite set of points
\[
    x_1,\ldots,x_n
\]
and every vector
\[
    \bm{a}\in\R^n,
\]
the matrix
\[
    K_{ij}=k(x_i,x_j)
\]
satisfies
\[
    \bm{a}^T K\bm{a}\geq 0.
\]

This condition ensures that any finite collection of model values
\[
    \bigl(m(x_1),\ldots,m(x_n)\bigr)
\]
can have a valid Gaussian covariance matrix.

However, a positive-definite kernel is not automatically the end of the story.
For a function-space prior, we must also ask what space the resulting random
function lives in, what regularity its samples have, and whether the associated
operator is trace class on the chosen model space.

\begin{warning}
A kernel that gives valid finite-dimensional covariance matrices still needs
to be checked against the chosen function space.
\end{warning}

\subsection{Smoothness encoded by kernels}

The smoothness of a covariance kernel is related to the smoothness of prior
samples, but the relationship must be handled carefully.

Very smooth kernels often produce smooth samples. Rougher kernels produce
rougher samples. Length-scale parameters determine how quickly correlations
decay with distance.

For example, a stationary kernel on $\R^d$ may have the form
\[
    k(x,x') = q(x-x'),
\]
so correlation depends only on separation. On bounded domains, boundary
effects and boundary conditions may break or modify this simple structure.

Kernels are often intuitive because they allow us to say:
\[
    \text{nearby points should be strongly correlated, distant points less so.}
\]
Operator constructions are often powerful because they connect directly to
differential equations, boundary conditions, and spectral decay.

Both perspectives are useful.

\section{Operator view versus kernel view}

\subsection{When each view is useful}

The kernel view is useful when prior information is naturally spatial:
\begin{enumerate}[label=\arabic*.]
    \item What is the variance at each point?
    \item How quickly should correlations decay?
    \item Should the prior be stationary?
    \item How smooth should the covariance function be?
\end{enumerate}

The operator view is useful when prior information is naturally geometric or
spectral:
\begin{enumerate}[label=\arabic*.]
    \item Which modes should be damped?
    \item How should boundary conditions enter?
    \item What differential operator controls smoothness?
    \item Does the covariance have finite trace?
\end{enumerate}

In numerical inverse problems, the operator view is often especially clean
because it tells us how to discretize the covariance in a way that respects the
continuous model. Rather than inventing a covariance matrix on coefficients, we
discretize an operator that already exists.

\subsection{How they describe the same object}

Under suitable conditions, the kernel and operator views describe the same
covariance.

Given a covariance kernel $k$, we may define
\[
    (C_0 f)(x)
    =
    \int_\Omega k(x,x')f(x')\,dx'.
\]

Conversely, given a covariance operator $C_0$, one may sometimes represent it
by a kernel:
\[
    C_0(x,x') = k(x,x').
\]

Spectrally, if
\[
    C_0 e_j=\lambda_j e_j,
\]
then the kernel may be written formally as
\[
    k(x,x')
    =
    \sum_{j=1}^{\infty}
    \lambda_j e_j(x)e_j(x').
\]

This representation shows the connection clearly:
\begin{enumerate}[label=\arabic*.]
    \item the eigenfunctions $e_j$ provide spatial modes;
    \item the eigenvalues $\lambda_j$ assign variance to those modes;
    \item the kernel records the resulting covariance between locations.
\end{enumerate}

The correct object to choose is not necessarily a matrix, a kernel, or a
differential expression in isolation. The correct object is a covariance
operator that defines a valid prior measure on the chosen model space.

\begin{lesson}
The prior covariance is a modelling object first and a computational object
second.
\end{lesson}

\section{The moral of Act VI}

The naive prior failed because it treated covariance as a matrix attached to a
mesh. In function space, this is the wrong starting point.

A prior must be a probability measure on the model space:
\[
    m\sim\prior.
\]
If the prior is Gaussian, then it is described by a mean and a covariance
operator:
\[
    \prior=\mathcal{N}(m_0,C_0).
\]
That covariance must be linear, symmetric, positive, and trace class. The trace
condition is the mathematical gatekeeper that prevents the prior from assigning
infinite total variance to the model.

The identity covariance failed because
\[
    \Tr(\sigma^2 I)=\infty
\]
in infinite dimensions. The finite-dimensional symptom was
\[
    \Tr(\sigma^2 I_N)=\sigma^2 N\to\infty.
\]

To build meaningful covariances, we may use differential operators, kernels, or
spectral constructions. But in every case, the covariance must define a valid
Gaussian measure on the chosen model space, and it must be discretized as an
approximation of that continuous object.

\begin{checkpoint}
We now have all the pieces: model space, data space, forward map, adjoint,
noise model, and prior measure. The final act puts them together and returns to
the opening toy problem, where the pathologies of the naive posterior disappear
once the formulation is done in the right order.
\end{checkpoint}
% ============================================================
% ACT VII
% ============================================================

\part{Act VII: Putting the Pieces Together}

\chapter{The Function-Space Bayesian Inverse Problem}

\section{The chosen ingredients}

We now return from the detour. The pieces are on the table: model space, data
space, forward map, noise model, prior measure, and the geometry needed to
define the adjoint.

The central lesson is that these objects must be chosen before discretization.
Only then does a matrix become a representation of something meaningful, rather
than an accidental substitute for the problem itself.

The function-space Bayesian inverse problem begins with the data equation
\[
    \obs = \G m + \noise.
\]
This equation now has a precise meaning because each symbol has a mathematical
home:
\[
    m\in\M,
    \qquad
    \obs\in\D,
    \qquad
    \G:\M\to\D,
    \qquad
    \noise\sim\mathcal{N}(0,C_D),
    \qquad
    m\sim\prior.
\]

\begin{lesson}
The same equation that looked innocent in Act I now has a full mathematical
skeleton. That skeleton is what prevents the discretization from inventing a
different creature.
\end{lesson}

\subsection{Model space}

The unknown model lives in a chosen Hilbert space:
\[
    \M = H^s(\Omega)
    \quad
    \text{or another chosen Hilbert space.}
\]

The precise choice depends on the problem. We might choose
\[
    \M=L^2(\Omega)
\]
if average-square size is the main quantity of interest. We might choose
\[
    \M=H^1(\Omega)
\]
if one weak derivative should be controlled. We might choose a higher-order or
fractional Sobolev space
\[
    \M=H^s(\Omega)
\]
if stronger or intermediate smoothness is required.

The model space tells us:
\begin{enumerate}[label=\arabic*.]
    \item which objects count as admissible models;
    \item what convergence means;
    \item what perturbations are small;
    \item what inner product defines the adjoint;
    \item what kind of Gaussian prior can live on the space.
\end{enumerate}

The important point is not that one particular Hilbert space is universally
correct. The important point is that a model space must be chosen deliberately.

\begin{warning}
A basis does not choose the model space for us. A basis represents a space after
the modelling decision has already been made.
\end{warning}

\subsection{Data space}

The data space is finite-dimensional:
\[
    \D = \R^K.
\]

The observations are collected into a vector
\[
    \obs =
    \begin{pmatrix}
        y_1\\
        \vdots\\
        y_K
    \end{pmatrix}.
\]

Since $\D$ is finite-dimensional, several technical issues become easier.
Completeness is automatic, all norms induce the same topology, and linear maps
between finite-dimensional spaces are continuous.

Nevertheless, the data space still has geometry. We must decide how to measure
data residuals. If the observational noise covariance is
\[
    C_D,
\]
then the statistically natural data misfit is
\[
    \|\obs-\G m\|_{C_D^{-1}}^2
    =
    (\obs-\G m)^T C_D^{-1}(\obs-\G m).
\]

This is not merely a numerical penalty. It is the quadratic term in the
negative log likelihood for Gaussian observational noise.

\subsection{Forward map}

The forward map sends models to predicted data:
\[
    \G : \M \to \D.
\]

In the toy problem, the components of $\G$ might be
\[
    [\G m]_i
    =
    \int_\Omega K_i(x)m(x)\,dx,
    \qquad
    i=1,\ldots,K.
\]

The map $\G$ must be meaningful on the chosen model space. For example, if
\[
    \M=L^2(\Omega)
\]
and
\[
    K_i\in L^2(\Omega),
\]
then
\[
    |[\G m]_i|
    \leq
    \|K_i\|_{L^2}\|m\|_{L^2},
\]
so each component of $\G$ is a continuous linear functional on $L^2(\Omega)$.

This compatibility matters. A forward map is not merely a matrix; it is an
operator between chosen spaces. The matrix appears only after discretization.

\subsection{Noise model}

The observational noise is Gaussian:
\[
    \noise \sim \mathcal{N}(0,C_D).
\]

Thus, conditional on the model $m$, the data are distributed as
\[
    \obs\mid m
    \sim
    \mathcal{N}(\G m,C_D).
\]

The likelihood is therefore
\[
    \pi(\obs\mid m)
    \propto
    \exp
    \left(
        -
        \frac{1}{2}
        \|\obs-\G m\|_{C_D^{-1}}^2
    \right).
\]

Equivalently, define the misfit potential
\[
    \Phi(m;\obs)
    =
    \frac{1}{2}
    \|\obs-\G m\|_{C_D^{-1}}^2.
\]
Then
\[
    \pi(\obs\mid m)
    \propto
    \exp(-\Phi(m;\obs)).
\]

\subsection{Prior measure}

The prior is a probability measure on the model space:
\[
    \prior = \mathcal{N}(m_0,C_0).
\]

Here
\[
    m_0\in\M
\]
is the prior mean, and
\[
    C_0:\M\to\M
\]
is a covariance operator.

For this to define a Gaussian measure on $\M$, the covariance operator must be
linear, symmetric, positive, and trace class:
\[
    \Tr(C_0)<\infty.
\]

This is the step that repairs the naive prior. We no longer assign an
independent variance $\sigma^2$ to every coefficient introduced by the mesh.
Instead, we choose a covariance operator that defines a genuine probability
measure on functions.

If
\[
    C_0 e_j=\lambda_j e_j
\]
for an orthonormal basis $\{e_j\}_{j=1}^{\infty}$ of $\M$, then the trace-class
condition is
\[
    \sum_{j=1}^{\infty}\lambda_j<\infty.
\]

The prior perturbation can then be represented formally as
\[
    m-m_0
    =
    \sum_{j=1}^{\infty}
    \sqrt{\lambda_j}\xi_j e_j,
    \qquad
    \xi_j\sim\mathcal{N}(0,1).
\]

The decay of the eigenvalues controls the regularity and total variance of
prior samples.

\begin{intuition}
The prior covariance is the instrument panel of the random function. It controls
amplitude, smoothness, correlation length, and boundary behaviour.
\end{intuition}

\section{The posterior}

Once the ingredients are chosen, the posterior is no longer a numerical guess.
It is a probability measure on the same model space as the prior.

Using the potential
\[
    \Phi(m;\obs)
    =
    \frac{1}{2}
    \|\obs-\G m\|_{C_D^{-1}}^2,
\]
Bayes' theorem in function space can be written as
\[
    \frac{d\post}{d\prior}(m)
    =
    \frac{1}{Z(\obs)}
    \exp(-\Phi(m;\obs)),
\]
where
\[
    Z(\obs)
    =
    \int_{\M}
    \exp(-\Phi(m;\obs))
    \,d\prior(m).
\]

This expression says that the posterior is obtained by reweighting the prior:
models that fit the data well receive more posterior mass, while models that
fit poorly receive less.

\subsection{Gaussian conditioning in function space}

In the linear Gaussian case, the posterior is again Gaussian:
\[
    \post
    =
    \mathcal{N}(m^y,C^y).
\]

This is the function-space analogue of the finite-dimensional Gaussian
conditioning formula. The difference is that all objects now have continuous
meaning. The covariance is an operator. The adjoint is the Hilbert adjoint. The
prior is a measure on $\M$.

The data covariance induced by the prior and noise is
\[
    \G C_0 \Gadj + C_D.
\]
Since
\[
    \D=\R^K,
\]
this is a finite-dimensional covariance operator on the data space. In
coordinates, it is a $K\times K$ matrix. The observational noise covariance
$C_D$ is positive definite, so this sum is invertible under the usual Gaussian
noise assumptions.

\subsection{Posterior covariance}

The posterior covariance is
\[
    C^y
    =
    C_0
    -
    C_0 \Gadj
    \left(
        \G C_0 \Gadj + C_D
    \right)^{-1}
    \G C_0.
\]

This formula has the same shape as the finite-dimensional formula, but its
meaning is different. Each symbol now refers to an operator between chosen
spaces:
\[
    C_0:\M\to\M,
    \qquad
    \G:\M\to\D,
    \qquad
    \Gadj:\D\to\M,
    \qquad
    C_D:\D\to\D.
\]

The term
\[
    C_0 \Gadj
    \left(
        \G C_0 \Gadj + C_D
    \right)^{-1}
    \G C_0
\]
is the amount of prior covariance removed by the data. Since the data space is
finite-dimensional, the update can only affect finitely many data-informed
directions. The remaining uncertainty is controlled by the prior covariance
operator.

Because the prior covariance is trace class, the posterior covariance remains
controlled. In particular, the posterior covariance is no larger than the prior
covariance in the positive-operator sense:
\[
    0 \leq C^y \leq C_0.
\]
Consequently,
\[
    \Tr(C^y)\leq \Tr(C_0)<\infty.
\]

This is the mathematical antidote to the divergence seen in the naive posterior.

\subsection{Posterior mean}

The posterior mean is
\[
    m^y
    =
    m_0
    +
    C_0 \Gadj
    \left(
        \G C_0 \Gadj + C_D
    \right)^{-1}
    \left(
        \obs - \G m_0
    \right).
\]

The residual
\[
    \obs-\G m_0
\]
lives in the data space. The inverse
\[
    \left(
        \G C_0 \Gadj + C_D
    \right)^{-1}
\]
weights it by the combined uncertainty from the prior and the noise. The
adjoint
\[
    \Gadj
\]
then pulls this data-space information back into the model space. Finally,
$C_0$ shapes the update according to the prior covariance.

Thus, the posterior mean is not merely a least-squares solution. It is a
geometry-aware, covariance-aware update of the prior mean.

\begin{lesson}
The posterior mean is the prior mean plus a data-driven correction. The adjoint
decides how data residuals enter model space, and the covariance decides which
model directions are allowed to respond.
\end{lesson}

\section{Discretization after formulation}

We now discretize, but the order has changed. We are no longer using
discretization to define the inverse problem. We are using discretization to
approximate an inverse problem that already exists.

\subsection{Choosing finite-dimensional subspaces}

Choose finite-dimensional approximation spaces
\[
    \M_N \subset \M.
\]

For example, we may take
\[
    \M_N
    =
    \operatorname{span}\{\phi_1,\ldots,\phi_N\},
\]
where the $\phi_j$ are hat functions on a mesh.

A discrete model is then
\[
    m_N(x)
    =
    \sum_{j=1}^N u_j\phi_j(x).
\]

The key difference from the naive approach is conceptual:
\[
    \M_N
    \text{ approximates }
    \M.
\]
It does not replace $\M$ as the starting point of the problem.

\subsection{Projecting the continuous problem}

Let
\[
    P_N:\M\to\M_N
\]
be a projection or approximation operator. The goal is to build finite
dimensional objects that approximate the continuous ones:
\[
    m_0^N \approx m_0,
    \qquad
    C_0^N \approx C_0,
    \qquad
    \G_N \approx \G.
\]

For instance, the discrete forward matrix may be assembled by applying the
continuous forward map to the basis functions:
\[
    G_{ij}
    =
    [\G\phi_j]_i.
\]

The prior covariance should be obtained from the continuous covariance operator,
not invented independently at each resolution. Depending on the discretization,
one may define its matrix representation through the action of $C_0$ on basis
or test functions.

The guiding principle is:
\[
    \text{discrete object}
    =
    \text{representation of continuous object}.
\]

\subsection{Respecting inner products}

If the basis functions are not orthonormal in $\M$, the coefficient-space
Euclidean dot product does not represent the model-space inner product.

The model-space Gram matrix is
\[
    M_{ij}
    =
    \langle \phi_i,\phi_j\rangle_{\M}.
\]

For coefficient vectors $\bm{u}$ and $\bm{v}$ representing
\[
    m_N=\sum_{j=1}^N u_j\phi_j,
    \qquad
    n_N=\sum_{j=1}^N v_j\phi_j,
\]
we have
\[
    \langle m_N,n_N\rangle_{\M}
    =
    \bm{u}^T M\bm{v}.
\]

Thus, the mass or Gram matrix is not optional garnish. It is the coordinate
expression of the model-space geometry.

Similarly, if the data-space inner product is represented by
\[
    W_D,
\]
then
\[
    \langle a,b\rangle_{\D}
    =
    a^T W_D b.
\]

The discrete adjoint must respect both geometries:
\[
    G^\ast_{\mathrm{disc}}
    =
    M^{-1}G^T W_D.
\]

\begin{warning}
Replacing $G^\ast_{\mathrm{disc}}$ by $G^T$ is valid only in the special case
where the chosen coordinates are orthonormal for the chosen inner products.
\end{warning}

\subsection{Respecting covariance operators}

The discrete covariance must approximate the continuous covariance operator
$C_0$.

Suppose $C_0$ has spectral representation
\[
    C_0 e_j=\lambda_j e_j.
\]
Then a discretization should reproduce, as accurately as appropriate, the
variance assigned to the corresponding continuous modes.

If $C_0$ is defined through a differential operator, such as
\[
    C_0
    =
    \alpha(\kappa^2 I-\Delta)^{-s},
\]
then the discrete covariance should come from a discretization of that operator
and its inverse power.

For example, if $A=\kappa^2 I-\Delta$ is discretized as $A_N$, then a natural
discrete covariance is based on
\[
    C_{0,N}
    \approx
    \alpha A_N^{-s},
\]
with the geometry of the model space included through the appropriate mass or
Gram matrices.

The important point is that $C_{0,N}$ is not chosen merely because it is a
positive matrix. It is chosen because it approximates $C_0$.

\subsection{Respecting boundary conditions}

If the prior covariance is built from a differential operator, then boundary
conditions are part of the prior.

For example,
\[
    C_0
    =
    \alpha(\kappa^2 I-\Delta)^{-s}
\]
is not fully defined until the domain of $\Delta$ is specified. Dirichlet,
Neumann, periodic, and mixed boundary conditions lead to different covariance
operators.

Therefore, the discretization must respect the same boundary conditions.

If the continuous prior imposes Dirichlet behaviour, the discrete basis or
operator should reflect that. If the continuous prior uses Neumann boundary
conditions, the discrete operator should encode those instead.

\begin{lesson}
Boundary conditions do not enter at the edge of the code. They enter at the
definition of the covariance operator.
\end{lesson}

\section{Why the initial problems melt away}

We can now revisit the pathologies of the naive posterior and see why they
disappear when the formulation is done in the right order.

\subsection{The prior has a continuous meaning}

In the naive construction, we chose
\[
    \bm{u}_N\sim\mathcal{N}(0,\sigma^2 I_N)
\]
for every $N$. Each prior was valid in its own finite-dimensional coordinate
space, but the sequence did not correspond to a meaningful Gaussian prior on
functions.

In the function-space construction, we choose
\[
    m\sim\mathcal{N}(m_0,C_0)
\]
first.

The discrete priors are then approximations of this same prior. They are not
independent inventions at each resolution.

This changes the meaning of refinement. Increasing $N$ no longer introduces a
new Bayesian model with more independent variance. Instead, it improves the
approximation of the same continuous prior.

\subsection{The covariance has bounded trace}

The naive identity covariance failed because
\[
    \Tr(\sigma^2 I_N)
    =
    \sigma^2 N
    \to\infty.
\]

The continuous version of this failure is
\[
    \Tr(\sigma^2 I)
    =
    \infty.
\]

In the corrected formulation, we require
\[
    \Tr(C_0)<\infty.
\]
Therefore,
\[
    \mathbb{E}\|m-m_0\|_{\M}^2
    =
    \Tr(C_0)
    <
    \infty.
\]

The posterior covariance satisfies
\[
    0\leq C^y\leq C_0,
\]
and therefore
\[
    \Tr(C^y)\leq \Tr(C_0)<\infty.
\]

Thus, the posterior uncertainty is controlled by construction.

\begin{proposition}
Let $\M$ be a Hilbert space, let $C_0$ be a positive symmetric trace-class
covariance operator on $\M$, and let $\G:\M\to\D$ be a bounded linear map into
the finite-dimensional data space $\D=\R^K$. Suppose $C_D$ is symmetric positive
definite. Then the posterior covariance
\[
    C^y
    =
    C_0
    -
    C_0\Gadj
    \left(
        \G C_0\Gadj + C_D
    \right)^{-1}
    \G C_0
\]
is positive and trace class, with
\[
    \Tr(C^y)\leq\Tr(C_0).
\]
\end{proposition}

\begin{proof}
The posterior covariance is obtained by conditioning a joint Gaussian measure.
Conditioning reduces uncertainty in the positive-operator sense, so
\[
    0\leq C^y\leq C_0.
\]
Since $C_0$ is trace class and $C^y$ is positive and dominated by $C_0$, it is
also trace class and satisfies
\[
    \Tr(C^y)\leq \Tr(C_0).
\]
\end{proof}

\subsection{The adjoint is geometry-aware}

In the naive discretization, the posterior formula used $G^T$ without asking
whether this was the correct adjoint.

In the function-space formulation, the adjoint is defined by
\[
    \langle \G m,d\rangle_{\D}
    =
    \langle m,\Gadj d\rangle_{\M}.
\]

After discretization, this becomes
\[
    G^\ast_{\mathrm{disc}}
    =
    M^{-1}G^T W_D,
\]
where
\[
    M_{ij}
    =
    \langle \phi_i,\phi_j\rangle_{\M}
\]
and $W_D$ represents the data-space inner product.

Thus, the discrete posterior formulas are built from the adjoint determined by
the spaces, not from a raw transpose.

The model-space geometry now travels with the computation.

\subsection{Samples remain controlled under refinement}

A sample from the continuous prior has the form
\[
    m
    =
    m_0
    +
    \sum_{j=1}^{\infty}
    \sqrt{\lambda_j}\xi_j e_j,
    \qquad
    \sum_{j=1}^{\infty}\lambda_j<\infty.
\]

A sample from a discretized approximation uses only finitely many degrees of
freedom, or a finite-dimensional projection of the continuous covariance. As
$N$ increases, the discretized samples add finer details in a way controlled by
the eigenvalue decay of $C_0$.

Unlike the naive prior, refinement does not add independent variance of fixed
size to every new direction. New directions receive the amount of variance
prescribed by the continuous covariance operator. For a trace-class covariance,
the total remaining variance in high modes becomes small.

If the posterior covariance is discretized consistently, then posterior samples
also remain controlled:
\[
    \Tr(C_N^y)
    \lesssim
    \Tr(C_0)
\]
in the appropriate discretized geometry.

\begin{intuition}
Refinement no longer opens new trapdoors of variance. It reveals finer pieces
of the same probability measure.
\end{intuition}

\subsection{The posterior converges}

The goal of discretization is now clear. The continuous posterior
\[
    \post=\mathcal{N}(m^y,C^y)
\]
is the target. The discrete posteriors
\[
    \post_N
\]
should approximate it.

In a well-designed discretization, one expects convergence in an appropriate
sense:
\[
    \post_N \to \post
    \qquad
    \text{as}
    \qquad
    N\to\infty.
\]

The exact meaning of convergence depends on the setting. It might be convergence
of posterior means,
\[
    m_N^y\to m^y,
\]
convergence of covariance operators,
\[
    C_N^y\to C^y,
\]
or convergence of probability measures in a metric such as the Hellinger
distance or another suitable topology.

The critical point is that there is now a target posterior to converge to.

This was missing in the naive workflow. There, the sequence
\[
    \post_1,\post_2,\post_3,\ldots
\]
was not clearly approximating anything. In the corrected workflow, the
continuous posterior is defined first, and the discrete posteriors are judged by
how well they approximate it.

\begin{lesson}
The cure is not a cleverer matrix.
The cure is choosing the mathematical object first and discretizing it second.
\end{lesson}

\section{The corrected workflow}

It is useful to contrast the naive workflow with the corrected one.

The naive workflow was:
\[
    \text{mesh}
    \quad\Rightarrow\quad
    \text{basis}
    \quad\Rightarrow\quad
    \text{coefficient vector}
    \quad\Rightarrow\quad
    \text{matrix prior}
    \quad\Rightarrow\quad
    \text{posterior}.
\]

The corrected workflow is:
\[
    \text{model space}
    \quad\Rightarrow\quad
    \text{data space}
    \quad\Rightarrow\quad
    \text{forward map}
    \quad\Rightarrow\quad
    \text{noise model}
    \quad\Rightarrow\quad
    \text{prior measure}
    \quad\Rightarrow\quad
    \text{posterior}
    \quad\Rightarrow\quad
    \text{discretization}.
\]

More explicitly:
\begin{enumerate}[label=\arabic*.]
    \item Choose the model space $\M$.
    \item Choose the data space $\D$.
    \item Define the forward map $\G:\M\to\D$.
    \item Specify the noise model $\noise\sim\mathcal{N}(0,C_D)$.
    \item Choose the prior measure $\prior=\mathcal{N}(m_0,C_0)$ on $\M$.
    \item Define the posterior measure $\post$ on $\M$.
    \item Choose approximation spaces $\M_N\subset\M$.
    \item Discretize $\G$, $\Gadj$, $C_0$, and the posterior consistently.
\end{enumerate}

This order protects the problem from accidental modelling choices.

\section{The moral of Act VII}

The opening example looked harmless because every finite-dimensional step was
legal. The matrix $G$ was legal. The diagonal covariance matrix was legal. The
Gaussian conditioning formula was legal. The posterior mean was computable.

But legal finite-dimensional algebra does not guarantee a coherent
infinite-dimensional inverse problem.

The corrected formulation changes the foundation. We first choose:
\[
    \M,
    \qquad
    \D,
    \qquad
    \G,
    \qquad
    C_D,
    \qquad
    \prior.
\]
Then we define:
\[
    \post
    =
    \mathcal{N}(m^y,C^y).
\]
Only then do we discretize.

The posterior mean, covariance, samples, and refinement behaviour now belong to
one mathematical story. The plot no longer has to hide a monster under the
mean.

\begin{checkpoint}
The final return to the toy problem is now straightforward. We will solve the
same opening example again, but this time the prior will be a genuine
function-space Gaussian, the adjoint will respect the chosen geometry, and mesh
refinement will approximate the posterior rather than rewriting the problem.
\end{checkpoint}
% ============================================================
% ACT VIII
% ============================================================

\part{Act VIII: Return to the Toy Problem}

\chapter{Re-solving the Opening Example}

\section{The original naive workflow}

We now return to the toy inverse problem from the beginning. The data equation
was
\[
    \obs = \G m + \noise,
\]
where the unknown model was a function on
\[
    \Omega=[0,L],
\]
the data lived in
\[
    \D=\R^K,
\]
and the observational noise was Gaussian:
\[
    \noise\sim\mathcal{N}(0,C_D).
\]

In Act I, we discretized immediately. The result seemed reasonable at first:
the posterior mean looked acceptable. But Act II revealed the trouble. When we
sampled the posterior, the uncertainty behaved badly under refinement. The
posterior was not converging toward a sensible function-space object.

The purpose of this final act is to solve the same problem again, but with the
order corrected.

\subsection{Basis first}

The naive workflow began by choosing a mesh
\[
    0=x_1<x_2<\cdots<x_N=L
\]
and a corresponding set of hat functions
\[
    \{\phi_1,\ldots,\phi_N\}.
\]
The model was then approximated by
\[
    m_N(x)
    =
    \sum_{j=1}^N u_j\phi_j(x).
\]

This step is not wrong by itself. Hat functions are perfectly legitimate
approximation tools. The problem was that they were asked to do too much. They
did not merely approximate a pre-existing function-space problem. They became
the starting point of the formulation.

The unknown was no longer treated as a function
\[
    m\in\M.
\]
It was treated as a coefficient vector
\[
    \bm{u}\in\R^N.
\]

This moved the problem from function space into coefficient space before the
model space, prior measure, and geometry had been chosen.

\begin{warning}
A basis should represent the unknown after the model space is chosen. It should
not choose the model space by stealth.
\end{warning}

\subsection{Diagonal covariance}

After introducing coefficients, the naive workflow placed the prior
\[
    \bm{u}\sim\mathcal{N}(0,\sigma^2 I_N)
\]
directly on the coefficient vector.

For fixed $N$, this is a valid finite-dimensional Gaussian distribution. But
as $N$ changes, the prior changes too:
\[
    \mathcal{N}(0,\sigma^2 I_1),
    \quad
    \mathcal{N}(0,\sigma^2 I_2),
    \quad
    \mathcal{N}(0,\sigma^2 I_3),
    \quad
    \ldots
\]

These are not automatically approximations of a single probability measure on
functions.

The trace of the prior covariance is
\[
    \Tr(\sigma^2 I_N)
    =
    \sigma^2 N.
\]
Therefore,
\[
    \Tr(\sigma^2 I_N)\to\infty
    \qquad
    \text{as}
    \qquad
    N\to\infty.
\]

This was the hidden reservoir of instability. Every mesh refinement introduced
new independent variance.

\subsection{Euclidean transpose}

The naive workflow also used the transpose
\[
    G^T
\]
as though it were automatically the adjoint of the forward operator.

But the adjoint is defined by
\[
    \langle \G m,d\rangle_{\D}
    =
    \langle m,\Gadj d\rangle_{\M}.
\]
It depends on the inner products of both the model space and the data space.

If the model is represented using non-orthonormal hat functions, then the
model-space inner product is represented by a Gram matrix
\[
    M_{ij}
    =
    \langle \phi_i,\phi_j\rangle_{\M}.
\]
If the data-space inner product is represented by
\[
    W_D,
\]
then the discrete adjoint is
\[
    G^\ast_{\mathrm{disc}}
    =
    M^{-1}G^T W_D.
\]

Thus, using $G^T$ alone corresponds to a specific Euclidean coordinate geometry.
If that geometry is not the one intended by the function-space problem, then
the posterior formula is solving a different problem.

\subsection{Unstable refinement}

The finite-dimensional posterior covariance in the naive workflow was
\[
    C_N^y
    =
    C_N
    -
    C_NG^T
    \left(
        GC_NG^T+C_D
    \right)^{-1}
    GC_N.
\]

For
\[
    C_N=\sigma^2 I_N,
\]
the posterior covariance satisfies
\[
    \Tr(C_N^y)
    \geq
    \sigma^2(N-K),
\]
because the data can constrain at most $K$ independent directions. As
\[
    N\to\infty
\]
with $K$ fixed,
\[
    \Tr(C_N^y)\to\infty.
\]

This explains the unstable samples. The data informed some directions, but the
number of uninformed directions grew with the mesh. Each new direction carried
variance of size $\sigma^2$.

The posterior mean looked fine because means can hide variance. The samples
showed the truth.

\begin{lesson}
The naive workflow did not approximate one continuous inverse problem. It
created a new finite-dimensional Bayesian problem at every resolution.
\end{lesson}

\section{The corrected workflow}

We now rebuild the same example in the proper order.

\subsection{Choose the model space}

First choose the model space. For example, let
\[
    \M=L^2(\Omega)
\]
if we want square-integrable models, or let
\[
    \M=H^s(\Omega)
\]
if we want to encode a particular degree of Sobolev regularity.

The choice of $\M$ should be justified by the physical meaning of the model and
by the measurements we wish to apply. If the forward map consists of integrals
against kernels,
\[
    [\G m]_i
    =
    \int_\Omega K_i(x)m(x)\,dx,
\]
then $L^2(\Omega)$ is natural when
\[
    K_i\in L^2(\Omega).
\]
Indeed,
\[
    |[\G m]_i|
    \leq
    \|K_i\|_{L^2}\|m\|_{L^2}.
\]

If we want additional regularity, or if the prior is naturally expressed using
a differential operator, then a Sobolev space may be more appropriate.

The important point is that the unknown is now a function-space object:
\[
    m\in\M.
\]

\subsection{Choose the data space}

The data space is
\[
    \D=\R^K.
\]

The data vector is
\[
    \obs =
    \begin{pmatrix}
        y_1\\
        \vdots\\
        y_K
    \end{pmatrix}.
\]

We choose an inner product on $\D$. In the simplest case,
\[
    \langle a,b\rangle_{\D}=a^Tb.
\]
More generally,
\[
    \langle a,b\rangle_{\D}=a^TW_Db,
\]
where $W_D$ is symmetric positive definite.

The noise covariance is
\[
    C_D.
\]
The natural data misfit is
\[
    \|\obs-\G m\|_{C_D^{-1}}^2
    =
    (\obs-\G m)^T C_D^{-1}(\obs-\G m)
\]
in coordinates.

\subsection{Choose the forward operator}

Next define the forward map:
\[
    \G:\M\to\D.
\]

For the toy problem,
\[
    [\G m]_i
    =
    \int_\Omega K_i(x)m(x)\,dx,
    \qquad
    i=1,\ldots,K.
\]

This map is chosen before discretization. The discrete matrix will later be a
representation of $\G$, not a replacement for it.

If the basis functions are $\phi_j$, then the matrix entries are
\[
    G_{ij}
    =
    [\G\phi_j]_i
    =
    \int_\Omega K_i(x)\phi_j(x)\,dx.
\]

The adjoint is also defined at the function-space level:
\[
    \langle \G m,d\rangle_{\D}
    =
    \langle m,\Gadj d\rangle_{\M}.
\]

Only after bases are chosen do we obtain the discrete representation
\[
    G^\ast_{\mathrm{disc}}
    =
    M^{-1}G^T W_D.
\]

\subsection{Choose the prior measure}

Now choose the prior as a probability measure on $\M$:
\[
    m\sim\prior.
\]

For a Gaussian prior,
\[
    \prior=\mathcal{N}(m_0,C_0),
\]
where
\[
    C_0:\M\to\M
\]
is a covariance operator satisfying
\[
    \Tr(C_0)<\infty.
\]

A useful choice is a covariance built from a differential operator:
\[
    C_0
    =
    \alpha(\kappa^2 I-\Delta)^{-s}.
\]
Here:
\begin{enumerate}[label=\arabic*.]
    \item $\alpha>0$ controls the overall variance scale;
    \item $\kappa>0$ controls the correlation length scale;
    \item $s>0$ controls smoothness;
    \item the boundary conditions for $\Delta$ encode boundary behaviour.
\end{enumerate}

If
\[
    -\Delta e_j=\rho_j e_j,
\]
then
\[
    C_0 e_j
    =
    \alpha(\kappa^2+\rho_j)^{-s}e_j.
\]
The prior eigenvalues are therefore
\[
    \lambda_j
    =
    \alpha(\kappa^2+\rho_j)^{-s}.
\]

The trace condition is
\[
    \sum_{j=1}^{\infty}\lambda_j<\infty.
\]

This ensures that the prior has finite expected energy:
\[
    \mathbb{E}\|m-m_0\|_{\M}^2
    =
    \Tr(C_0)
    <
    \infty.
\]

\subsection{Only then discretize}

Only now do we choose a mesh and basis:
\[
    \M_N
    =
    \operatorname{span}\{\phi_1,\ldots,\phi_N\}
    \subset \M.
\]

The discrete model is
\[
    m_N(x)
    =
    \sum_{j=1}^N u_j\phi_j(x).
\]

But now the coefficients $\bm{u}$ are coordinates of a function in $\M_N$.
They are not the primary unknown. The primary unknown remains
\[
    m\in\M.
\]

We discretize each continuous object:
\[
    \G
    \leadsto
    G,
    \qquad
    \Gadj
    \leadsto
    G^\ast_{\mathrm{disc}},
    \qquad
    C_0
    \leadsto
    C_{0,N}.
\]

The model-space Gram matrix is
\[
    M_{ij}
    =
    \langle \phi_i,\phi_j\rangle_{\M}.
\]

The discrete adjoint is
\[
    G^\ast_{\mathrm{disc}}
    =
    M^{-1}G^T W_D.
\]

The discrete prior covariance is obtained by discretizing $C_0$, not by
choosing an unrelated diagonal matrix. The discrete posterior then approximates
the continuous posterior
\[
    \post=\mathcal{N}(m^y,C^y).
\]

\begin{lesson}
In the corrected workflow, the mesh is a microscope. In the naive workflow, the
mesh was the lawgiver.
\end{lesson}

\section{Comparison}

We now compare the two workflows.

\subsection{Posterior means}

In both workflows, we can compute a posterior mean.

The naive posterior mean is obtained from
\[
    \bar{\bm{u}}_N^y
    =
    \bm{u}_0
    +
    C_NG^T
    \left(
        GC_NG^T+C_D
    \right)^{-1}
    (\bm{y}-G\bm{u}_0).
\]

The corrected posterior mean is the discretization of the function-space
formula
\[
    m^y
    =
    m_0
    +
    C_0\Gadj
    \left(
        \G C_0\Gadj+C_D
    \right)^{-1}
    (\obs-\G m_0).
\]

The difference is not merely cosmetic. In the corrected formula:
\begin{enumerate}[label=\arabic*.]
    \item $C_0$ is a covariance operator on the model space;
    \item $\Gadj$ is determined by the inner products;
    \item the discretized covariance approximates the continuous covariance;
    \item the posterior mean approximates a function-space posterior mean.
\end{enumerate}

The two posterior means may look superficially similar in plots, especially at
coarse resolution. But only the corrected mean belongs to a posterior measure
with a coherent continuous limit.

\subsection{Posterior samples}

The contrast is sharper for samples.

In the naive workflow, samples are drawn from
\[
    \mathcal{N}(\bar{\bm{u}}_N^y,C_N^y)
\]
and reconstructed as
\[
    m_N^{(s)}(x)
    =
    \sum_{j=1}^N u_j^{(s)}\phi_j(x).
\]

As $N$ increases, new coefficient directions receive independent prior
variance. Many of these directions are weakly constrained or invisible to the
data. Consequently, samples can become increasingly oscillatory or unstable.

In the corrected workflow, samples are drawn from a discretization of
\[
    \post=\mathcal{N}(m^y,C^y).
\]
The continuous covariance controls how much variance is available in each
direction. If
\[
    C_0 e_j=\lambda_j e_j
\]
with
\[
    \sum_{j=1}^{\infty}\lambda_j<\infty,
\]
then high-index directions receive decaying variance.

Thus, refinement adds controlled detail rather than uncontrolled noise.

\subsection{Posterior variance}

The naive prior covariance has trace
\[
    \Tr(\sigma^2 I_N)=\sigma^2 N.
\]
The naive posterior covariance satisfies
\[
    \Tr(C_N^y)
    \geq
    \sigma^2(N-K).
\]
Hence, as
\[
    N\to\infty,
\]
the total posterior variance diverges.

In the corrected workflow, the prior covariance satisfies
\[
    \Tr(C_0)<\infty.
\]
The posterior covariance satisfies
\[
    0\leq C^y\leq C_0,
\]
and therefore
\[
    \Tr(C^y)\leq\Tr(C_0)<\infty.
\]

After discretization, the discrete posterior variance approximates the
continuous posterior variance rather than growing without bound.

This is the heart of the cure:
\[
    \text{finite trace}
    \quad\Rightarrow\quad
    \text{controlled uncertainty}.
\]

\subsection{Mesh refinement}

In the naive workflow, mesh refinement changes the prior:
\[
    \mathcal{N}(0,\sigma^2 I_N)
    \quad
    \text{changes with } N.
\]
Each new basis function receives variance $\sigma^2$, so refinement injects
new uncertainty.

In the corrected workflow, mesh refinement improves the approximation of fixed
objects:
\[
    \M_N\to\M,
    \qquad
    G_N\to\G,
    \qquad
    C_{0,N}\to C_0,
    \qquad
    \post_N\to\post.
\]

The posterior at resolution $N$ is no longer a separate model. It is an
approximation to the same function-space posterior.

A useful way to phrase the distinction is:
\[
    \text{naive refinement}
    =
    \text{more coefficients, more prior variance},
\]
whereas
\[
    \text{function-space refinement}
    =
    \text{better approximation of the same probability measure}.
\]

\begin{intuition}
In the naive workflow, refinement keeps building new rooms in the house and
forgetting to install walls. In the corrected workflow, refinement sharpens the
blueprint that was already there.
\end{intuition}

\section{What the example teaches}

The toy example teaches a simple lesson, but not a small one.

A discretization can play two very different roles. It can approximate a
continuous inverse problem that has already been formulated, or it can quietly
define a finite-dimensional problem in place of the continuous one.

The same hat functions may appear in both workflows. The same forward matrix
entries
\[
    G_{ij}
    =
    \int_\Omega K_i(x)\phi_j(x)\,dx
\]
may be computed. The same data may be used. The same posterior mean formula may
seem to appear.

But the mathematical meaning is different.

In the naive workflow:
\[
    \text{basis}
    \Rightarrow
    \text{coefficients}
    \Rightarrow
    \text{matrix prior}
    \Rightarrow
    \text{posterior}.
\]

In the corrected workflow:
\[
    \text{model space}
    \Rightarrow
    \text{prior measure}
    \Rightarrow
    \text{operator posterior}
    \Rightarrow
    \text{discretization}.
\]

The first workflow gives a sequence of finite-dimensional Bayesian problems.
The second gives a sequence of approximations to one Bayesian inverse problem.

\begin{checkpoint}
The same numerical basis can either approximate a coherent function-space
problem or define a sequence of unrelated finite-dimensional problems.
The difference is the order in which the modelling choices are made.
\end{checkpoint}

\section{Final summary}

The opening mistake was not obvious because every finite-dimensional operation
looked valid. The diagonal prior was a valid covariance matrix. The posterior
mean was computable. The plot looked plausible.

But the posterior samples exposed the missing structure. The prior did not have
a continuous meaning. The covariance did not have bounded trace. The adjoint was
treated as a transpose. Refinement changed the problem rather than improving the
approximation.

The corrected formulation fixes the order:
\[
    \boxed{
    \text{Think first, discretize later.}
    }
\]

More explicitly:
\begin{enumerate}[label=\arabic*.]
    \item Choose the model space $\M$.
    \item Choose the data space $\D$.
    \item Define the forward map $\G:\M\to\D$.
    \item Specify the noise covariance $C_D$.
    \item Choose the prior measure $\prior=\mathcal{N}(m_0,C_0)$.
    \item Define the posterior measure $\post$.
    \item Discretize the already-defined objects.
\end{enumerate}

Only then do the posterior mean, covariance, samples, and refinement behaviour
belong to one coherent mathematical story.

\begin{lesson}
The final answer is not that matrices are bad.
The final answer is that matrices should be shadows cast by mathematical
objects, not the objects themselves.
\end{lesson}
% ============================================================
% Back matter
% ============================================================

\backmatter

\chapter{Summary of the Argument}

We now collect the whole story in one place.

The central claim of this document is simple:

\[
    \text{A discretization should approximate an inverse problem.}
\]
It should not define the inverse problem by accident.

The opening example showed how easily this accident can happen. We began with
an unknown function, introduced a basis, placed a Gaussian prior directly on the
coefficients, used an ordinary transpose as if it were automatically the
adjoint, and computed a finite-dimensional posterior. The posterior mean looked
reasonable. But posterior samples and mesh refinement revealed the problem:
the prior did not correspond to a coherent probability measure on a function
space.

The repair was not to find a cleverer diagonal matrix. The repair was to return
to the modelling choices that should have been made before discretization.

\section{The five modelling choices}

A function-space Bayesian inverse problem should begin with five explicit
choices.

\begin{enumerate}[label=\arabic*.]
    \item Choose the model space $\M$.
    \item Choose the data space $\D$.
    \item Choose the forward map $\G$.
    \item Choose the noise model.
    \item Choose the prior measure $\prior$.
\end{enumerate}

These choices define the mathematical problem.

The model space tells us what kind of object the unknown is:
\[
    m\in\M.
\]
The data space tells us what kind of object the observations are:
\[
    \obs\in\D.
\]
The forward map connects them:
\[
    \G:\M\to\D.
\]
The noise model describes observational uncertainty:
\[
    \noise\sim\mathcal{N}(0,C_D).
\]
The prior measure describes uncertainty before observing the data:
\[
    m\sim\prior.
\]

Once these ingredients are chosen, the posterior is defined by Bayes' theorem.
For Gaussian observational noise, with
\[
    \Phi(m;\obs)
    =
    \frac{1}{2}
    \|\obs-\G m\|_{C_D^{-1}}^2,
\]
the posterior satisfies
\[
    \frac{d\post}{d\prior}(m)
    =
    \frac{1}{Z(\obs)}
    \exp(-\Phi(m;\obs)).
\]

Only after this formula has meaning should we introduce a mesh, basis, matrix,
or quadrature rule.

\begin{lesson}
The discretization is the numerical approximation of the Bayesian inverse
problem. It is not the Bayesian inverse problem itself.
\end{lesson}

\section{The three common mistakes}

The same failure appeared in three disguises.

\begin{enumerate}[label=\arabic*.]
    \item Treating the prior as just a matrix.
    \item Treating the adjoint as just a transpose.
    \item Treating discretization as formulation.
\end{enumerate}

\subsection*{Mistake 1: Treating the prior as just a matrix}

A covariance matrix
\[
    C_N\in\R^{N\times N}
\]
may be perfectly valid for every fixed $N$, but this does not mean that the
sequence of priors defines a probability measure on functions.

The naive choice
\[
    C_N=\sigma^2 I_N
\]
has
\[
    \Tr(C_N)=\sigma^2 N.
\]
Thus,
\[
    \Tr(C_N)\to\infty
    \qquad
    \text{as}
    \qquad
    N\to\infty.
\]

In function space, the corresponding identity covariance would satisfy
\[
    \Tr(\sigma^2 I)=\infty.
\]
It is therefore not a valid covariance operator for a Gaussian measure living
on an infinite-dimensional Hilbert space.

A function-space Gaussian prior
\[
    \prior=\mathcal{N}(m_0,C_0)
\]
requires a covariance operator that is linear, symmetric, positive, and trace
class:
\[
    \Tr(C_0)<\infty.
\]

\subsection*{Mistake 2: Treating the adjoint as just a transpose}

The adjoint is defined by the inner products:
\[
    \langle \G m,d\rangle_{\D}
    =
    \langle m,\Gadj d\rangle_{\M}.
\]

If a finite-dimensional model space is represented by a non-orthonormal basis
\[
    \{\phi_1,\ldots,\phi_N\},
\]
then the model-space inner product is represented by the Gram matrix
\[
    M_{ij}
    =
    \langle \phi_i,\phi_j\rangle_{\M}.
\]

If the data-space inner product is represented by
\[
    W_D,
\]
then the discrete adjoint is
\[
    G^\ast_{\mathrm{disc}}
    =
    M^{-1}G^T W_D.
\]

Only in Euclidean coordinates with orthonormal bases does this reduce to
\[
    G^\ast_{\mathrm{disc}}=G^T.
\]

\subsection*{Mistake 3: Treating discretization as formulation}

Discretization is necessary for computation, but it should enter after the
continuous problem exists.

The dangerous order is
\[
    \text{basis}
    \quad\Rightarrow\quad
    \text{coefficients}
    \quad\Rightarrow\quad
    \text{matrix prior}
    \quad\Rightarrow\quad
    \text{posterior}.
\]

The coherent order is
\[
    \text{model space}
    \quad\Rightarrow\quad
    \text{prior measure}
    \quad\Rightarrow\quad
    \text{posterior measure}
    \quad\Rightarrow\quad
    \text{discretization}.
\]

\section{The final slogan}

\begin{lesson}
Think first.
Discretize later.
\end{lesson}

% ============================================================
% Appendices
% ============================================================

\appendix

\chapter{Appendix A: Finite-Dimensional Gaussian Conditioning}

This appendix records the finite-dimensional Gaussian conditioning formulas used
throughout the document. These formulas are familiar, but writing them out helps
clarify exactly where the posterior mean and covariance come from.

\section{Block Gaussian formulas}

Let
\[
    \begin{pmatrix}
        x\\
        y
    \end{pmatrix}
    \sim
    \mathcal{N}
    \left(
    \begin{pmatrix}
        \bar{x}\\
        \bar{y}
    \end{pmatrix},
    \begin{pmatrix}
        C_{xx} & C_{xy}\\
        C_{yx} & C_{yy}
    \end{pmatrix}
    \right),
\]
where
\[
    C_{yx}=C_{xy}^T
\]
and $C_{yy}$ is invertible.

Then the conditional distribution of $x$ given $y$ is Gaussian:
\[
    x\mid y
    \sim
    \mathcal{N}
    \left(
        \bar{x}
        +
        C_{xy}C_{yy}^{-1}(y-\bar{y}),
        \,
        C_{xx}
        -
        C_{xy}C_{yy}^{-1}C_{yx}
    \right).
\]

Thus,
\[
    \mathbb{E}[x\mid y]
    =
    \bar{x}
    +
    C_{xy}C_{yy}^{-1}(y-\bar{y}),
\]
and
\[
    \operatorname{Cov}(x\mid y)
    =
    C_{xx}
    -
    C_{xy}C_{yy}^{-1}C_{yx}.
\]

\begin{intuition}
Conditioning a joint Gaussian is a projection of uncertainty. The mean moves
toward values compatible with the observed component, and the covariance loses
the uncertainty explained by that observation.
\end{intuition}

\section{Derivation of the posterior mean}

Consider the finite-dimensional Bayesian linear model
\[
    \bm{u}\sim\mathcal{N}(\bm{u}_0,C_N),
\]
and
\[
    \bm{y}=G\bm{u}+\bm{\eta},
    \qquad
    \bm{\eta}\sim\mathcal{N}(0,C_D),
\]
with $\bm{u}$ and $\bm{\eta}$ independent.

The mean of $\bm{y}$ is
\[
    \mathbb{E}[\bm{y}]
    =
    G\bm{u}_0.
\]

The covariance of $\bm{y}$ is
\[
    \operatorname{Cov}(\bm{y})
    =
    \operatorname{Cov}(G\bm{u}+\bm{\eta})
    =
    GC_NG^T+C_D.
\]

The cross-covariance between $\bm{u}$ and $\bm{y}$ is
\[
    \operatorname{Cov}(\bm{u},\bm{y})
    =
    \operatorname{Cov}(\bm{u},G\bm{u}+\bm{\eta})
    =
    C_NG^T.
\]

Therefore,
\[
    \begin{pmatrix}
        \bm{u}\\
        \bm{y}
    \end{pmatrix}
    \sim
    \mathcal{N}
    \left(
    \begin{pmatrix}
        \bm{u}_0\\
        G\bm{u}_0
    \end{pmatrix},
    \begin{pmatrix}
        C_N & C_NG^T\\
        GC_N & GC_NG^T+C_D
    \end{pmatrix}
    \right).
\]

Using the block Gaussian conditioning formula with
\[
    x=\bm{u},
    \qquad
    y=\bm{y},
\]
we obtain
\[
    \bar{\bm{u}}_N^y
    =
    \bm{u}_0
    +
    C_NG^T
    \left(
        GC_NG^T+C_D
    \right)^{-1}
    (\bm{y}-G\bm{u}_0).
\]

\section{Derivation of the posterior covariance}

Using the same block Gaussian formula, the posterior covariance is
\[
    C_N^y
    =
    C_N
    -
    C_NG^T
    \left(
        GC_NG^T+C_D
    \right)^{-1}
    GC_N.
\]

This can be interpreted as
\[
    \text{posterior covariance}
    =
    \text{prior covariance}
    -
    \text{data-informed covariance reduction}.
\]

The matrix
\[
    GC_NG^T+C_D
\]
is the covariance of the predicted data, including both prior uncertainty
pushed through the forward model and observational noise.

The term
\[
    C_NG^T
    \left(
        GC_NG^T+C_D
    \right)^{-1}
    GC_N
\]
is positive semidefinite. Hence conditioning cannot increase covariance:
\[
    0\leq C_N^y\leq C_N.
\]

\begin{warning}
This finite-dimensional formula is correct as linear algebra. The danger begins
when $C_N$, $G^T$, and the coefficient space are treated as if they automatically
represent a coherent function-space problem.
\end{warning}

\chapter{Appendix B: Gaussian Measures on Hilbert Spaces}

This appendix collects a few facts about Gaussian measures on Hilbert spaces.
The goal is not to give a complete measure-theoretic treatment, but to record
the concepts needed in the main text.

\section{Karhunen--Loeve expansions}

Let $\M$ be a Hilbert space, and let
\[
    m\sim\mathcal{N}(m_0,C_0),
\]
where $C_0$ is a positive, symmetric, trace-class covariance operator.

Suppose $C_0$ has eigenpairs
\[
    C_0e_j=\lambda_j e_j,
\]
where
\[
    \{e_j\}_{j=1}^{\infty}
\]
is an orthonormal basis of $\M$, and
\[
    \lambda_j\geq 0.
\]

Then the Gaussian random model may be represented as
\[
    m
    =
    m_0
    +
    \sum_{j=1}^{\infty}
    \sqrt{\lambda_j}\xi_j e_j,
\]
where
\[
    \xi_j\sim\mathcal{N}(0,1)
\]
are independent standard normal random variables.

This is the Karhunen--Loeve expansion.

The expected squared norm of the perturbation is
\[
    \mathbb{E}\|m-m_0\|_{\M}^2
    =
    \sum_{j=1}^{\infty}\lambda_j
    =
    \Tr(C_0).
\]

Thus, the expansion defines a random element of $\M$ when
\[
    \sum_{j=1}^{\infty}\lambda_j<\infty.
\]

\begin{intuition}
The eigenfunctions are the directions of prior variation. The eigenvalues are
the variances assigned to those directions.
\end{intuition}

\section{Trace-class covariance operators}

A positive symmetric operator $C_0:\M\to\M$ is trace class if
\[
    \Tr(C_0)
    =
    \sum_{j=1}^{\infty}
    \langle C_0e_j,e_j\rangle_{\M}
    <
    \infty
\]
for one, and hence every, orthonormal basis $\{e_j\}$ of $\M$.

If $C_0e_j=\lambda_j e_j$, then
\[
    \Tr(C_0)
    =
    \sum_{j=1}^{\infty}\lambda_j.
\]

Trace class is the infinite-dimensional replacement for the finite-dimensional
fact that covariance matrices have finite trace.

The identity operator is not trace class on an infinite-dimensional Hilbert
space:
\[
    \Tr(I)
    =
    \sum_{j=1}^{\infty}1
    =
    \infty.
\]
Therefore,
\[
    \Tr(\sigma^2 I)
    =
    \infty.
\]

This is why a formal Gaussian prior with covariance $\sigma^2 I$ does not
define a Gaussian probability measure on the Hilbert space itself.

\begin{lesson}
Trace class means finite total prior variance. Without it, the Gaussian cloud
does not fit inside the chosen Hilbert space.
\end{lesson}

\section{Cameron--Martin space}

The Cameron--Martin space describes the directions in which a Gaussian measure
can be shifted in an especially regular way.

Let
\[
    \mu=\mathcal{N}(0,C_0)
\]
be a Gaussian measure on a Hilbert space $\M$. Suppose
\[
    C_0e_j=\lambda_j e_j,
\]
with $\lambda_j>0$ on the active directions of the prior.

The Cameron--Martin space is formally
\[
    \mathcal{H}_{\mu}
    =
    \operatorname{Range}(C_0^{1/2}),
\]
equipped with the norm
\[
    \|h\|_{\mathcal{H}_{\mu}}^2
    =
    \sum_{j=1}^{\infty}
    \frac{h_j^2}{\lambda_j},
\]
where
\[
    h=\sum_{j=1}^{\infty}h_j e_j.
\]

This space is often smoother or smaller than the space where typical prior
samples live. This can feel surprising at first.

For example, a Gaussian process may have rough typical samples, while its
Cameron--Martin space contains smoother functions. This distinction is central
in infinite-dimensional Bayesian analysis.

\begin{warning}
The Cameron--Martin space is not usually the same thing as the set of typical
samples. It describes admissible shifts of the Gaussian measure, not the rough
texture of an ordinary draw.
\end{warning}

\chapter{Appendix C: Discrete Bases and Gram Matrices}

This appendix records the coordinate bookkeeping needed when a function-space
problem is represented using a finite-dimensional basis.

\section{Coordinate vectors versus functions}

Let
\[
    \M_N
    =
    \operatorname{span}\{\phi_1,\ldots,\phi_N\}
    \subset \M.
\]
A coefficient vector
\[
    \bm{u}
    =
    \begin{pmatrix}
        u_1\\
        \vdots\\
        u_N
    \end{pmatrix}
    \in\R^N
\]
represents the function
\[
    m_N
    =
    \sum_{j=1}^N u_j\phi_j.
\]

The vector $\bm{u}$ is not the function itself. It is a coordinate
representation of the function in the chosen basis.

If the basis changes, the coordinate vector changes, even if the represented
function remains the same.

\begin{lesson}
Coordinates are names for objects. Changing the names should not change the
object.
\end{lesson}

\section{Mass matrices}

The model-space Gram matrix is
\[
    M_{ij}
    =
    \langle \phi_i,\phi_j\rangle_{\M}.
\]

If
\[
    m_N=\sum_{j=1}^N u_j\phi_j,
    \qquad
    n_N=\sum_{j=1}^N v_j\phi_j,
\]
then
\[
    \langle m_N,n_N\rangle_{\M}
    =
    \bm{u}^T M\bm{v}.
\]

For
\[
    \M=L^2(\Omega),
\]
the Gram matrix is the mass matrix
\[
    M_{ij}
    =
    \int_\Omega \phi_i(x)\phi_j(x)\,dx.
\]

For
\[
    \M=H^1(\Omega),
\]
with inner product
\[
    \langle f,g\rangle_{H^1}
    =
    \int_\Omega f(x)g(x)\,dx
    +
    \int_\Omega f'(x)g'(x)\,dx,
\]
the Gram matrix is
\[
    M_{ij}
    =
    \int_\Omega \phi_i(x)\phi_j(x)\,dx
    +
    \int_\Omega \phi_i'(x)\phi_j'(x)\,dx.
\]

Thus, the same basis has different Gram matrices in different model-space
geometries.

\section{Discrete adjoints}

Let
\[
    G\in\R^{K\times N}
\]
represent the forward map from coefficients to data:
\[
    \bm{d}=G\bm{u}.
\]

Let the model-space inner product be represented by
\[
    M
\]
and the data-space inner product by
\[
    W_D.
\]

The discrete adjoint is the matrix $G^\ast_{\mathrm{disc}}$ satisfying
\[
    \langle G\bm{u},\bm{d}\rangle_{\D}
    =
    \langle \bm{u},G^\ast_{\mathrm{disc}}\bm{d}\rangle_{\M}
\]
for all
\[
    \bm{u}\in\R^N,
    \qquad
    \bm{d}\in\R^K.
\]

Expanding both sides gives
\[
    (G\bm{u})^T W_D\bm{d}
    =
    \bm{u}^T M G^\ast_{\mathrm{disc}}\bm{d}.
\]
Hence
\[
    \bm{u}^T G^TW_D\bm{d}
    =
    \bm{u}^T M G^\ast_{\mathrm{disc}}\bm{d}.
\]
Since this must hold for all $\bm{u}$ and $\bm{d}$,
\[
    G^TW_D
    =
    M G^\ast_{\mathrm{disc}}.
\]
Therefore,
\[
    G^\ast_{\mathrm{disc}}
    =
    M^{-1}G^T W_D.
\]

Only if
\[
    M=I_N
    \qquad
    \text{and}
    \qquad
    W_D=I_K
\]
do we recover
\[
    G^\ast_{\mathrm{disc}}=G^T.
\]

\section{Projection and interpolation}

There are two common ways to associate a continuous model with a discrete
function.

The first is interpolation. Given nodal values
\[
    m(x_j),
\]
we define
\[
    I_Nm
    =
    \sum_{j=1}^N m(x_j)\phi_j.
\]

This is natural when pointwise values are meaningful.

The second is projection. The orthogonal projection $P_Nm\in\M_N$ is defined by
\[
    \langle m-P_Nm,v_N\rangle_{\M}=0
    \qquad
    \forall v_N\in\M_N.
\]

Writing
\[
    P_Nm=\sum_{j=1}^N u_j\phi_j,
\]
the coefficients satisfy
\[
    \sum_{j=1}^N
    M_{ij}u_j
    =
    \langle m,\phi_i\rangle_{\M},
    \qquad
    i=1,\ldots,N.
\]

In matrix form,
\[
    M\bm{u}=\bm{b},
\]
where
\[
    b_i=\langle m,\phi_i\rangle_{\M}.
\]

Projection depends on the model-space inner product. Interpolation depends on
pointwise evaluation. These are not the same operation.

\begin{warning}
Interpolation and projection can produce different discrete functions. Choosing
between them is a modelling and numerical decision, not harmless punctuation.
\end{warning}

\chapter{Appendix D: Suggested Exercises}

The exercises below are designed to reinforce the main argument. They are not
intended as a separate problem set in numerical analysis, but as guided checks
on the conceptual machinery.

\section{Exercise set for Act I}

\begin{enumerate}[label=\textbf{I.\arabic*.}]
    \item Let $\Omega=[0,1]$ and define
    \[
        [\G m]_i
        =
        \int_0^1 K_i(x)m(x)\,dx.
    \]
    Show that if $K_i\in L^2(0,1)$, then $\G_i$ is a bounded linear functional
    on $L^2(0,1)$.

    \item Let $\{\phi_j\}_{j=1}^N$ be hat functions on a uniform mesh. Write
    down the matrix entries
    \[
        G_{ij}=[\G\phi_j]_i.
    \]
    Explain why this matrix is a representation of $\G$ only after the basis
    has been chosen.

    \item Suppose
    \[
        \bm{u}\sim\mathcal{N}(0,\sigma^2 I_N).
    \]
    What assumptions about the coefficients are encoded by this prior?

    \item Explain why a plausible posterior mean does not guarantee a
    plausible posterior measure.

    \item In your own words, distinguish between the unknown function $m$ and
    its coefficient vector $\bm{u}$.
\end{enumerate}

\section{Exercise set for Act II}

\begin{enumerate}[label=\textbf{II.\arabic*.}]
    \item Let
    \[
        \bm{u}\sim\mathcal{N}(\bar{\bm{u}},C).
    \]
    If $C=LL^T$ and $\bm{\xi}\sim\mathcal{N}(0,I)$, show that
    \[
        \bar{\bm{u}}+L\bm{\xi}
    \]
    has mean $\bar{\bm{u}}$ and covariance $C$.

    \item Let $G\in\R^{K\times N}$ with $N>K$. Show that
    \[
        \dim\operatorname{null}(G)\geq N-K.
    \]

    \item For the naive prior $C_N=\sigma^2 I_N$, show that if
    \[
        \bm{v}\in\operatorname{null}(G),
    \]
    then
    \[
        C_N^y\bm{v}=\sigma^2\bm{v}.
    \]

    \item Prove that
    \[
        \Tr(C_N^y)\geq \sigma^2(N-K)
    \]
    under the assumptions of the previous exercise.

    \item Explain why increasing $N$ with fixed $K$ creates a problem for the
    naive posterior.
\end{enumerate}

\section{Exercise set for Act III}

\begin{enumerate}[label=\textbf{III.\arabic*.}]
    \item List the five modelling choices required before discretization.

    \item Give an example of a model space $\M$ suitable for functions on
    $\Omega=[0,1]$. Explain what kind of models it admits.

    \item Suppose the data are point values
    \[
        y_i=m(x_i)+\eta_i.
    \]
    Why might $L^2(\Omega)$ be an insufficient model space for this forward map?

    \item Define the likelihood potential
    \[
        \Phi(m;\obs)
        =
        \frac{1}{2}
        \|\obs-\G m\|_{C_D^{-1}}^2.
    \]
    Explain how it enters the posterior density relative to the prior.

    \item Explain the difference between deterministic uniqueness of a model
    and uniqueness of a posterior measure.
\end{enumerate}

\section{Exercise set for Acts IV--VI}

\begin{enumerate}[label=\textbf{IV--VI.\arabic*.}]
    \item Define a normed space, a Banach space, and a Hilbert space. What extra
    structure is added at each step?

    \item Let
    \[
        \langle f,g\rangle_{L^2}
        =
        \int_\Omega f(x)g(x)\,dx.
    \]
    Show that this defines an inner product on $L^2(\Omega)$, up to equality
    almost everywhere.

    \item Explain why pointwise values are not generally meaningful for
    elements of $L^2(\Omega)$.

    \item Let
    \[
        C_0e_j=\lambda_j e_j
    \]
    with $\{e_j\}$ orthonormal. Show that
    \[
        \mathbb{E}\|m-m_0\|_{\M}^2
        =
        \sum_{j=1}^{\infty}\lambda_j
    \]
    for the formal expansion
    \[
        m=m_0+\sum_{j=1}^{\infty}\sqrt{\lambda_j}\xi_j e_j.
    \]

    \item Explain why
    \[
        C_0=\sigma^2 I
    \]
    is not trace class on an infinite-dimensional Hilbert space.

    \item Let
    \[
        C_0=\alpha(\kappa^2I-\Delta)^{-s}.
    \]
    Explain the roles of $\alpha$, $\kappa$, $s$, and the boundary conditions.

    \item Give an example of a covariance kernel $k(x,x')$ and explain what
    qualitative prior information it encodes.
\end{enumerate}

\section{Exercise set for Act VII}

\begin{enumerate}[label=\textbf{VII.\arabic*.}]
    \item Starting from
    \[
        \obs=\G m+\noise,
    \]
    write the function-space posterior mean and covariance for a linear
    Gaussian inverse problem.

    \item Explain why
    \[
        \G C_0\Gadj+C_D
    \]
    is a data-space covariance.

    \item Suppose $C_0$ is trace class and
    \[
        0\leq C^y\leq C_0.
    \]
    Explain why $C^y$ is also trace class.

    \item Let $\M_N\subset\M$ be a finite-dimensional approximation space.
    Explain what it means for $\M_N$ to approximate $\M$ rather than replace it.

    \item Describe how the discrete adjoint changes when the basis functions
    are not orthonormal.

    \item Compare the naive workflow
    \[
        \text{basis}\Rightarrow\text{matrix prior}\Rightarrow\text{posterior}
    \]
    with the corrected workflow
    \[
        \text{model space}\Rightarrow\text{prior measure}
        \Rightarrow\text{posterior}\Rightarrow\text{discretization}.
    \]

    \item In one paragraph, explain the slogan:
    \[
        \text{Think first. Discretize later.}
    \]
\end{enumerate}

% ============================================================
% End of document
% ============================================================

\bibliographystyle{plain}
\bibliography{bibliography}

\end{document}

\bibliographystyle{plain}
\bibliography{bibliography}

\end{document}